
\documentclass{article}
\usepackage{colortbl}
\usepackage{makecell}
\usepackage{multirow}
\usepackage{supertabular}

\begin{document}

\newcounter{utterance}

\centering \large Interaction Transcript for game `clean\_up', experiment `0\_easy\_7obj\_tr', episode 1 with Apertus{-}v1.5{-}8B.
\vspace{24pt}

{ \footnotesize  \setcounter{utterance}{1}
\setlength{\tabcolsep}{0pt}
\begin{supertabular}{c@{$\;$}|p{.15\linewidth}@{}p{.15\linewidth}p{.15\linewidth}p{.15\linewidth}p{.15\linewidth}p{.15\linewidth}}
    \# & $\;$A & \multicolumn{4}{c}{Game Master} & $\;\:$B\\
    \hline

    \theutterance \stepcounter{utterance}  
    & & \multicolumn{4}{p{0.6\linewidth}}{
        \cellcolor[rgb]{0.9,0.9,0.9}{
            \makecell[{{p{\linewidth}}}]{
                \texttt{\tiny{[P1$\langle$GM]}}
                \texttt{Ben senin oyun yöneticinim ve aşağıdaki ızgarayı oyun tahtası olarak kullandığınız işbirlikçi bir oyun oynuyorsunuz.} \\
\texttt{```} \\
\texttt{ 1234567} \\
\texttt{╔═══════╗} \\
\texttt{║◌◌◌◌◌T◌║ 1} \\
\texttt{║P◌◌◌◌◌◌║ 2} \\
\texttt{║◌◌AD◌┌─╢ 3} \\
\texttt{║◌◌◌┌─┤S║ 4} \\
\texttt{║O◌◌│◌├─╢ 5} \\
\texttt{╟───┼─┘◌║ 6} \\
\texttt{║◌M◌│◌◌◌║ 7} \\
\texttt{╚═══╧═══╝} \\
\\ 
\texttt{```} \\
\\ 
\texttt{* Üst kenarda x koordinatları sağa doğru artar, sağ kenarda ise y koordinatları aşağı doğru artar.} \\
\texttt{* Aşağıdaki nesneler oyun tahtasına rastgele yerleştirilmiştir: 'P', 'O', 'T', 'S', 'D', 'A', 'M'.} \\
\\ 
\\ 
\texttt{Diğer oyuncu, nesnelerin farklı rastgele konumlara yerleştirildiği oyun tahtasının bir varyasyonunu görür. Sen diğer oyuncunun tahtasını göremezsin ve onlar da seninkini göremez.} \\
\texttt{**Amaç:**} \\
\\ 
\texttt{Her iki oyuncunun da, kendi oyun tahtalarındaki nesneleri hareket ettirerek aynı türden nesnelerin aynı koordinatlara gelmesini sağlaması gerekir. Ortak bir hedef düzen üzerinde anlaşmak için diğer oyuncuyla iletişim kurmalısınız.} \\
\\ 
\\ 
\texttt{**Kurallar:**} \\
\\ 
\texttt{* Her turda, aşağıdaki iki komuttan yalnızca birini gönderebilirsin:} \\
\texttt{ 1. `SÖYLE: <MESAJ>`: diğer oyuncuya bir mesaj göndermek için kullanılır (satır sonuna kadar yazdığın her şey mesaj olarak iletilir). Ben bu mesajı partnerine ileteceğim.} \\
\texttt{2. `HAMLE: <NESNE>, (<X>, <Y>)`: bir nesneyi yeni bir konuma taşımak için kullanılır. Burada <X> sütunu, <Y> ise satırı ifade eder. Yaptığın hamlenin başarılı olup olmadığını sana söyleyeceğim.} \\
\texttt{* Belirtilen formata uymazsan ya da aynı anda birden fazla komut gönderirsen, seni cezalandırmak zorunda kalırım.} \\
\texttt{* Her iki oyuncu da 16 sayısından fazla ceza alırsa, oyunu birlikte kaybedersiniz.} \\
\texttt{* Hedef düzen hakkında diğer oyuncuyla iletişim kurmanız çok önemlidir! Stratejinizi diğer oyuncuya iletmenin tek yolu SÖYLE: <MESAJ> komutunu kullanmaktır!} \\
\\ 
\\ 
\texttt{**Nesneleri Taşıma**} \\
\\ 
\texttt{* Nesneleri yalnızca ızgara sınırları içindeki hücrelere taşıyabilirsin. Hedef hücre boş olmalıdır; yani yalnızca ‘◌’ sembolünü içermelidir.} \\
\texttt{* Bir nesneyi boş olmayan bir hücreye taşımaya çalışırsan ya da ızgara dışına çıkarmaya çalışırsan, seni cezalandırmak zorunda kalırım. Ardından tekrar deneme hakkın olur.} \\
\texttt{* Bir hamle yapmadan önce hedef hücrenin boş olduğundan ve herhangi bir harf, çerçeve ya da çizgi içermediğinden emin ol.} \\
\\ 
\\ 
\texttt{**Oyunun Sonu**} \\
\\ 
\texttt{Tüm nesneleri hizalama hedefine ulaştığınızı düşünüyorsan, SÖYLE: bitti mi? göndererek diğer oyuncudan oyunu bitirmesini isteyebilirsin. Eğer diğer oyuncu senden oyunu bitirmeni isterse ve sen SÖYLE: bitti! ile yanıt verirsen, oyun sona erer.} \\
\\ 
\texttt{Oyun, 28 tur içinde sona ererse her iki oyuncu da kazanır. Bir tur, iki oyuncunun da geçerli bir komut göndermesi olarak tanımlanır.} \\
\\ 
\texttt{**Puanlama:**} \\
\\ 
\texttt{Aynı nesneler her iki oyun tahtasında ne kadar yakın konumlanırsa, o kadar fazla puan kazanırsınız. Cezalar puanınızı düşürür. Rekoru kırabilir misiniz?} \\
\\ 
\\ 
\texttt{Oyuna başlamak için diğer oyuncuya bir mesaj gönder!} \\
            }
        }
    }
    & \\ \\

    \theutterance \stepcounter{utterance}  
    & \multicolumn{4}{p{0.6\linewidth}}{
        \cellcolor[rgb]{0.9,0.9,0.9}{
            \makecell[{{p{\linewidth}}}]{
                \texttt{\tiny{[P1$\rangle$GM]}}
                \texttt{SÖYLE: Merhaba! Hedefimiz, her türden nesnenin aynı koordinatlarda toplanması. Öncelikle stratejimizi belirleyelim.} \\
            }
        }
    }
    & & \\ \\

    \theutterance \stepcounter{utterance}  
    & & & \multicolumn{2}{p{0.3\linewidth}}{
        \cellcolor[rgb]{0.9,0.9,0.9}{
            \makecell[{{p{\linewidth}}}]{
                \texttt{\tiny{[GM$|$GM]}}
                \texttt{SÖYLE: Merhaba! Hedefimiz, her türden nesnenin aynı koordinatlarda toplanması. Öncelikle stratejimizi belirleyelim.} \\
            }
        }
    }
    & & \\ \\

    \theutterance \stepcounter{utterance}  
    & & \multicolumn{4}{p{0.6\linewidth}}{
        \cellcolor[rgb]{0.9,0.9,0.9}{
            \makecell[{{p{\linewidth}}}]{
                \texttt{\tiny{[GM$\rangle$P2]}}
                \texttt{Ben senin oyun yöneticinim ve aşağıdaki ızgarayı oyun tahtası olarak kullandığınız işbirlikçi bir oyun oynuyorsunuz.} \\
\texttt{```} \\
\texttt{ 1234567} \\
\texttt{╔═══════╗} \\
\texttt{║◌◌◌◌◌S◌║ 1} \\
\texttt{║◌◌◌◌O◌◌║ 2} \\
\texttt{║◌◌◌M◌┌─╢ 3} \\
\texttt{║◌◌A┌─┤◌║ 4} \\
\texttt{║◌◌◌│◌├─╢ 5} \\
\texttt{╟───┼─┘◌║ 6} \\
\texttt{║D◌T│P◌◌║ 7} \\
\texttt{╚═══╧═══╝} \\
\\ 
\texttt{```} \\
\\ 
\texttt{* Üst kenarda x koordinatları sağa doğru artar, sağ kenarda ise y koordinatları aşağı doğru artar.} \\
\texttt{* Aşağıdaki nesneler oyun tahtasına rastgele yerleştirilmiştir: 'P', 'O', 'T', 'S', 'D', 'A', 'M'.} \\
\\ 
\\ 
\texttt{Diğer oyuncu, nesnelerin farklı rastgele konumlara yerleştirildiği oyun tahtasının bir varyasyonunu görür. Sen diğer oyuncunun tahtasını göremezsin ve onlar da seninkini göremez.} \\
\texttt{**Amaç:**} \\
\\ 
\texttt{Her iki oyuncunun da, kendi oyun tahtalarındaki nesneleri hareket ettirerek aynı türden nesnelerin aynı koordinatlara gelmesini sağlaması gerekir. Ortak bir hedef düzen üzerinde anlaşmak için diğer oyuncuyla iletişim kurmalısınız.} \\
\\ 
\\ 
\texttt{**Kurallar:**} \\
\\ 
\texttt{* Her turda, aşağıdaki iki komuttan yalnızca birini gönderebilirsin:} \\
\texttt{ 1. `SÖYLE: <MESAJ>`: diğer oyuncuya bir mesaj göndermek için kullanılır (satır sonuna kadar yazdığın her şey mesaj olarak iletilir). Ben bu mesajı partnerine ileteceğim.} \\
\texttt{2. `HAMLE: <NESNE>, (<X>, <Y>)`: bir nesneyi yeni bir konuma taşımak için kullanılır. Burada <X> sütunu, <Y> ise satırı ifade eder. Yaptığın hamlenin başarılı olup olmadığını sana söyleyeceğim.} \\
\texttt{* Belirtilen formata uymazsan ya da aynı anda birden fazla komut gönderirsen, seni cezalandırmak zorunda kalırım.} \\
\texttt{* Her iki oyuncu da 16 sayısından fazla ceza alırsa, oyunu birlikte kaybedersiniz.} \\
\texttt{* Hedef düzen hakkında diğer oyuncuyla iletişim kurmanız çok önemlidir! Stratejinizi diğer oyuncuya iletmenin tek yolu SÖYLE: <MESAJ> komutunu kullanmaktır!} \\
\\ 
\\ 
\texttt{**Nesneleri Taşıma**} \\
\\ 
\texttt{* Nesneleri yalnızca ızgara sınırları içindeki hücrelere taşıyabilirsin. Hedef hücre boş olmalıdır; yani yalnızca ‘◌’ sembolünü içermelidir.} \\
\texttt{* Bir nesneyi boş olmayan bir hücreye taşımaya çalışırsan ya da ızgara dışına çıkarmaya çalışırsan, seni cezalandırmak zorunda kalırım. Ardından tekrar deneme hakkın olur.} \\
\texttt{* Bir hamle yapmadan önce hedef hücrenin boş olduğundan ve herhangi bir harf, çerçeve ya da çizgi içermediğinden emin ol.} \\
\\ 
\\ 
\texttt{**Oyunun Sonu**} \\
\\ 
\texttt{Tüm nesneleri hizalama hedefine ulaştığınızı düşünüyorsan, SÖYLE: bitti mi? göndererek diğer oyuncudan oyunu bitirmesini isteyebilirsin. Eğer diğer oyuncu senden oyunu bitirmeni isterse ve sen SÖYLE: bitti! ile yanıt verirsen, oyun sona erer.} \\
\\ 
\texttt{Oyun, 28 tur içinde sona ererse her iki oyuncu da kazanır. Bir tur, iki oyuncunun da geçerli bir komut göndermesi olarak tanımlanır.} \\
\\ 
\texttt{**Puanlama:**} \\
\\ 
\texttt{Aynı nesneler her iki oyun tahtasında ne kadar yakın konumlanırsa, o kadar fazla puan kazanırsınız. Cezalar puanınızı düşürür. Rekoru kırabilir misiniz?} \\
\\ 
\texttt{Diğer oyuncu şu mesajı göndererek oyunu başlattı:} \\
\\ 
\texttt{"Merhaba! Hedefimiz, her türden nesnenin aynı koordinatlarda toplanması. Öncelikle stratejimizi belirleyelim."} \\
\\ 
\texttt{İlk komutun ne?} \\
            }
        }
    }
    & \\ \\

    \theutterance \stepcounter{utterance}  
    & & & \multicolumn{4}{p{0.6\linewidth}}{
        \cellcolor[rgb]{0.9,0.9,0.9}{
            \makecell[{{p{\linewidth}}}]{
                \texttt{\tiny{[GM$\langle$P2]}}
                \texttt{SÖYLE: Hedefimiz, her türden nesneyi (P, O, T, S, D, A, M) tek bir koordinat üzerinde hizalamak. Öncelikle nesnelerin mevcut konumlarını ve hedef koordinatını paylaşacağız. Örnek bir düzen düşünelim: tüm nesneleri 4. sırada toplayalım. Stratejimiz, her nesneyi adım adım 4. satıra taşımak olabilir. İlk adım olarak S nesnesinin konumunu kontrol edelim.} \\
            }
        }
    }
     \\ \\

    \theutterance \stepcounter{utterance}  
    & & & \multicolumn{2}{p{0.3\linewidth}}{
        \cellcolor[rgb]{0.9,0.9,0.9}{
            \makecell[{{p{\linewidth}}}]{
                \texttt{\tiny{[GM$|$GM]}}
                \texttt{SÖYLE: Hedefimiz, her türden nesneyi (P, O, T, S, D, A, M) tek bir koordinat üzerinde hizalamak. Öncelikle nesnelerin mevcut konumlarını ve hedef koordinatını paylaşacağız. Örnek bir düzen düşünelim: tüm nesneleri 4. sırada toplayalım. Stratejimiz, her nesneyi adım adım 4. satıra taşımak olabilir. İlk adım olarak S nesnesinin konumunu kontrol edelim.} \\
            }
        }
    }
    & & \\ \\

    \theutterance \stepcounter{utterance}  
    & & \multicolumn{4}{p{0.6\linewidth}}{
        \cellcolor[rgb]{0.9,0.9,0.9}{
            \makecell[{{p{\linewidth}}}]{
                \texttt{\tiny{[P1$\langle$GM]}}
                \texttt{Mesajın diğer oyuncuya iletildi.} \\
\\ 
\\ 
\texttt{Şu anda en fazla 28 turdan 1. turu oynuyorsunuz.} \\
\texttt{Toplamda 16 cezanın 0 kadarını biriktirdiniz.Diğer oyuncu şu mesajı gönderdi:} \\
\\ 
\\ 
\texttt{"Hedefimiz, her türden nesneyi (P, O, T, S, D, A, M) tek bir koordinat üzerinde hizalamak. Öncelikle nesnelerin mevcut konumlarını ve hedef koordinatını paylaşacağız. Örnek bir düzen düşünelim: tüm nesneleri 4. sırada toplayalım. Stratejimiz, her nesneyi adım adım 4. satıra taşımak olabilir. İlk adım olarak S nesnesinin konumunu kontrol edelim."} \\
\\ 
\\ 
\texttt{Bir sonraki komutun nedir?} \\
            }
        }
    }
    & \\ \\

    \theutterance \stepcounter{utterance}  
    & \multicolumn{4}{p{0.6\linewidth}}{
        \cellcolor[rgb]{0.9,0.9,0.9}{
            \makecell[{{p{\linewidth}}}]{
                \texttt{\tiny{[P1$\rangle$GM]}}
                \texttt{HAMLE: S, (4, 3)} \\
            }
        }
    }
    & & \\ \\

    \theutterance \stepcounter{utterance}  
    & & & \multicolumn{2}{p{0.3\linewidth}}{
        \cellcolor[rgb]{0.9,0.9,0.9}{
            \makecell[{{p{\linewidth}}}]{
                \texttt{\tiny{[GM$|$GM]}}
                \texttt{HAMLE: S, (4, 3)} \\
            }
        }
    }
    & & \\ \\

    \theutterance \stepcounter{utterance}  
    & & & \multicolumn{2}{p{0.3\linewidth}}{
        \cellcolor[rgb]{0.9,0.9,0.9}{
            \makecell[{{p{\linewidth}}}]{
                \texttt{\tiny{[GM$|$GM]}}
                \texttt{Ceza: (4,3) boş değil; "D" içeriyor.} \\
\\ 
\texttt{Toplamda 16 cezanın 1 kadarını biriktirdiniz.Lütfen tekrar deneyin!} \\
            }
        }
    }
    & & \\ \\

    \theutterance \stepcounter{utterance}  
    & & \multicolumn{4}{p{0.6\linewidth}}{
        \cellcolor[rgb]{0.9,0.9,0.9}{
            \makecell[{{p{\linewidth}}}]{
                \texttt{\tiny{[P1$\langle$GM]}}
                \texttt{Ceza: (4,3) boş değil; "D" içeriyor.} \\
\\ 
\texttt{Toplamda 16 cezanın 1 kadarını biriktirdiniz.Lütfen tekrar deneyin!} \\
            }
        }
    }
    & \\ \\

    \theutterance \stepcounter{utterance}  
    & \multicolumn{4}{p{0.6\linewidth}}{
        \cellcolor[rgb]{0.9,0.9,0.9}{
            \makecell[{{p{\linewidth}}}]{
                \texttt{\tiny{[P1$\rangle$GM]}}
                \texttt{HAMLE: S, (4, 4)} \\
            }
        }
    }
    & & \\ \\

    \theutterance \stepcounter{utterance}  
    & & & \multicolumn{2}{p{0.3\linewidth}}{
        \cellcolor[rgb]{0.9,0.9,0.9}{
            \makecell[{{p{\linewidth}}}]{
                \texttt{\tiny{[GM$|$GM]}}
                \texttt{HAMLE: S, (4, 4)} \\
            }
        }
    }
    & & \\ \\

    \theutterance \stepcounter{utterance}  
    & & & \multicolumn{2}{p{0.3\linewidth}}{
        \cellcolor[rgb]{0.9,0.9,0.9}{
            \makecell[{{p{\linewidth}}}]{
                \texttt{\tiny{[GM$|$GM]}}
                \texttt{Ceza: (4,4) boş değil; "┌" içeriyor.} \\
\\ 
\texttt{Toplamda 16 cezanın 2 kadarını biriktirdiniz.Lütfen tekrar deneyin!} \\
            }
        }
    }
    & & \\ \\

    \theutterance \stepcounter{utterance}  
    & & \multicolumn{4}{p{0.6\linewidth}}{
        \cellcolor[rgb]{0.9,0.9,0.9}{
            \makecell[{{p{\linewidth}}}]{
                \texttt{\tiny{[P1$\langle$GM]}}
                \texttt{Ceza: (4,4) boş değil; "┌" içeriyor.} \\
\\ 
\texttt{Toplamda 16 cezanın 2 kadarını biriktirdiniz.Lütfen tekrar deneyin!} \\
            }
        }
    }
    & \\ \\

    \theutterance \stepcounter{utterance}  
    & \multicolumn{4}{p{0.6\linewidth}}{
        \cellcolor[rgb]{0.9,0.9,0.9}{
            \makecell[{{p{\linewidth}}}]{
                \texttt{\tiny{[P1$\rangle$GM]}}
                \texttt{HAMLE: S, (4, 5)} \\
            }
        }
    }
    & & \\ \\

    \theutterance \stepcounter{utterance}  
    & & & \multicolumn{2}{p{0.3\linewidth}}{
        \cellcolor[rgb]{0.9,0.9,0.9}{
            \makecell[{{p{\linewidth}}}]{
                \texttt{\tiny{[GM$|$GM]}}
                \texttt{HAMLE: S, (4, 5)} \\
            }
        }
    }
    & & \\ \\

    \theutterance \stepcounter{utterance}  
    & & & \multicolumn{2}{p{0.3\linewidth}}{
        \cellcolor[rgb]{0.9,0.9,0.9}{
            \makecell[{{p{\linewidth}}}]{
                \texttt{\tiny{[GM$|$GM]}}
                \texttt{Ceza: (4,5) boş değil; "│" içeriyor.} \\
\\ 
\texttt{Toplamda 16 cezanın 3 kadarını biriktirdiniz.Lütfen tekrar deneyin!} \\
            }
        }
    }
    & & \\ \\

    \theutterance \stepcounter{utterance}  
    & & \multicolumn{4}{p{0.6\linewidth}}{
        \cellcolor[rgb]{0.9,0.9,0.9}{
            \makecell[{{p{\linewidth}}}]{
                \texttt{\tiny{[P1$\langle$GM]}}
                \texttt{Ceza: (4,5) boş değil; "│" içeriyor.} \\
\\ 
\texttt{Toplamda 16 cezanın 3 kadarını biriktirdiniz.Lütfen tekrar deneyin!} \\
            }
        }
    }
    & \\ \\

    \theutterance \stepcounter{utterance}  
    & \multicolumn{4}{p{0.6\linewidth}}{
        \cellcolor[rgb]{0.9,0.9,0.9}{
            \makecell[{{p{\linewidth}}}]{
                \texttt{\tiny{[P1$\rangle$GM]}}
                \texttt{HAMLE: S, (4, 6)} \\
            }
        }
    }
    & & \\ \\

    \theutterance \stepcounter{utterance}  
    & & & \multicolumn{2}{p{0.3\linewidth}}{
        \cellcolor[rgb]{0.9,0.9,0.9}{
            \makecell[{{p{\linewidth}}}]{
                \texttt{\tiny{[GM$|$GM]}}
                \texttt{HAMLE: S, (4, 6)} \\
            }
        }
    }
    & & \\ \\

    \theutterance \stepcounter{utterance}  
    & & & \multicolumn{2}{p{0.3\linewidth}}{
        \cellcolor[rgb]{0.9,0.9,0.9}{
            \makecell[{{p{\linewidth}}}]{
                \texttt{\tiny{[GM$|$GM]}}
                \texttt{Ceza: (4,6) boş değil; "┼" içeriyor.} \\
\\ 
\texttt{Toplamda 16 cezanın 4 kadarını biriktirdiniz.Lütfen tekrar deneyin!} \\
            }
        }
    }
    & & \\ \\

    \theutterance \stepcounter{utterance}  
    & & \multicolumn{4}{p{0.6\linewidth}}{
        \cellcolor[rgb]{0.9,0.9,0.9}{
            \makecell[{{p{\linewidth}}}]{
                \texttt{\tiny{[P1$\langle$GM]}}
                \texttt{Ceza: (4,6) boş değil; "┼" içeriyor.} \\
\\ 
\texttt{Toplamda 16 cezanın 4 kadarını biriktirdiniz.Lütfen tekrar deneyin!} \\
            }
        }
    }
    & \\ \\

    \theutterance \stepcounter{utterance}  
    & \multicolumn{4}{p{0.6\linewidth}}{
        \cellcolor[rgb]{0.9,0.9,0.9}{
            \makecell[{{p{\linewidth}}}]{
                \texttt{\tiny{[P1$\rangle$GM]}}
                \texttt{HAMLE: S, (4, 7)} \\
            }
        }
    }
    & & \\ \\

    \theutterance \stepcounter{utterance}  
    & & & \multicolumn{2}{p{0.3\linewidth}}{
        \cellcolor[rgb]{0.9,0.9,0.9}{
            \makecell[{{p{\linewidth}}}]{
                \texttt{\tiny{[GM$|$GM]}}
                \texttt{HAMLE: S, (4, 7)} \\
            }
        }
    }
    & & \\ \\

    \theutterance \stepcounter{utterance}  
    & & & \multicolumn{2}{p{0.3\linewidth}}{
        \cellcolor[rgb]{0.9,0.9,0.9}{
            \makecell[{{p{\linewidth}}}]{
                \texttt{\tiny{[GM$|$GM]}}
                \texttt{Ceza: (4,7) boş değil; "│" içeriyor.} \\
\\ 
\texttt{Toplamda 16 cezanın 5 kadarını biriktirdiniz.Lütfen tekrar deneyin!} \\
            }
        }
    }
    & & \\ \\

    \theutterance \stepcounter{utterance}  
    & & \multicolumn{4}{p{0.6\linewidth}}{
        \cellcolor[rgb]{0.9,0.9,0.9}{
            \makecell[{{p{\linewidth}}}]{
                \texttt{\tiny{[P1$\langle$GM]}}
                \texttt{Ceza: (4,7) boş değil; "│" içeriyor.} \\
\\ 
\texttt{Toplamda 16 cezanın 5 kadarını biriktirdiniz.Lütfen tekrar deneyin!} \\
            }
        }
    }
    & \\ \\

    \theutterance \stepcounter{utterance}  
    & \multicolumn{4}{p{0.6\linewidth}}{
        \cellcolor[rgb]{0.9,0.9,0.9}{
            \makecell[{{p{\linewidth}}}]{
                \texttt{\tiny{[P1$\rangle$GM]}}
                \texttt{HAMLE: S, (4, 5)} \\
            }
        }
    }
    & & \\ \\

    \theutterance \stepcounter{utterance}  
    & & & \multicolumn{2}{p{0.3\linewidth}}{
        \cellcolor[rgb]{0.9,0.9,0.9}{
            \makecell[{{p{\linewidth}}}]{
                \texttt{\tiny{[GM$|$GM]}}
                \texttt{HAMLE: S, (4, 5)} \\
            }
        }
    }
    & & \\ \\

    \theutterance \stepcounter{utterance}  
    & & & \multicolumn{2}{p{0.3\linewidth}}{
        \cellcolor[rgb]{0.9,0.9,0.9}{
            \makecell[{{p{\linewidth}}}]{
                \texttt{\tiny{[GM$|$GM]}}
                \texttt{Ceza: (4,5) boş değil; "│" içeriyor.} \\
\\ 
\texttt{Toplamda 16 cezanın 6 kadarını biriktirdiniz.Lütfen tekrar deneyin!} \\
            }
        }
    }
    & & \\ \\

    \theutterance \stepcounter{utterance}  
    & & \multicolumn{4}{p{0.6\linewidth}}{
        \cellcolor[rgb]{0.9,0.9,0.9}{
            \makecell[{{p{\linewidth}}}]{
                \texttt{\tiny{[P1$\langle$GM]}}
                \texttt{Ceza: (4,5) boş değil; "│" içeriyor.} \\
\\ 
\texttt{Toplamda 16 cezanın 6 kadarını biriktirdiniz.Lütfen tekrar deneyin!} \\
            }
        }
    }
    & \\ \\

    \theutterance \stepcounter{utterance}  
    & \multicolumn{4}{p{0.6\linewidth}}{
        \cellcolor[rgb]{0.9,0.9,0.9}{
            \makecell[{{p{\linewidth}}}]{
                \texttt{\tiny{[P1$\rangle$GM]}}
                \texttt{HAMLE: S, (4, 3)} \\
            }
        }
    }
    & & \\ \\

    \theutterance \stepcounter{utterance}  
    & & & \multicolumn{2}{p{0.3\linewidth}}{
        \cellcolor[rgb]{0.9,0.9,0.9}{
            \makecell[{{p{\linewidth}}}]{
                \texttt{\tiny{[GM$|$GM]}}
                \texttt{HAMLE: S, (4, 3)} \\
            }
        }
    }
    & & \\ \\

    \theutterance \stepcounter{utterance}  
    & & & \multicolumn{2}{p{0.3\linewidth}}{
        \cellcolor[rgb]{0.9,0.9,0.9}{
            \makecell[{{p{\linewidth}}}]{
                \texttt{\tiny{[GM$|$GM]}}
                \texttt{Ceza: (4,3) boş değil; "D" içeriyor.} \\
\\ 
\texttt{Toplamda 16 cezanın 7 kadarını biriktirdiniz.Lütfen tekrar deneyin!} \\
            }
        }
    }
    & & \\ \\

    \theutterance \stepcounter{utterance}  
    & & \multicolumn{4}{p{0.6\linewidth}}{
        \cellcolor[rgb]{0.9,0.9,0.9}{
            \makecell[{{p{\linewidth}}}]{
                \texttt{\tiny{[P1$\langle$GM]}}
                \texttt{Ceza: (4,3) boş değil; "D" içeriyor.} \\
\\ 
\texttt{Toplamda 16 cezanın 7 kadarını biriktirdiniz.Lütfen tekrar deneyin!} \\
            }
        }
    }
    & \\ \\

    \theutterance \stepcounter{utterance}  
    & \multicolumn{4}{p{0.6\linewidth}}{
        \cellcolor[rgb]{0.9,0.9,0.9}{
            \makecell[{{p{\linewidth}}}]{
                \texttt{\tiny{[P1$\rangle$GM]}}
                \texttt{HAMLE: S, (4, 4)} \\
            }
        }
    }
    & & \\ \\

    \theutterance \stepcounter{utterance}  
    & & & \multicolumn{2}{p{0.3\linewidth}}{
        \cellcolor[rgb]{0.9,0.9,0.9}{
            \makecell[{{p{\linewidth}}}]{
                \texttt{\tiny{[GM$|$GM]}}
                \texttt{HAMLE: S, (4, 4)} \\
            }
        }
    }
    & & \\ \\

    \theutterance \stepcounter{utterance}  
    & & & \multicolumn{2}{p{0.3\linewidth}}{
        \cellcolor[rgb]{0.9,0.9,0.9}{
            \makecell[{{p{\linewidth}}}]{
                \texttt{\tiny{[GM$|$GM]}}
                \texttt{Ceza: (4,4) boş değil; "┌" içeriyor.} \\
\\ 
\texttt{Toplamda 16 cezanın 8 kadarını biriktirdiniz.Lütfen tekrar deneyin!} \\
            }
        }
    }
    & & \\ \\

    \theutterance \stepcounter{utterance}  
    & & \multicolumn{4}{p{0.6\linewidth}}{
        \cellcolor[rgb]{0.9,0.9,0.9}{
            \makecell[{{p{\linewidth}}}]{
                \texttt{\tiny{[P1$\langle$GM]}}
                \texttt{Ceza: (4,4) boş değil; "┌" içeriyor.} \\
\\ 
\texttt{Toplamda 16 cezanın 8 kadarını biriktirdiniz.Lütfen tekrar deneyin!} \\
            }
        }
    }
    & \\ \\

    \theutterance \stepcounter{utterance}  
    & \multicolumn{4}{p{0.6\linewidth}}{
        \cellcolor[rgb]{0.9,0.9,0.9}{
            \makecell[{{p{\linewidth}}}]{
                \texttt{\tiny{[P1$\rangle$GM]}}
                \texttt{HAMLE: S, (3, 4)} \\
            }
        }
    }
    & & \\ \\

    \theutterance \stepcounter{utterance}  
    & & & \multicolumn{2}{p{0.3\linewidth}}{
        \cellcolor[rgb]{0.9,0.9,0.9}{
            \makecell[{{p{\linewidth}}}]{
                \texttt{\tiny{[GM$|$GM]}}
                \texttt{HAMLE: S, (3, 4)} \\
            }
        }
    }
    & & \\ \\

    \theutterance \stepcounter{utterance}  
    & & & \multicolumn{2}{p{0.3\linewidth}}{
        \cellcolor[rgb]{0.9,0.9,0.9}{
            \makecell[{{p{\linewidth}}}]{
                \texttt{\tiny{[GM$|$GM]}}
                \texttt{S moved to (3, 4)} \\
            }
        }
    }
    & & \\ \\

    \theutterance \stepcounter{utterance}  
    & & \multicolumn{4}{p{0.6\linewidth}}{
        \cellcolor[rgb]{0.9,0.9,0.9}{
            \makecell[{{p{\linewidth}}}]{
                \texttt{\tiny{[GM$\rangle$P2]}}
                \texttt{Mesajın diğer oyuncuya iletildi.} \\
\\ 
\\ 
\texttt{Şu anda en fazla 28 turdan 2. turu oynuyorsunuz.} \\
\texttt{Toplamda 16 cezanın 8 kadarını biriktirdiniz.Diğer oyuncu kendi ızgarasında bir nesneyi hareket ettirdi.} \\
\\ 
\\ 
\texttt{Bir sonraki komutun nedir?} \\
            }
        }
    }
    & \\ \\

    \theutterance \stepcounter{utterance}  
    & & & \multicolumn{4}{p{0.6\linewidth}}{
        \cellcolor[rgb]{0.9,0.9,0.9}{
            \makecell[{{p{\linewidth}}}]{
                \texttt{\tiny{[GM$\langle$P2]}}
                \texttt{SÖYLE: S nesnesinin mevcut konumunu ve hedef konumunu kontrol edelim. Önce S'nin şu anda nerede olduğunu öğrenelim; ardından 4. satıra hareket ettirmek için doğru X ve Y değerlerini belirleyelim.} \\
            }
        }
    }
     \\ \\

    \theutterance \stepcounter{utterance}  
    & & & \multicolumn{2}{p{0.3\linewidth}}{
        \cellcolor[rgb]{0.9,0.9,0.9}{
            \makecell[{{p{\linewidth}}}]{
                \texttt{\tiny{[GM$|$GM]}}
                \texttt{SÖYLE: S nesnesinin mevcut konumunu ve hedef konumunu kontrol edelim. Önce S'nin şu anda nerede olduğunu öğrenelim; ardından 4. satıra hareket ettirmek için doğru X ve Y değerlerini belirleyelim.} \\
            }
        }
    }
    & & \\ \\

    \theutterance \stepcounter{utterance}  
    & & \multicolumn{4}{p{0.6\linewidth}}{
        \cellcolor[rgb]{0.9,0.9,0.9}{
            \makecell[{{p{\linewidth}}}]{
                \texttt{\tiny{[P1$\langle$GM]}}
                \texttt{S’ başarıyla (3,4) koordinatlarına taşındı. Güncellenmiş ızgaran şu şekilde:} \\
\texttt{```} \\
\texttt{ 1234567} \\
\texttt{╔═══════╗} \\
\texttt{║◌◌◌◌◌T◌║ 1} \\
\texttt{║P◌◌◌◌◌◌║ 2} \\
\texttt{║◌◌AD◌┌─╢ 3} \\
\texttt{║◌◌S┌─┤◌║ 4} \\
\texttt{║O◌◌│◌├─╢ 5} \\
\texttt{╟───┼─┘◌║ 6} \\
\texttt{║◌M◌│◌◌◌║ 7} \\
\texttt{╚═══╧═══╝} \\
\texttt{```} \\
\texttt{Şu anda en fazla 28 turdan 2. turu oynuyorsunuz.} \\
\texttt{Toplamda 16 cezanın 8 kadarını biriktirdiniz.Diğer oyuncu şu mesajı gönderdi:} \\
\\ 
\\ 
\texttt{"S nesnesinin mevcut konumunu ve hedef konumunu kontrol edelim. Önce S'nin şu anda nerede olduğunu öğrenelim; ardından 4. satıra hareket ettirmek için doğru X ve Y değerlerini belirleyelim."} \\
\\ 
\\ 
\texttt{Bir sonraki komutun nedir?} \\
            }
        }
    }
    & \\ \\

    \theutterance \stepcounter{utterance}  
    & \multicolumn{4}{p{0.6\linewidth}}{
        \cellcolor[rgb]{0.9,0.9,0.9}{
            \makecell[{{p{\linewidth}}}]{
                \texttt{\tiny{[P1$\rangle$GM]}}
                \texttt{SÖYLE: S'nin mevcut konumu, koordinat (3,4). Hedefimiz 4. sırada hizalama, bu nedenle S'yi 4. satırda uygun bir X'e taşımayı düşünüyorum.} \\
            }
        }
    }
    & & \\ \\

    \theutterance \stepcounter{utterance}  
    & & & \multicolumn{2}{p{0.3\linewidth}}{
        \cellcolor[rgb]{0.9,0.9,0.9}{
            \makecell[{{p{\linewidth}}}]{
                \texttt{\tiny{[GM$|$GM]}}
                \texttt{SÖYLE: S'nin mevcut konumu, koordinat (3,4). Hedefimiz 4. sırada hizalama, bu nedenle S'yi 4. satırda uygun bir X'e taşımayı düşünüyorum.} \\
            }
        }
    }
    & & \\ \\

    \theutterance \stepcounter{utterance}  
    & & \multicolumn{4}{p{0.6\linewidth}}{
        \cellcolor[rgb]{0.9,0.9,0.9}{
            \makecell[{{p{\linewidth}}}]{
                \texttt{\tiny{[GM$\rangle$P2]}}
                \texttt{Mesajın diğer oyuncuya iletildi.} \\
\\ 
\\ 
\texttt{Şu anda en fazla 28 turdan 3. turu oynuyorsunuz.} \\
\texttt{Toplamda 16 cezanın 8 kadarını biriktirdiniz.Diğer oyuncu şu mesajı gönderdi:} \\
\\ 
\\ 
\texttt{"S'nin mevcut konumu, koordinat (3,4). Hedefimiz 4. sırada hizalama, bu nedenle S'yi 4. satırda uygun bir X'e taşımayı düşünüyorum."} \\
\\ 
\\ 
\texttt{Bir sonraki komutun nedir?} \\
            }
        }
    }
    & \\ \\

    \theutterance \stepcounter{utterance}  
    & & & \multicolumn{4}{p{0.6\linewidth}}{
        \cellcolor[rgb]{0.9,0.9,0.9}{
            \makecell[{{p{\linewidth}}}]{
                \texttt{\tiny{[GM$\langle$P2]}}
                \texttt{HAMLE: S, (4, 4)} \\
            }
        }
    }
     \\ \\

    \theutterance \stepcounter{utterance}  
    & & & \multicolumn{2}{p{0.3\linewidth}}{
        \cellcolor[rgb]{0.9,0.9,0.9}{
            \makecell[{{p{\linewidth}}}]{
                \texttt{\tiny{[GM$|$GM]}}
                \texttt{HAMLE: S, (4, 4)} \\
            }
        }
    }
    & & \\ \\

    \theutterance \stepcounter{utterance}  
    & & & \multicolumn{2}{p{0.3\linewidth}}{
        \cellcolor[rgb]{0.9,0.9,0.9}{
            \makecell[{{p{\linewidth}}}]{
                \texttt{\tiny{[GM$|$GM]}}
                \texttt{Ceza: (4,4) boş değil; "┌" içeriyor.} \\
\\ 
\texttt{Toplamda 16 cezanın 9 kadarını biriktirdiniz.Lütfen tekrar deneyin!} \\
            }
        }
    }
    & & \\ \\

    \theutterance \stepcounter{utterance}  
    & & \multicolumn{4}{p{0.6\linewidth}}{
        \cellcolor[rgb]{0.9,0.9,0.9}{
            \makecell[{{p{\linewidth}}}]{
                \texttt{\tiny{[GM$\rangle$P2]}}
                \texttt{Ceza: (4,4) boş değil; "┌" içeriyor.} \\
\\ 
\texttt{Toplamda 16 cezanın 9 kadarını biriktirdiniz.Lütfen tekrar deneyin!} \\
            }
        }
    }
    & \\ \\

    \theutterance \stepcounter{utterance}  
    & & & \multicolumn{4}{p{0.6\linewidth}}{
        \cellcolor[rgb]{0.9,0.9,0.9}{
            \makecell[{{p{\linewidth}}}]{
                \texttt{\tiny{[GM$\langle$P2]}}
                \texttt{SÖYLE: S'nin mevcut konumu (3,4). Hedef X'i kontrol edelim; 4. satırda boş hücre olup olmadığını ve uygun X değerini belirleyelim.} \\
            }
        }
    }
     \\ \\

    \theutterance \stepcounter{utterance}  
    & & & \multicolumn{2}{p{0.3\linewidth}}{
        \cellcolor[rgb]{0.9,0.9,0.9}{
            \makecell[{{p{\linewidth}}}]{
                \texttt{\tiny{[GM$|$GM]}}
                \texttt{SÖYLE: S'nin mevcut konumu (3,4). Hedef X'i kontrol edelim; 4. satırda boş hücre olup olmadığını ve uygun X değerini belirleyelim.} \\
            }
        }
    }
    & & \\ \\

    \theutterance \stepcounter{utterance}  
    & & \multicolumn{4}{p{0.6\linewidth}}{
        \cellcolor[rgb]{0.9,0.9,0.9}{
            \makecell[{{p{\linewidth}}}]{
                \texttt{\tiny{[P1$\langle$GM]}}
                \texttt{Mesajın diğer oyuncuya iletildi.} \\
\\ 
\\ 
\texttt{Şu anda en fazla 28 turdan 3. turu oynuyorsunuz.} \\
\texttt{Toplamda 16 cezanın 9 kadarını biriktirdiniz.Diğer oyuncu şu mesajı gönderdi:} \\
\\ 
\\ 
\texttt{"S'nin mevcut konumu (3,4). Hedef X'i kontrol edelim; 4. satırda boş hücre olup olmadığını ve uygun X değerini belirleyelim."} \\
\\ 
\\ 
\texttt{Bir sonraki komutun nedir?} \\
            }
        }
    }
    & \\ \\

    \theutterance \stepcounter{utterance}  
    & \multicolumn{4}{p{0.6\linewidth}}{
        \cellcolor[rgb]{0.9,0.9,0.9}{
            \makecell[{{p{\linewidth}}}]{
                \texttt{\tiny{[P1$\rangle$GM]}}
                \texttt{HAMLE: S, (4, 4)} \\
            }
        }
    }
    & & \\ \\

    \theutterance \stepcounter{utterance}  
    & & & \multicolumn{2}{p{0.3\linewidth}}{
        \cellcolor[rgb]{0.9,0.9,0.9}{
            \makecell[{{p{\linewidth}}}]{
                \texttt{\tiny{[GM$|$GM]}}
                \texttt{HAMLE: S, (4, 4)} \\
            }
        }
    }
    & & \\ \\

    \theutterance \stepcounter{utterance}  
    & & & \multicolumn{2}{p{0.3\linewidth}}{
        \cellcolor[rgb]{0.9,0.9,0.9}{
            \makecell[{{p{\linewidth}}}]{
                \texttt{\tiny{[GM$|$GM]}}
                \texttt{Ceza: (4,4) boş değil; "┌" içeriyor.} \\
\\ 
\texttt{Toplamda 16 cezanın 10 kadarını biriktirdiniz.Lütfen tekrar deneyin!} \\
            }
        }
    }
    & & \\ \\

    \theutterance \stepcounter{utterance}  
    & & \multicolumn{4}{p{0.6\linewidth}}{
        \cellcolor[rgb]{0.9,0.9,0.9}{
            \makecell[{{p{\linewidth}}}]{
                \texttt{\tiny{[P1$\langle$GM]}}
                \texttt{Ceza: (4,4) boş değil; "┌" içeriyor.} \\
\\ 
\texttt{Toplamda 16 cezanın 10 kadarını biriktirdiniz.Lütfen tekrar deneyin!} \\
            }
        }
    }
    & \\ \\

    \theutterance \stepcounter{utterance}  
    & \multicolumn{4}{p{0.6\linewidth}}{
        \cellcolor[rgb]{0.9,0.9,0.9}{
            \makecell[{{p{\linewidth}}}]{
                \texttt{\tiny{[P1$\rangle$GM]}}
                \texttt{HAMLE: S, (4, 5)} \\
            }
        }
    }
    & & \\ \\

    \theutterance \stepcounter{utterance}  
    & & & \multicolumn{2}{p{0.3\linewidth}}{
        \cellcolor[rgb]{0.9,0.9,0.9}{
            \makecell[{{p{\linewidth}}}]{
                \texttt{\tiny{[GM$|$GM]}}
                \texttt{HAMLE: S, (4, 5)} \\
            }
        }
    }
    & & \\ \\

    \theutterance \stepcounter{utterance}  
    & & & \multicolumn{2}{p{0.3\linewidth}}{
        \cellcolor[rgb]{0.9,0.9,0.9}{
            \makecell[{{p{\linewidth}}}]{
                \texttt{\tiny{[GM$|$GM]}}
                \texttt{Ceza: (4,5) boş değil; "│" içeriyor.} \\
\\ 
\texttt{Toplamda 16 cezanın 11 kadarını biriktirdiniz.Lütfen tekrar deneyin!} \\
            }
        }
    }
    & & \\ \\

    \theutterance \stepcounter{utterance}  
    & & \multicolumn{4}{p{0.6\linewidth}}{
        \cellcolor[rgb]{0.9,0.9,0.9}{
            \makecell[{{p{\linewidth}}}]{
                \texttt{\tiny{[P1$\langle$GM]}}
                \texttt{Ceza: (4,5) boş değil; "│" içeriyor.} \\
\\ 
\texttt{Toplamda 16 cezanın 11 kadarını biriktirdiniz.Lütfen tekrar deneyin!} \\
            }
        }
    }
    & \\ \\

    \theutterance \stepcounter{utterance}  
    & \multicolumn{4}{p{0.6\linewidth}}{
        \cellcolor[rgb]{0.9,0.9,0.9}{
            \makecell[{{p{\linewidth}}}]{
                \texttt{\tiny{[P1$\rangle$GM]}}
                \texttt{HAMLE: S, (4, 6)} \\
            }
        }
    }
    & & \\ \\

    \theutterance \stepcounter{utterance}  
    & & & \multicolumn{2}{p{0.3\linewidth}}{
        \cellcolor[rgb]{0.9,0.9,0.9}{
            \makecell[{{p{\linewidth}}}]{
                \texttt{\tiny{[GM$|$GM]}}
                \texttt{HAMLE: S, (4, 6)} \\
            }
        }
    }
    & & \\ \\

    \theutterance \stepcounter{utterance}  
    & & & \multicolumn{2}{p{0.3\linewidth}}{
        \cellcolor[rgb]{0.9,0.9,0.9}{
            \makecell[{{p{\linewidth}}}]{
                \texttt{\tiny{[GM$|$GM]}}
                \texttt{Ceza: (4,6) boş değil; "┼" içeriyor.} \\
\\ 
\texttt{Toplamda 16 cezanın 12 kadarını biriktirdiniz.Lütfen tekrar deneyin!} \\
            }
        }
    }
    & & \\ \\

    \theutterance \stepcounter{utterance}  
    & & \multicolumn{4}{p{0.6\linewidth}}{
        \cellcolor[rgb]{0.9,0.9,0.9}{
            \makecell[{{p{\linewidth}}}]{
                \texttt{\tiny{[P1$\langle$GM]}}
                \texttt{Ceza: (4,6) boş değil; "┼" içeriyor.} \\
\\ 
\texttt{Toplamda 16 cezanın 12 kadarını biriktirdiniz.Lütfen tekrar deneyin!} \\
            }
        }
    }
    & \\ \\

    \theutterance \stepcounter{utterance}  
    & \multicolumn{4}{p{0.6\linewidth}}{
        \cellcolor[rgb]{0.9,0.9,0.9}{
            \makecell[{{p{\linewidth}}}]{
                \texttt{\tiny{[P1$\rangle$GM]}}
                \texttt{HAMLE: S, (4, 7)} \\
            }
        }
    }
    & & \\ \\

    \theutterance \stepcounter{utterance}  
    & & & \multicolumn{2}{p{0.3\linewidth}}{
        \cellcolor[rgb]{0.9,0.9,0.9}{
            \makecell[{{p{\linewidth}}}]{
                \texttt{\tiny{[GM$|$GM]}}
                \texttt{HAMLE: S, (4, 7)} \\
            }
        }
    }
    & & \\ \\

    \theutterance \stepcounter{utterance}  
    & & & \multicolumn{2}{p{0.3\linewidth}}{
        \cellcolor[rgb]{0.9,0.9,0.9}{
            \makecell[{{p{\linewidth}}}]{
                \texttt{\tiny{[GM$|$GM]}}
                \texttt{Ceza: (4,7) boş değil; "│" içeriyor.} \\
\\ 
\texttt{Toplamda 16 cezanın 13 kadarını biriktirdiniz.Lütfen tekrar deneyin!} \\
            }
        }
    }
    & & \\ \\

    \theutterance \stepcounter{utterance}  
    & & \multicolumn{4}{p{0.6\linewidth}}{
        \cellcolor[rgb]{0.9,0.9,0.9}{
            \makecell[{{p{\linewidth}}}]{
                \texttt{\tiny{[P1$\langle$GM]}}
                \texttt{Ceza: (4,7) boş değil; "│" içeriyor.} \\
\\ 
\texttt{Toplamda 16 cezanın 13 kadarını biriktirdiniz.Lütfen tekrar deneyin!} \\
            }
        }
    }
    & \\ \\

    \theutterance \stepcounter{utterance}  
    & \multicolumn{4}{p{0.6\linewidth}}{
        \cellcolor[rgb]{0.9,0.9,0.9}{
            \makecell[{{p{\linewidth}}}]{
                \texttt{\tiny{[P1$\rangle$GM]}}
                \texttt{HAMLE: S, (5, 4)} \\
            }
        }
    }
    & & \\ \\

    \theutterance \stepcounter{utterance}  
    & & & \multicolumn{2}{p{0.3\linewidth}}{
        \cellcolor[rgb]{0.9,0.9,0.9}{
            \makecell[{{p{\linewidth}}}]{
                \texttt{\tiny{[GM$|$GM]}}
                \texttt{HAMLE: S, (5, 4)} \\
            }
        }
    }
    & & \\ \\

    \theutterance \stepcounter{utterance}  
    & & & \multicolumn{2}{p{0.3\linewidth}}{
        \cellcolor[rgb]{0.9,0.9,0.9}{
            \makecell[{{p{\linewidth}}}]{
                \texttt{\tiny{[GM$|$GM]}}
                \texttt{Ceza: (5,4) boş değil; "─" içeriyor.} \\
\\ 
\texttt{Toplamda 16 cezanın 14 kadarını biriktirdiniz.Lütfen tekrar deneyin!} \\
            }
        }
    }
    & & \\ \\

    \theutterance \stepcounter{utterance}  
    & & \multicolumn{4}{p{0.6\linewidth}}{
        \cellcolor[rgb]{0.9,0.9,0.9}{
            \makecell[{{p{\linewidth}}}]{
                \texttt{\tiny{[P1$\langle$GM]}}
                \texttt{Ceza: (5,4) boş değil; "─" içeriyor.} \\
\\ 
\texttt{Toplamda 16 cezanın 14 kadarını biriktirdiniz.Lütfen tekrar deneyin!} \\
            }
        }
    }
    & \\ \\

    \theutterance \stepcounter{utterance}  
    & \multicolumn{4}{p{0.6\linewidth}}{
        \cellcolor[rgb]{0.9,0.9,0.9}{
            \makecell[{{p{\linewidth}}}]{
                \texttt{\tiny{[P1$\rangle$GM]}}
                \texttt{HAMLE: S, (5, 5)} \\
            }
        }
    }
    & & \\ \\

    \theutterance \stepcounter{utterance}  
    & & & \multicolumn{2}{p{0.3\linewidth}}{
        \cellcolor[rgb]{0.9,0.9,0.9}{
            \makecell[{{p{\linewidth}}}]{
                \texttt{\tiny{[GM$|$GM]}}
                \texttt{HAMLE: S, (5, 5)} \\
            }
        }
    }
    & & \\ \\

    \theutterance \stepcounter{utterance}  
    & & & \multicolumn{2}{p{0.3\linewidth}}{
        \cellcolor[rgb]{0.9,0.9,0.9}{
            \makecell[{{p{\linewidth}}}]{
                \texttt{\tiny{[GM$|$GM]}}
                \texttt{S moved to (5, 5)} \\
            }
        }
    }
    & & \\ \\

    \theutterance \stepcounter{utterance}  
    & & \multicolumn{4}{p{0.6\linewidth}}{
        \cellcolor[rgb]{0.9,0.9,0.9}{
            \makecell[{{p{\linewidth}}}]{
                \texttt{\tiny{[GM$\rangle$P2]}}
                \texttt{Mesajın diğer oyuncuya iletildi.} \\
\\ 
\\ 
\texttt{Şu anda en fazla 28 turdan 4. turu oynuyorsunuz.} \\
\texttt{Toplamda 16 cezanın 14 kadarını biriktirdiniz.Diğer oyuncu kendi ızgarasında bir nesneyi hareket ettirdi.} \\
\\ 
\\ 
\texttt{Bir sonraki komutun nedir?} \\
            }
        }
    }
    & \\ \\

    \theutterance \stepcounter{utterance}  
    & & & \multicolumn{4}{p{0.6\linewidth}}{
        \cellcolor[rgb]{0.9,0.9,0.9}{
            \makecell[{{p{\linewidth}}}]{
                \texttt{\tiny{[GM$\langle$P2]}}
                \texttt{SÖYLE: Hedef X değeri için 4. satırda (4,4) dışındaki bir X'i kontrol edelim. Örneğin, (4,2) veya (4,5) gibi boş hücreler olup olmadığını kontrol edelim ve uygun X ile S'yi 4. satıra taşıyabileceğimizi belirleyelim.} \\
            }
        }
    }
     \\ \\

    \theutterance \stepcounter{utterance}  
    & & & \multicolumn{2}{p{0.3\linewidth}}{
        \cellcolor[rgb]{0.9,0.9,0.9}{
            \makecell[{{p{\linewidth}}}]{
                \texttt{\tiny{[GM$|$GM]}}
                \texttt{SÖYLE: Hedef X değeri için 4. satırda (4,4) dışındaki bir X'i kontrol edelim. Örneğin, (4,2) veya (4,5) gibi boş hücreler olup olmadığını kontrol edelim ve uygun X ile S'yi 4. satıra taşıyabileceğimizi belirleyelim.} \\
            }
        }
    }
    & & \\ \\

    \theutterance \stepcounter{utterance}  
    & & \multicolumn{4}{p{0.6\linewidth}}{
        \cellcolor[rgb]{0.9,0.9,0.9}{
            \makecell[{{p{\linewidth}}}]{
                \texttt{\tiny{[P1$\langle$GM]}}
                \texttt{S’ başarıyla (5,5) koordinatlarına taşındı. Güncellenmiş ızgaran şu şekilde:} \\
\texttt{```} \\
\texttt{ 1234567} \\
\texttt{╔═══════╗} \\
\texttt{║◌◌◌◌◌T◌║ 1} \\
\texttt{║P◌◌◌◌◌◌║ 2} \\
\texttt{║◌◌AD◌┌─╢ 3} \\
\texttt{║◌◌◌┌─┤◌║ 4} \\
\texttt{║O◌◌│S├─╢ 5} \\
\texttt{╟───┼─┘◌║ 6} \\
\texttt{║◌M◌│◌◌◌║ 7} \\
\texttt{╚═══╧═══╝} \\
\texttt{```} \\
\texttt{Şu anda en fazla 28 turdan 4. turu oynuyorsunuz.} \\
\texttt{Toplamda 16 cezanın 14 kadarını biriktirdiniz.Diğer oyuncu şu mesajı gönderdi:} \\
\\ 
\\ 
\texttt{"Hedef X değeri için 4. satırda (4,4) dışındaki bir X'i kontrol edelim. Örneğin, (4,2) veya (4,5) gibi boş hücreler olup olmadığını kontrol edelim ve uygun X ile S'yi 4. satıra taşıyabileceğimizi belirleyelim."} \\
\\ 
\\ 
\texttt{Bir sonraki komutun nedir?} \\
            }
        }
    }
    & \\ \\

    \theutterance \stepcounter{utterance}  
    & \multicolumn{4}{p{0.6\linewidth}}{
        \cellcolor[rgb]{0.9,0.9,0.9}{
            \makecell[{{p{\linewidth}}}]{
                \texttt{\tiny{[P1$\rangle$GM]}}
                \texttt{HAMLE: S, (4, 2)} \\
            }
        }
    }
    & & \\ \\

    \theutterance \stepcounter{utterance}  
    & & & \multicolumn{2}{p{0.3\linewidth}}{
        \cellcolor[rgb]{0.9,0.9,0.9}{
            \makecell[{{p{\linewidth}}}]{
                \texttt{\tiny{[GM$|$GM]}}
                \texttt{HAMLE: S, (4, 2)} \\
            }
        }
    }
    & & \\ \\

    \theutterance \stepcounter{utterance}  
    & & & \multicolumn{2}{p{0.3\linewidth}}{
        \cellcolor[rgb]{0.9,0.9,0.9}{
            \makecell[{{p{\linewidth}}}]{
                \texttt{\tiny{[GM$|$GM]}}
                \texttt{S moved to (4, 2)} \\
            }
        }
    }
    & & \\ \\

    \theutterance \stepcounter{utterance}  
    & & \multicolumn{4}{p{0.6\linewidth}}{
        \cellcolor[rgb]{0.9,0.9,0.9}{
            \makecell[{{p{\linewidth}}}]{
                \texttt{\tiny{[GM$\rangle$P2]}}
                \texttt{Mesajın diğer oyuncuya iletildi.} \\
\\ 
\\ 
\texttt{Şu anda en fazla 28 turdan 5. turu oynuyorsunuz.} \\
\texttt{Toplamda 16 cezanın 14 kadarını biriktirdiniz.Diğer oyuncu kendi ızgarasında bir nesneyi hareket ettirdi.} \\
\\ 
\\ 
\texttt{Bir sonraki komutun nedir?} \\
            }
        }
    }
    & \\ \\

    \theutterance \stepcounter{utterance}  
    & & & \multicolumn{4}{p{0.6\linewidth}}{
        \cellcolor[rgb]{0.9,0.9,0.9}{
            \makecell[{{p{\linewidth}}}]{
                \texttt{\tiny{[GM$\langle$P2]}}
                \texttt{HAMLE: S, (4, 3)} \\
            }
        }
    }
     \\ \\

    \theutterance \stepcounter{utterance}  
    & & & \multicolumn{2}{p{0.3\linewidth}}{
        \cellcolor[rgb]{0.9,0.9,0.9}{
            \makecell[{{p{\linewidth}}}]{
                \texttt{\tiny{[GM$|$GM]}}
                \texttt{HAMLE: S, (4, 3)} \\
            }
        }
    }
    & & \\ \\

    \theutterance \stepcounter{utterance}  
    & & & \multicolumn{2}{p{0.3\linewidth}}{
        \cellcolor[rgb]{0.9,0.9,0.9}{
            \makecell[{{p{\linewidth}}}]{
                \texttt{\tiny{[GM$|$GM]}}
                \texttt{Ceza: (4,3) boş değil; "M" içeriyor.} \\
\\ 
\texttt{Toplamda 16 cezanın 15 kadarını biriktirdiniz.Lütfen tekrar deneyin!} \\
            }
        }
    }
    & & \\ \\

    \theutterance \stepcounter{utterance}  
    & & \multicolumn{4}{p{0.6\linewidth}}{
        \cellcolor[rgb]{0.9,0.9,0.9}{
            \makecell[{{p{\linewidth}}}]{
                \texttt{\tiny{[GM$\rangle$P2]}}
                \texttt{Ceza: (4,3) boş değil; "M" içeriyor.} \\
\\ 
\texttt{Toplamda 16 cezanın 15 kadarını biriktirdiniz.Lütfen tekrar deneyin!} \\
            }
        }
    }
    & \\ \\

    \theutterance \stepcounter{utterance}  
    & & & \multicolumn{4}{p{0.6\linewidth}}{
        \cellcolor[rgb]{0.9,0.9,0.9}{
            \makecell[{{p{\linewidth}}}]{
                \texttt{\tiny{[GM$\langle$P2]}}
                \texttt{SÖYLE: Hedef X değeri için 4. satırda (4, 3) dışındaki X'leri kontrol edelim. Örneğin, (4, 1), (4, 2), (4, 4), (4, 5), (4, 6) veya (4, 7) gibi boş hücreler olup olmadığını kontrol edelim ve uygun X ile S'yi 4. satıra taşıyabileceğimizi belirleyelim.} \\
            }
        }
    }
     \\ \\

    \theutterance \stepcounter{utterance}  
    & & & \multicolumn{2}{p{0.3\linewidth}}{
        \cellcolor[rgb]{0.9,0.9,0.9}{
            \makecell[{{p{\linewidth}}}]{
                \texttt{\tiny{[GM$|$GM]}}
                \texttt{SÖYLE: Hedef X değeri için 4. satırda (4, 3) dışındaki X'leri kontrol edelim. Örneğin, (4, 1), (4, 2), (4, 4), (4, 5), (4, 6) veya (4, 7) gibi boş hücreler olup olmadığını kontrol edelim ve uygun X ile S'yi 4. satıra taşıyabileceğimizi belirleyelim.} \\
            }
        }
    }
    & & \\ \\

    \theutterance \stepcounter{utterance}  
    & & \multicolumn{4}{p{0.6\linewidth}}{
        \cellcolor[rgb]{0.9,0.9,0.9}{
            \makecell[{{p{\linewidth}}}]{
                \texttt{\tiny{[P1$\langle$GM]}}
                \texttt{S’ başarıyla (4,2) koordinatlarına taşındı. Güncellenmiş ızgaran şu şekilde:} \\
\texttt{```} \\
\texttt{ 1234567} \\
\texttt{╔═══════╗} \\
\texttt{║◌◌◌◌◌T◌║ 1} \\
\texttt{║P◌◌S◌◌◌║ 2} \\
\texttt{║◌◌AD◌┌─╢ 3} \\
\texttt{║◌◌◌┌─┤◌║ 4} \\
\texttt{║O◌◌│◌├─╢ 5} \\
\texttt{╟───┼─┘◌║ 6} \\
\texttt{║◌M◌│◌◌◌║ 7} \\
\texttt{╚═══╧═══╝} \\
\texttt{```} \\
\texttt{Şu anda en fazla 28 turdan 5. turu oynuyorsunuz.} \\
\texttt{Toplamda 16 cezanın 15 kadarını biriktirdiniz.Diğer oyuncu şu mesajı gönderdi:} \\
\\ 
\\ 
\texttt{"Hedef X değeri için 4. satırda (4, 3) dışındaki X'leri kontrol edelim. Örneğin, (4, 1), (4, 2), (4, 4), (4, 5), (4, 6) veya (4, 7) gibi boş hücreler olup olmadığını kontrol edelim ve uygun X ile S'yi 4. satıra taşıyabileceğimizi belirleyelim."} \\
\\ 
\\ 
\texttt{Bir sonraki komutun nedir?} \\
            }
        }
    }
    & \\ \\

    \theutterance \stepcounter{utterance}  
    & \multicolumn{4}{p{0.6\linewidth}}{
        \cellcolor[rgb]{0.9,0.9,0.9}{
            \makecell[{{p{\linewidth}}}]{
                \texttt{\tiny{[P1$\rangle$GM]}}
                \texttt{HAMLE: S, (4, 4)} \\
            }
        }
    }
    & & \\ \\

    \theutterance \stepcounter{utterance}  
    & & & \multicolumn{2}{p{0.3\linewidth}}{
        \cellcolor[rgb]{0.9,0.9,0.9}{
            \makecell[{{p{\linewidth}}}]{
                \texttt{\tiny{[GM$|$GM]}}
                \texttt{HAMLE: S, (4, 4)} \\
            }
        }
    }
    & & \\ \\

    \theutterance \stepcounter{utterance}  
    & & & \multicolumn{2}{p{0.3\linewidth}}{
        \cellcolor[rgb]{0.9,0.9,0.9}{
            \makecell[{{p{\linewidth}}}]{
                \texttt{\tiny{[GM$|$GM]}}
                \texttt{Ceza: (4,4) boş değil; "┌" içeriyor.} \\
\\ 
\texttt{Toplamda 16 cezanın 16 kadarını biriktirdiniz.Lütfen tekrar deneyin!} \\
            }
        }
    }
    & & \\ \\

    \theutterance \stepcounter{utterance}  
    & & \multicolumn{4}{p{0.6\linewidth}}{
        \cellcolor[rgb]{0.9,0.9,0.9}{
            \makecell[{{p{\linewidth}}}]{
                \texttt{\tiny{[P1$\langle$GM]}}
                \texttt{Ceza: (4,4) boş değil; "┌" içeriyor.} \\
\\ 
\texttt{Toplamda 16 cezanın 16 kadarını biriktirdiniz.Lütfen tekrar deneyin!} \\
            }
        }
    }
    & \\ \\

    \theutterance \stepcounter{utterance}  
    & \multicolumn{4}{p{0.6\linewidth}}{
        \cellcolor[rgb]{0.9,0.9,0.9}{
            \makecell[{{p{\linewidth}}}]{
                \texttt{\tiny{[P1$\rangle$GM]}}
                \texttt{HAMLE: S, (4, 5)} \\
            }
        }
    }
    & & \\ \\

    \theutterance \stepcounter{utterance}  
    & & & \multicolumn{2}{p{0.3\linewidth}}{
        \cellcolor[rgb]{0.9,0.9,0.9}{
            \makecell[{{p{\linewidth}}}]{
                \texttt{\tiny{[GM$|$GM]}}
                \texttt{HAMLE: S, (4, 5)} \\
            }
        }
    }
    & & \\ \\

    \theutterance \stepcounter{utterance}  
    & & & \multicolumn{2}{p{0.3\linewidth}}{
        \cellcolor[rgb]{0.9,0.9,0.9}{
            \makecell[{{p{\linewidth}}}]{
                \texttt{\tiny{[GM$|$GM]}}
                \texttt{Ceza: (4,5) boş değil; "│" içeriyor.} \\
\\ 
\texttt{Toplamda 16 cezanın 17 kadarını biriktirdiniz.Lütfen tekrar deneyin!} \\
            }
        }
    }
    & & \\ \\

    \theutterance \stepcounter{utterance}  
    & & \multicolumn{4}{p{0.6\linewidth}}{
        \cellcolor[rgb]{0.9,0.9,0.9}{
            \makecell[{{p{\linewidth}}}]{
                \texttt{\tiny{[P1$\langle$GM]}}
                \texttt{Ceza: (4,5) boş değil; "│" içeriyor.} \\
\\ 
\texttt{Toplamda 16 cezanın 17 kadarını biriktirdiniz.Lütfen tekrar deneyin!} \\
            }
        }
    }
    & \\ \\

    \theutterance \stepcounter{utterance}  
    & \multicolumn{4}{p{0.6\linewidth}}{
        \cellcolor[rgb]{0.9,0.9,0.9}{
            \makecell[{{p{\linewidth}}}]{
                \texttt{\tiny{[P1$\rangle$GM]}}
                \texttt{HAMLE: S, (4, 6)} \\
            }
        }
    }
    & & \\ \\

    \theutterance \stepcounter{utterance}  
    & & & \multicolumn{2}{p{0.3\linewidth}}{
        \cellcolor[rgb]{0.9,0.9,0.9}{
            \makecell[{{p{\linewidth}}}]{
                \texttt{\tiny{[GM$|$GM]}}
                \texttt{HAMLE: S, (4, 6)} \\
            }
        }
    }
    & & \\ \\

    \theutterance \stepcounter{utterance}  
    & & & \multicolumn{2}{p{0.3\linewidth}}{
        \cellcolor[rgb]{0.9,0.9,0.9}{
            \makecell[{{p{\linewidth}}}]{
                \texttt{\tiny{[GM$|$GM]}}
                \texttt{Ceza: (4,6) boş değil; "┼" içeriyor.} \\
\\ 
\texttt{Toplamda 16 cezanın 18 kadarını biriktirdiniz.Lütfen tekrar deneyin!} \\
            }
        }
    }
    & & \\ \\

    \theutterance \stepcounter{utterance}  
    & & \multicolumn{4}{p{0.6\linewidth}}{
        \cellcolor[rgb]{0.9,0.9,0.9}{
            \makecell[{{p{\linewidth}}}]{
                \texttt{\tiny{[P1$\langle$GM]}}
                \texttt{Ceza: (4,6) boş değil; "┼" içeriyor.} \\
\\ 
\texttt{Toplamda 16 cezanın 18 kadarını biriktirdiniz.Lütfen tekrar deneyin!} \\
            }
        }
    }
    & \\ \\

    \theutterance \stepcounter{utterance}  
    & \multicolumn{4}{p{0.6\linewidth}}{
        \cellcolor[rgb]{0.9,0.9,0.9}{
            \makecell[{{p{\linewidth}}}]{
                \texttt{\tiny{[P1$\rangle$GM]}}
                \texttt{HAMLE: S, (4, 7)} \\
            }
        }
    }
    & & \\ \\

    \theutterance \stepcounter{utterance}  
    & & & \multicolumn{2}{p{0.3\linewidth}}{
        \cellcolor[rgb]{0.9,0.9,0.9}{
            \makecell[{{p{\linewidth}}}]{
                \texttt{\tiny{[GM$|$GM]}}
                \texttt{HAMLE: S, (4, 7)} \\
            }
        }
    }
    & & \\ \\

    \theutterance \stepcounter{utterance}  
    & & & \multicolumn{2}{p{0.3\linewidth}}{
        \cellcolor[rgb]{0.9,0.9,0.9}{
            \makecell[{{p{\linewidth}}}]{
                \texttt{\tiny{[GM$|$GM]}}
                \texttt{Ceza: (4,7) boş değil; "│" içeriyor.} \\
\\ 
\texttt{Toplamda 16 cezanın 19 kadarını biriktirdiniz.Lütfen tekrar deneyin!} \\
            }
        }
    }
    & & \\ \\

    \theutterance \stepcounter{utterance}  
    & & \multicolumn{4}{p{0.6\linewidth}}{
        \cellcolor[rgb]{0.9,0.9,0.9}{
            \makecell[{{p{\linewidth}}}]{
                \texttt{\tiny{[P1$\langle$GM]}}
                \texttt{Ceza: (4,7) boş değil; "│" içeriyor.} \\
\\ 
\texttt{Toplamda 16 cezanın 19 kadarını biriktirdiniz.Lütfen tekrar deneyin!} \\
            }
        }
    }
    & \\ \\

    \theutterance \stepcounter{utterance}  
    & \multicolumn{4}{p{0.6\linewidth}}{
        \cellcolor[rgb]{0.9,0.9,0.9}{
            \makecell[{{p{\linewidth}}}]{
                \texttt{\tiny{[P1$\rangle$GM]}}
                \texttt{HAMLE: S, (3, 4)} \\
            }
        }
    }
    & & \\ \\

    \theutterance \stepcounter{utterance}  
    & & & \multicolumn{2}{p{0.3\linewidth}}{
        \cellcolor[rgb]{0.9,0.9,0.9}{
            \makecell[{{p{\linewidth}}}]{
                \texttt{\tiny{[GM$|$GM]}}
                \texttt{HAMLE: S, (3, 4)} \\
            }
        }
    }
    & & \\ \\

    \theutterance \stepcounter{utterance}  
    & & & \multicolumn{2}{p{0.3\linewidth}}{
        \cellcolor[rgb]{0.9,0.9,0.9}{
            \makecell[{{p{\linewidth}}}]{
                \texttt{\tiny{[GM$|$GM]}}
                \texttt{S moved to (3, 4)} \\
            }
        }
    }
    & & \\ \\

    \theutterance \stepcounter{utterance}  
    & & \multicolumn{4}{p{0.6\linewidth}}{
        \cellcolor[rgb]{0.9,0.9,0.9}{
            \makecell[{{p{\linewidth}}}]{
                \texttt{\tiny{[GM$\rangle$P2]}}
                \texttt{Mesajın diğer oyuncuya iletildi.} \\
\\ 
\\ 
\texttt{Şu anda en fazla 28 turdan 6. turu oynuyorsunuz.} \\
\texttt{Toplamda 16 cezanın 19 kadarını biriktirdiniz.Diğer oyuncu kendi ızgarasında bir nesneyi hareket ettirdi.} \\
\\ 
\\ 
\texttt{Bir sonraki komutun nedir?} \\
            }
        }
    }
    & \\ \\

    \theutterance \stepcounter{utterance}  
    & & & \multicolumn{4}{p{0.6\linewidth}}{
        \cellcolor[rgb]{0.9,0.9,0.9}{
            \makecell[{{p{\linewidth}}}]{
                \texttt{\tiny{[GM$\langle$P2]}}
                \texttt{HAMLE: S, (4, 1)} \\
            }
        }
    }
     \\ \\

    \theutterance \stepcounter{utterance}  
    & & & \multicolumn{2}{p{0.3\linewidth}}{
        \cellcolor[rgb]{0.9,0.9,0.9}{
            \makecell[{{p{\linewidth}}}]{
                \texttt{\tiny{[GM$|$GM]}}
                \texttt{HAMLE: S, (4, 1)} \\
            }
        }
    }
    & & \\ \\

    \theutterance \stepcounter{utterance}  
    & & & \multicolumn{2}{p{0.3\linewidth}}{
        \cellcolor[rgb]{0.9,0.9,0.9}{
            \makecell[{{p{\linewidth}}}]{
                \texttt{\tiny{[GM$|$GM]}}
                \texttt{S moved to (4, 1)} \\
            }
        }
    }
    & & \\ \\

    \theutterance \stepcounter{utterance}  
    & & \multicolumn{4}{p{0.6\linewidth}}{
        \cellcolor[rgb]{0.9,0.9,0.9}{
            \makecell[{{p{\linewidth}}}]{
                \texttt{\tiny{[P1$\langle$GM]}}
                \texttt{S’ başarıyla (3,4) koordinatlarına taşındı. Güncellenmiş ızgaran şu şekilde:} \\
\texttt{```} \\
\texttt{ 1234567} \\
\texttt{╔═══════╗} \\
\texttt{║◌◌◌◌◌T◌║ 1} \\
\texttt{║P◌◌◌◌◌◌║ 2} \\
\texttt{║◌◌AD◌┌─╢ 3} \\
\texttt{║◌◌S┌─┤◌║ 4} \\
\texttt{║O◌◌│◌├─╢ 5} \\
\texttt{╟───┼─┘◌║ 6} \\
\texttt{║◌M◌│◌◌◌║ 7} \\
\texttt{╚═══╧═══╝} \\
\texttt{```} \\
\texttt{Şu anda en fazla 28 turdan 6. turu oynuyorsunuz.} \\
\texttt{Toplamda 16 cezanın 19 kadarını biriktirdiniz.Diğer oyuncu kendi ızgarasında bir nesneyi hareket ettirdi.} \\
\\ 
\\ 
\texttt{Bir sonraki komutun nedir?} \\
            }
        }
    }
    & \\ \\

    \theutterance \stepcounter{utterance}  
    & \multicolumn{4}{p{0.6\linewidth}}{
        \cellcolor[rgb]{0.9,0.9,0.9}{
            \makecell[{{p{\linewidth}}}]{
                \texttt{\tiny{[P1$\rangle$GM]}}
                \texttt{SÖYLE: S’nin mevcut konumu (3,4). Hedefimiz S’yi 4. satıra taşımak. 4. satırın boş hücrelerini kontrol edelim, örn. (4,2), (4,3), etc.} \\
            }
        }
    }
    & & \\ \\

    \theutterance \stepcounter{utterance}  
    & & & \multicolumn{2}{p{0.3\linewidth}}{
        \cellcolor[rgb]{0.9,0.9,0.9}{
            \makecell[{{p{\linewidth}}}]{
                \texttt{\tiny{[GM$|$GM]}}
                \texttt{SÖYLE: S’nin mevcut konumu (3,4). Hedefimiz S’yi 4. satıra taşımak. 4. satırın boş hücrelerini kontrol edelim, örn. (4,2), (4,3), etc.} \\
            }
        }
    }
    & & \\ \\

    \theutterance \stepcounter{utterance}  
    & & \multicolumn{4}{p{0.6\linewidth}}{
        \cellcolor[rgb]{0.9,0.9,0.9}{
            \makecell[{{p{\linewidth}}}]{
                \texttt{\tiny{[GM$\rangle$P2]}}
                \texttt{S’ başarıyla (4,1) koordinatlarına taşındı. Güncellenmiş ızgaran şu şekilde:} \\
\texttt{```} \\
\texttt{ 1234567} \\
\texttt{╔═══════╗} \\
\texttt{║◌◌◌S◌◌◌║ 1} \\
\texttt{║◌◌◌◌O◌◌║ 2} \\
\texttt{║◌◌◌M◌┌─╢ 3} \\
\texttt{║◌◌A┌─┤◌║ 4} \\
\texttt{║◌◌◌│◌├─╢ 5} \\
\texttt{╟───┼─┘◌║ 6} \\
\texttt{║D◌T│P◌◌║ 7} \\
\texttt{╚═══╧═══╝} \\
\texttt{```} \\
\texttt{Şu anda en fazla 28 turdan 7. turu oynuyorsunuz.} \\
\texttt{Toplamda 16 cezanın 19 kadarını biriktirdiniz.Diğer oyuncu şu mesajı gönderdi:} \\
\\ 
\\ 
\texttt{"S’nin mevcut konumu (3,4). Hedefimiz S’yi 4. satıra taşımak. 4. satırın boş hücrelerini kontrol edelim, örn. (4,2), (4,3), etc."} \\
\\ 
\\ 
\texttt{Bir sonraki komutun nedir?} \\
            }
        }
    }
    & \\ \\

    \theutterance \stepcounter{utterance}  
    & & & \multicolumn{4}{p{0.6\linewidth}}{
        \cellcolor[rgb]{0.9,0.9,0.9}{
            \makecell[{{p{\linewidth}}}]{
                \texttt{\tiny{[GM$\langle$P2]}}
                \texttt{SÖYLE: 4. satırda boş hücreler: (4,2) ve (4,3) olup olmadığını kontrol edelim. S’yi (4,2) veya (4,3) konumuna taşıyabiliriz. Öncelikle (4,2) hücresinin hangi nesneyi içerdiğini kontrol edelim.} \\
            }
        }
    }
     \\ \\

    \theutterance \stepcounter{utterance}  
    & & & \multicolumn{2}{p{0.3\linewidth}}{
        \cellcolor[rgb]{0.9,0.9,0.9}{
            \makecell[{{p{\linewidth}}}]{
                \texttt{\tiny{[GM$|$GM]}}
                \texttt{SÖYLE: 4. satırda boş hücreler: (4,2) ve (4,3) olup olmadığını kontrol edelim. S’yi (4,2) veya (4,3) konumuna taşıyabiliriz. Öncelikle (4,2) hücresinin hangi nesneyi içerdiğini kontrol edelim.} \\
            }
        }
    }
    & & \\ \\

    \theutterance \stepcounter{utterance}  
    & & \multicolumn{4}{p{0.6\linewidth}}{
        \cellcolor[rgb]{0.9,0.9,0.9}{
            \makecell[{{p{\linewidth}}}]{
                \texttt{\tiny{[P1$\langle$GM]}}
                \texttt{Mesajın diğer oyuncuya iletildi.} \\
\\ 
\\ 
\texttt{Şu anda en fazla 28 turdan 7. turu oynuyorsunuz.} \\
\texttt{Toplamda 16 cezanın 19 kadarını biriktirdiniz.Diğer oyuncu şu mesajı gönderdi:} \\
\\ 
\\ 
\texttt{"4. satırda boş hücreler: (4,2) ve (4,3) olup olmadığını kontrol edelim. S’yi (4,2) veya (4,3) konumuna taşıyabiliriz. Öncelikle (4,2) hücresinin hangi nesneyi içerdiğini kontrol edelim."} \\
\\ 
\\ 
\texttt{Bir sonraki komutun nedir?} \\
            }
        }
    }
    & \\ \\

    \theutterance \stepcounter{utterance}  
    & \multicolumn{4}{p{0.6\linewidth}}{
        \cellcolor[rgb]{0.9,0.9,0.9}{
            \makecell[{{p{\linewidth}}}]{
                \texttt{\tiny{[P1$\rangle$GM]}}
                \texttt{HAMLE: S, (4,2)} \\
            }
        }
    }
    & & \\ \\

    \theutterance \stepcounter{utterance}  
    & & & \multicolumn{2}{p{0.3\linewidth}}{
        \cellcolor[rgb]{0.9,0.9,0.9}{
            \makecell[{{p{\linewidth}}}]{
                \texttt{\tiny{[GM$|$GM]}}
                \texttt{HAMLE: S, (4,2)} \\
            }
        }
    }
    & & \\ \\

    \theutterance \stepcounter{utterance}  
    & & & \multicolumn{2}{p{0.3\linewidth}}{
        \cellcolor[rgb]{0.9,0.9,0.9}{
            \makecell[{{p{\linewidth}}}]{
                \texttt{\tiny{[GM$|$GM]}}
                \texttt{S moved to (4, 2)} \\
            }
        }
    }
    & & \\ \\

    \theutterance \stepcounter{utterance}  
    & & \multicolumn{4}{p{0.6\linewidth}}{
        \cellcolor[rgb]{0.9,0.9,0.9}{
            \makecell[{{p{\linewidth}}}]{
                \texttt{\tiny{[GM$\rangle$P2]}}
                \texttt{Mesajın diğer oyuncuya iletildi.} \\
\\ 
\\ 
\texttt{Şu anda en fazla 28 turdan 8. turu oynuyorsunuz.} \\
\texttt{Toplamda 16 cezanın 19 kadarını biriktirdiniz.Diğer oyuncu kendi ızgarasında bir nesneyi hareket ettirdi.} \\
\\ 
\\ 
\texttt{Bir sonraki komutun nedir?} \\
            }
        }
    }
    & \\ \\

    \theutterance \stepcounter{utterance}  
    & & & \multicolumn{4}{p{0.6\linewidth}}{
        \cellcolor[rgb]{0.9,0.9,0.9}{
            \makecell[{{p{\linewidth}}}]{
                \texttt{\tiny{[GM$\langle$P2]}}
                \texttt{HAMLE: S, (4, 2)} \\
            }
        }
    }
     \\ \\

    \theutterance \stepcounter{utterance}  
    & & & \multicolumn{2}{p{0.3\linewidth}}{
        \cellcolor[rgb]{0.9,0.9,0.9}{
            \makecell[{{p{\linewidth}}}]{
                \texttt{\tiny{[GM$|$GM]}}
                \texttt{HAMLE: S, (4, 2)} \\
            }
        }
    }
    & & \\ \\

    \theutterance \stepcounter{utterance}  
    & & & \multicolumn{2}{p{0.3\linewidth}}{
        \cellcolor[rgb]{0.9,0.9,0.9}{
            \makecell[{{p{\linewidth}}}]{
                \texttt{\tiny{[GM$|$GM]}}
                \texttt{S moved to (4, 2)} \\
            }
        }
    }
    & & \\ \\

    \theutterance \stepcounter{utterance}  
    & & \multicolumn{4}{p{0.6\linewidth}}{
        \cellcolor[rgb]{0.9,0.9,0.9}{
            \makecell[{{p{\linewidth}}}]{
                \texttt{\tiny{[P1$\langle$GM]}}
                \texttt{S’ başarıyla (4,2) koordinatlarına taşındı. Güncellenmiş ızgaran şu şekilde:} \\
\texttt{```} \\
\texttt{ 1234567} \\
\texttt{╔═══════╗} \\
\texttt{║◌◌◌◌◌T◌║ 1} \\
\texttt{║P◌◌S◌◌◌║ 2} \\
\texttt{║◌◌AD◌┌─╢ 3} \\
\texttt{║◌◌◌┌─┤◌║ 4} \\
\texttt{║O◌◌│◌├─╢ 5} \\
\texttt{╟───┼─┘◌║ 6} \\
\texttt{║◌M◌│◌◌◌║ 7} \\
\texttt{╚═══╧═══╝} \\
\texttt{```} \\
\texttt{Şu anda en fazla 28 turdan 8. turu oynuyorsunuz.} \\
\texttt{Toplamda 16 cezanın 19 kadarını biriktirdiniz.Diğer oyuncu kendi ızgarasında bir nesneyi hareket ettirdi.} \\
\\ 
\\ 
\texttt{Bir sonraki komutun nedir?} \\
            }
        }
    }
    & \\ \\

    \theutterance \stepcounter{utterance}  
    & \multicolumn{4}{p{0.6\linewidth}}{
        \cellcolor[rgb]{0.9,0.9,0.9}{
            \makecell[{{p{\linewidth}}}]{
                \texttt{\tiny{[P1$\rangle$GM]}}
                \texttt{SÖYLE: S’yi şu an (4,2) konumunda, hedefimiz 4. satırda hizalama. Teşekkürler!} \\
            }
        }
    }
    & & \\ \\

    \theutterance \stepcounter{utterance}  
    & & & \multicolumn{2}{p{0.3\linewidth}}{
        \cellcolor[rgb]{0.9,0.9,0.9}{
            \makecell[{{p{\linewidth}}}]{
                \texttt{\tiny{[GM$|$GM]}}
                \texttt{SÖYLE: S’yi şu an (4,2) konumunda, hedefimiz 4. satırda hizalama. Teşekkürler!} \\
            }
        }
    }
    & & \\ \\

    \theutterance \stepcounter{utterance}  
    & & \multicolumn{4}{p{0.6\linewidth}}{
        \cellcolor[rgb]{0.9,0.9,0.9}{
            \makecell[{{p{\linewidth}}}]{
                \texttt{\tiny{[GM$\rangle$P2]}}
                \texttt{S’ başarıyla (4,2) koordinatlarına taşındı. Güncellenmiş ızgaran şu şekilde:} \\
\texttt{```} \\
\texttt{ 1234567} \\
\texttt{╔═══════╗} \\
\texttt{║◌◌◌◌◌◌◌║ 1} \\
\texttt{║◌◌◌SO◌◌║ 2} \\
\texttt{║◌◌◌M◌┌─╢ 3} \\
\texttt{║◌◌A┌─┤◌║ 4} \\
\texttt{║◌◌◌│◌├─╢ 5} \\
\texttt{╟───┼─┘◌║ 6} \\
\texttt{║D◌T│P◌◌║ 7} \\
\texttt{╚═══╧═══╝} \\
\texttt{```} \\
\texttt{Şu anda en fazla 28 turdan 9. turu oynuyorsunuz.} \\
\texttt{Toplamda 16 cezanın 19 kadarını biriktirdiniz.Diğer oyuncu şu mesajı gönderdi:} \\
\\ 
\\ 
\texttt{"S’yi şu an (4,2) konumunda, hedefimiz 4. satırda hizalama. Teşekkürler!"} \\
\\ 
\\ 
\texttt{Bir sonraki komutun nedir?} \\
            }
        }
    }
    & \\ \\

    \theutterance \stepcounter{utterance}  
    & & & \multicolumn{4}{p{0.6\linewidth}}{
        \cellcolor[rgb]{0.9,0.9,0.9}{
            \makecell[{{p{\linewidth}}}]{
                \texttt{\tiny{[GM$\langle$P2]}}
                \texttt{SÖYLE: Hücrelerinizi gözden geçirelim. 4. satırda hangi boş hücreler kaldı? Örn. (4,5) veya (4,6) veya (4,7). Öncelikle 4. satırda tamamı kapalı mı, yoksa başka taşınabilir hücreler mi var.} \\
            }
        }
    }
     \\ \\

    \theutterance \stepcounter{utterance}  
    & & & \multicolumn{2}{p{0.3\linewidth}}{
        \cellcolor[rgb]{0.9,0.9,0.9}{
            \makecell[{{p{\linewidth}}}]{
                \texttt{\tiny{[GM$|$GM]}}
                \texttt{SÖYLE: Hücrelerinizi gözden geçirelim. 4. satırda hangi boş hücreler kaldı? Örn. (4,5) veya (4,6) veya (4,7). Öncelikle 4. satırda tamamı kapalı mı, yoksa başka taşınabilir hücreler mi var.} \\
            }
        }
    }
    & & \\ \\

    \theutterance \stepcounter{utterance}  
    & & \multicolumn{4}{p{0.6\linewidth}}{
        \cellcolor[rgb]{0.9,0.9,0.9}{
            \makecell[{{p{\linewidth}}}]{
                \texttt{\tiny{[P1$\langle$GM]}}
                \texttt{Mesajın diğer oyuncuya iletildi.} \\
\\ 
\\ 
\texttt{Şu anda en fazla 28 turdan 9. turu oynuyorsunuz.} \\
\texttt{Toplamda 16 cezanın 19 kadarını biriktirdiniz.Diğer oyuncu şu mesajı gönderdi:} \\
\\ 
\\ 
\texttt{"Hücrelerinizi gözden geçirelim. 4. satırda hangi boş hücreler kaldı? Örn. (4,5) veya (4,6) veya (4,7). Öncelikle 4. satırda tamamı kapalı mı, yoksa başka taşınabilir hücreler mi var."} \\
\\ 
\\ 
\texttt{Bir sonraki komutun nedir?} \\
            }
        }
    }
    & \\ \\

    \theutterance \stepcounter{utterance}  
    & \multicolumn{4}{p{0.6\linewidth}}{
        \cellcolor[rgb]{0.9,0.9,0.9}{
            \makecell[{{p{\linewidth}}}]{
                \texttt{\tiny{[P1$\rangle$GM]}}
                \texttt{SÖYLE: 4. satırda (4,5), (4,6) ve (4,7) hücrelerinin hepsinin kapalı olduğunu kontrol ediyorum. Taşınabilir boş hücre olmadığını düşünüyorum.} \\
            }
        }
    }
    & & \\ \\

    \theutterance \stepcounter{utterance}  
    & & & \multicolumn{2}{p{0.3\linewidth}}{
        \cellcolor[rgb]{0.9,0.9,0.9}{
            \makecell[{{p{\linewidth}}}]{
                \texttt{\tiny{[GM$|$GM]}}
                \texttt{SÖYLE: 4. satırda (4,5), (4,6) ve (4,7) hücrelerinin hepsinin kapalı olduğunu kontrol ediyorum. Taşınabilir boş hücre olmadığını düşünüyorum.} \\
            }
        }
    }
    & & \\ \\

    \theutterance \stepcounter{utterance}  
    & & \multicolumn{4}{p{0.6\linewidth}}{
        \cellcolor[rgb]{0.9,0.9,0.9}{
            \makecell[{{p{\linewidth}}}]{
                \texttt{\tiny{[GM$\rangle$P2]}}
                \texttt{Mesajın diğer oyuncuya iletildi.} \\
\\ 
\\ 
\texttt{Şu anda en fazla 28 turdan 10. turu oynuyorsunuz.} \\
\texttt{Toplamda 16 cezanın 19 kadarını biriktirdiniz.Diğer oyuncu şu mesajı gönderdi:} \\
\\ 
\\ 
\texttt{"4. satırda (4,5), (4,6) ve (4,7) hücrelerinin hepsinin kapalı olduğunu kontrol ediyorum. Taşınabilir boş hücre olmadığını düşünüyorum."} \\
\\ 
\\ 
\texttt{Bir sonraki komutun nedir?} \\
            }
        }
    }
    & \\ \\

    \theutterance \stepcounter{utterance}  
    & & & \multicolumn{4}{p{0.6\linewidth}}{
        \cellcolor[rgb]{0.9,0.9,0.9}{
            \makecell[{{p{\linewidth}}}]{
                \texttt{\tiny{[GM$\langle$P2]}}
                \texttt{SÖYLE: Teşekkürler! Stratejimizi gözden geçirelim. S’nin şu anda (4,2) olduğu için, diğer nesneleri 4. satıra doğru hareket ettirebilir miyiz? Örn. O, A, M, D, T, P gibi nesneleri sırayla 4. satıra taşıyabiliriz. Öncelikle O'nun mevcut konumunu kontrol edelim.} \\
            }
        }
    }
     \\ \\

    \theutterance \stepcounter{utterance}  
    & & & \multicolumn{2}{p{0.3\linewidth}}{
        \cellcolor[rgb]{0.9,0.9,0.9}{
            \makecell[{{p{\linewidth}}}]{
                \texttt{\tiny{[GM$|$GM]}}
                \texttt{SÖYLE: Teşekkürler! Stratejimizi gözden geçirelim. S’nin şu anda (4,2) olduğu için, diğer nesneleri 4. satıra doğru hareket ettirebilir miyiz? Örn. O, A, M, D, T, P gibi nesneleri sırayla 4. satıra taşıyabiliriz. Öncelikle O'nun mevcut konumunu kontrol edelim.} \\
            }
        }
    }
    & & \\ \\

    \theutterance \stepcounter{utterance}  
    & & \multicolumn{4}{p{0.6\linewidth}}{
        \cellcolor[rgb]{0.9,0.9,0.9}{
            \makecell[{{p{\linewidth}}}]{
                \texttt{\tiny{[P1$\langle$GM]}}
                \texttt{Mesajın diğer oyuncuya iletildi.} \\
\\ 
\\ 
\texttt{Şu anda en fazla 28 turdan 10. turu oynuyorsunuz.} \\
\texttt{Toplamda 16 cezanın 19 kadarını biriktirdiniz.Diğer oyuncu şu mesajı gönderdi:} \\
\\ 
\\ 
\texttt{"Teşekkürler! Stratejimizi gözden geçirelim. S’nin şu anda (4,2) olduğu için, diğer nesneleri 4. satıra doğru hareket ettirebilir miyiz? Örn. O, A, M, D, T, P gibi nesneleri sırayla 4. satıra taşıyabiliriz. Öncelikle O'nun mevcut konumunu kontrol edelim."} \\
\\ 
\\ 
\texttt{Bir sonraki komutun nedir?} \\
            }
        }
    }
    & \\ \\

    \theutterance \stepcounter{utterance}  
    & \multicolumn{4}{p{0.6\linewidth}}{
        \cellcolor[rgb]{0.9,0.9,0.9}{
            \makecell[{{p{\linewidth}}}]{
                \texttt{\tiny{[P1$\rangle$GM]}}
                \texttt{SÖYLE: O'nun mevcut konumunu kontrol edelim.} \\
            }
        }
    }
    & & \\ \\

    \theutterance \stepcounter{utterance}  
    & & & \multicolumn{2}{p{0.3\linewidth}}{
        \cellcolor[rgb]{0.9,0.9,0.9}{
            \makecell[{{p{\linewidth}}}]{
                \texttt{\tiny{[GM$|$GM]}}
                \texttt{SÖYLE: O'nun mevcut konumunu kontrol edelim.} \\
            }
        }
    }
    & & \\ \\

    \theutterance \stepcounter{utterance}  
    & & \multicolumn{4}{p{0.6\linewidth}}{
        \cellcolor[rgb]{0.9,0.9,0.9}{
            \makecell[{{p{\linewidth}}}]{
                \texttt{\tiny{[GM$\rangle$P2]}}
                \texttt{Mesajın diğer oyuncuya iletildi.} \\
\\ 
\\ 
\texttt{Şu anda en fazla 28 turdan 11. turu oynuyorsunuz.} \\
\texttt{Toplamda 16 cezanın 19 kadarını biriktirdiniz.Diğer oyuncu şu mesajı gönderdi:} \\
\\ 
\\ 
\texttt{"O'nun mevcut konumunu kontrol edelim."} \\
\\ 
\\ 
\texttt{Bir sonraki komutun nedir?} \\
            }
        }
    }
    & \\ \\

    \theutterance \stepcounter{utterance}  
    & & & \multicolumn{4}{p{0.6\linewidth}}{
        \cellcolor[rgb]{0.9,0.9,0.9}{
            \makecell[{{p{\linewidth}}}]{
                \texttt{\tiny{[GM$\langle$P2]}}
                \texttt{SÖYLE: O'nun şu anki konumunu öğrenelim. Oyuncunun belirttiği gibi kontrollere başlayalım.} \\
            }
        }
    }
     \\ \\

    \theutterance \stepcounter{utterance}  
    & & & \multicolumn{2}{p{0.3\linewidth}}{
        \cellcolor[rgb]{0.9,0.9,0.9}{
            \makecell[{{p{\linewidth}}}]{
                \texttt{\tiny{[GM$|$GM]}}
                \texttt{SÖYLE: O'nun şu anki konumunu öğrenelim. Oyuncunun belirttiği gibi kontrollere başlayalım.} \\
            }
        }
    }
    & & \\ \\

    \theutterance \stepcounter{utterance}  
    & & \multicolumn{4}{p{0.6\linewidth}}{
        \cellcolor[rgb]{0.9,0.9,0.9}{
            \makecell[{{p{\linewidth}}}]{
                \texttt{\tiny{[P1$\langle$GM]}}
                \texttt{Mesajın diğer oyuncuya iletildi.} \\
\\ 
\\ 
\texttt{Şu anda en fazla 28 turdan 11. turu oynuyorsunuz.} \\
\texttt{Toplamda 16 cezanın 19 kadarını biriktirdiniz.Diğer oyuncu şu mesajı gönderdi:} \\
\\ 
\\ 
\texttt{"O'nun şu anki konumunu öğrenelim. Oyuncunun belirttiği gibi kontrollere başlayalım."} \\
\\ 
\\ 
\texttt{Bir sonraki komutun nedir?} \\
            }
        }
    }
    & \\ \\

    \theutterance \stepcounter{utterance}  
    & \multicolumn{4}{p{0.6\linewidth}}{
        \cellcolor[rgb]{0.9,0.9,0.9}{
            \makecell[{{p{\linewidth}}}]{
                \texttt{\tiny{[P1$\rangle$GM]}}
                \texttt{HAMLE: O, (5, 5)} \\
            }
        }
    }
    & & \\ \\

    \theutterance \stepcounter{utterance}  
    & & & \multicolumn{2}{p{0.3\linewidth}}{
        \cellcolor[rgb]{0.9,0.9,0.9}{
            \makecell[{{p{\linewidth}}}]{
                \texttt{\tiny{[GM$|$GM]}}
                \texttt{HAMLE: O, (5, 5)} \\
            }
        }
    }
    & & \\ \\

    \theutterance \stepcounter{utterance}  
    & & & \multicolumn{2}{p{0.3\linewidth}}{
        \cellcolor[rgb]{0.9,0.9,0.9}{
            \makecell[{{p{\linewidth}}}]{
                \texttt{\tiny{[GM$|$GM]}}
                \texttt{O moved to (5, 5)} \\
            }
        }
    }
    & & \\ \\

    \theutterance \stepcounter{utterance}  
    & & \multicolumn{4}{p{0.6\linewidth}}{
        \cellcolor[rgb]{0.9,0.9,0.9}{
            \makecell[{{p{\linewidth}}}]{
                \texttt{\tiny{[GM$\rangle$P2]}}
                \texttt{Mesajın diğer oyuncuya iletildi.} \\
\\ 
\\ 
\texttt{Şu anda en fazla 28 turdan 12. turu oynuyorsunuz.} \\
\texttt{Toplamda 16 cezanın 19 kadarını biriktirdiniz.Diğer oyuncu kendi ızgarasında bir nesneyi hareket ettirdi.} \\
\\ 
\\ 
\texttt{Bir sonraki komutun nedir?} \\
            }
        }
    }
    & \\ \\

    \theutterance \stepcounter{utterance}  
    & & & \multicolumn{4}{p{0.6\linewidth}}{
        \cellcolor[rgb]{0.9,0.9,0.9}{
            \makecell[{{p{\linewidth}}}]{
                \texttt{\tiny{[GM$\langle$P2]}}
                \texttt{HAMLE: O, (4, 3)} \\
            }
        }
    }
     \\ \\

    \theutterance \stepcounter{utterance}  
    & & & \multicolumn{2}{p{0.3\linewidth}}{
        \cellcolor[rgb]{0.9,0.9,0.9}{
            \makecell[{{p{\linewidth}}}]{
                \texttt{\tiny{[GM$|$GM]}}
                \texttt{HAMLE: O, (4, 3)} \\
            }
        }
    }
    & & \\ \\

    \theutterance \stepcounter{utterance}  
    & & & \multicolumn{2}{p{0.3\linewidth}}{
        \cellcolor[rgb]{0.9,0.9,0.9}{
            \makecell[{{p{\linewidth}}}]{
                \texttt{\tiny{[GM$|$GM]}}
                \texttt{Ceza: (4,3) boş değil; "M" içeriyor.} \\
\\ 
\texttt{Toplamda 16 cezanın 20 kadarını biriktirdiniz.Lütfen tekrar deneyin!} \\
            }
        }
    }
    & & \\ \\

    \theutterance \stepcounter{utterance}  
    & & \multicolumn{4}{p{0.6\linewidth}}{
        \cellcolor[rgb]{0.9,0.9,0.9}{
            \makecell[{{p{\linewidth}}}]{
                \texttt{\tiny{[GM$\rangle$P2]}}
                \texttt{Ceza: (4,3) boş değil; "M" içeriyor.} \\
\\ 
\texttt{Toplamda 16 cezanın 20 kadarını biriktirdiniz.Lütfen tekrar deneyin!} \\
            }
        }
    }
    & \\ \\

    \theutterance \stepcounter{utterance}  
    & & & \multicolumn{4}{p{0.6\linewidth}}{
        \cellcolor[rgb]{0.9,0.9,0.9}{
            \makecell[{{p{\linewidth}}}]{
                \texttt{\tiny{[GM$\langle$P2]}}
                \texttt{SÖYLE: O'nun mevcut konumu (2, 2). Hedef X için 4. satırda boş hücre olup olmadığını kontrol edelim. Örn. (4, 1), (4, 2), (4, 5), (4, 6) veya (4, 7) gibi hücrelerden hangisinin boş olduğunu kontrol edelim.} \\
            }
        }
    }
     \\ \\

    \theutterance \stepcounter{utterance}  
    & & & \multicolumn{2}{p{0.3\linewidth}}{
        \cellcolor[rgb]{0.9,0.9,0.9}{
            \makecell[{{p{\linewidth}}}]{
                \texttt{\tiny{[GM$|$GM]}}
                \texttt{SÖYLE: O'nun mevcut konumu (2, 2). Hedef X için 4. satırda boş hücre olup olmadığını kontrol edelim. Örn. (4, 1), (4, 2), (4, 5), (4, 6) veya (4, 7) gibi hücrelerden hangisinin boş olduğunu kontrol edelim.} \\
            }
        }
    }
    & & \\ \\

    \theutterance \stepcounter{utterance}  
    & & \multicolumn{4}{p{0.6\linewidth}}{
        \cellcolor[rgb]{0.9,0.9,0.9}{
            \makecell[{{p{\linewidth}}}]{
                \texttt{\tiny{[P1$\langle$GM]}}
                \texttt{O’ başarıyla (5,5) koordinatlarına taşındı. Güncellenmiş ızgaran şu şekilde:} \\
\texttt{```} \\
\texttt{ 1234567} \\
\texttt{╔═══════╗} \\
\texttt{║◌◌◌◌◌T◌║ 1} \\
\texttt{║P◌◌S◌◌◌║ 2} \\
\texttt{║◌◌AD◌┌─╢ 3} \\
\texttt{║◌◌◌┌─┤◌║ 4} \\
\texttt{║◌◌◌│O├─╢ 5} \\
\texttt{╟───┼─┘◌║ 6} \\
\texttt{║◌M◌│◌◌◌║ 7} \\
\texttt{╚═══╧═══╝} \\
\texttt{```} \\
\texttt{Şu anda en fazla 28 turdan 12. turu oynuyorsunuz.} \\
\texttt{Toplamda 16 cezanın 20 kadarını biriktirdiniz.Diğer oyuncu şu mesajı gönderdi:} \\
\\ 
\\ 
\texttt{"O'nun mevcut konumu (2, 2). Hedef X için 4. satırda boş hücre olup olmadığını kontrol edelim. Örn. (4, 1), (4, 2), (4, 5), (4, 6) veya (4, 7) gibi hücrelerden hangisinin boş olduğunu kontrol edelim."} \\
\\ 
\\ 
\texttt{Bir sonraki komutun nedir?} \\
            }
        }
    }
    & \\ \\

    \theutterance \stepcounter{utterance}  
    & \multicolumn{4}{p{0.6\linewidth}}{
        \cellcolor[rgb]{0.9,0.9,0.9}{
            \makecell[{{p{\linewidth}}}]{
                \texttt{\tiny{[P1$\rangle$GM]}}
                \texttt{HAMLE: O, (4, 2)} \\
            }
        }
    }
    & & \\ \\

    \theutterance \stepcounter{utterance}  
    & & & \multicolumn{2}{p{0.3\linewidth}}{
        \cellcolor[rgb]{0.9,0.9,0.9}{
            \makecell[{{p{\linewidth}}}]{
                \texttt{\tiny{[GM$|$GM]}}
                \texttt{HAMLE: O, (4, 2)} \\
            }
        }
    }
    & & \\ \\

    \theutterance \stepcounter{utterance}  
    & & & \multicolumn{2}{p{0.3\linewidth}}{
        \cellcolor[rgb]{0.9,0.9,0.9}{
            \makecell[{{p{\linewidth}}}]{
                \texttt{\tiny{[GM$|$GM]}}
                \texttt{Ceza: (4,2) boş değil; "S" içeriyor.} \\
\\ 
\texttt{Toplamda 16 cezanın 21 kadarını biriktirdiniz.Lütfen tekrar deneyin!} \\
            }
        }
    }
    & & \\ \\

    \theutterance \stepcounter{utterance}  
    & & \multicolumn{4}{p{0.6\linewidth}}{
        \cellcolor[rgb]{0.9,0.9,0.9}{
            \makecell[{{p{\linewidth}}}]{
                \texttt{\tiny{[P1$\langle$GM]}}
                \texttt{Ceza: (4,2) boş değil; "S" içeriyor.} \\
\\ 
\texttt{Toplamda 16 cezanın 21 kadarını biriktirdiniz.Lütfen tekrar deneyin!} \\
            }
        }
    }
    & \\ \\

    \theutterance \stepcounter{utterance}  
    & \multicolumn{4}{p{0.6\linewidth}}{
        \cellcolor[rgb]{0.9,0.9,0.9}{
            \makecell[{{p{\linewidth}}}]{
                \texttt{\tiny{[P1$\rangle$GM]}}
                \texttt{HAMLE: O, (4, 5)} \\
            }
        }
    }
    & & \\ \\

    \theutterance \stepcounter{utterance}  
    & & & \multicolumn{2}{p{0.3\linewidth}}{
        \cellcolor[rgb]{0.9,0.9,0.9}{
            \makecell[{{p{\linewidth}}}]{
                \texttt{\tiny{[GM$|$GM]}}
                \texttt{HAMLE: O, (4, 5)} \\
            }
        }
    }
    & & \\ \\

    \theutterance \stepcounter{utterance}  
    & & & \multicolumn{2}{p{0.3\linewidth}}{
        \cellcolor[rgb]{0.9,0.9,0.9}{
            \makecell[{{p{\linewidth}}}]{
                \texttt{\tiny{[GM$|$GM]}}
                \texttt{Ceza: (4,5) boş değil; "│" içeriyor.} \\
\\ 
\texttt{Toplamda 16 cezanın 22 kadarını biriktirdiniz.Lütfen tekrar deneyin!} \\
            }
        }
    }
    & & \\ \\

    \theutterance \stepcounter{utterance}  
    & & \multicolumn{4}{p{0.6\linewidth}}{
        \cellcolor[rgb]{0.9,0.9,0.9}{
            \makecell[{{p{\linewidth}}}]{
                \texttt{\tiny{[P1$\langle$GM]}}
                \texttt{Ceza: (4,5) boş değil; "│" içeriyor.} \\
\\ 
\texttt{Toplamda 16 cezanın 22 kadarını biriktirdiniz.Lütfen tekrar deneyin!} \\
            }
        }
    }
    & \\ \\

    \theutterance \stepcounter{utterance}  
    & \multicolumn{4}{p{0.6\linewidth}}{
        \cellcolor[rgb]{0.9,0.9,0.9}{
            \makecell[{{p{\linewidth}}}]{
                \texttt{\tiny{[P1$\rangle$GM]}}
                \texttt{HAMLE: O, (4, 6)} \\
            }
        }
    }
    & & \\ \\

    \theutterance \stepcounter{utterance}  
    & & & \multicolumn{2}{p{0.3\linewidth}}{
        \cellcolor[rgb]{0.9,0.9,0.9}{
            \makecell[{{p{\linewidth}}}]{
                \texttt{\tiny{[GM$|$GM]}}
                \texttt{HAMLE: O, (4, 6)} \\
            }
        }
    }
    & & \\ \\

    \theutterance \stepcounter{utterance}  
    & & & \multicolumn{2}{p{0.3\linewidth}}{
        \cellcolor[rgb]{0.9,0.9,0.9}{
            \makecell[{{p{\linewidth}}}]{
                \texttt{\tiny{[GM$|$GM]}}
                \texttt{Ceza: (4,6) boş değil; "┼" içeriyor.} \\
\\ 
\texttt{Toplamda 16 cezanın 23 kadarını biriktirdiniz.Lütfen tekrar deneyin!} \\
            }
        }
    }
    & & \\ \\

    \theutterance \stepcounter{utterance}  
    & & \multicolumn{4}{p{0.6\linewidth}}{
        \cellcolor[rgb]{0.9,0.9,0.9}{
            \makecell[{{p{\linewidth}}}]{
                \texttt{\tiny{[P1$\langle$GM]}}
                \texttt{Ceza: (4,6) boş değil; "┼" içeriyor.} \\
\\ 
\texttt{Toplamda 16 cezanın 23 kadarını biriktirdiniz.Lütfen tekrar deneyin!} \\
            }
        }
    }
    & \\ \\

    \theutterance \stepcounter{utterance}  
    & \multicolumn{4}{p{0.6\linewidth}}{
        \cellcolor[rgb]{0.9,0.9,0.9}{
            \makecell[{{p{\linewidth}}}]{
                \texttt{\tiny{[P1$\rangle$GM]}}
                \texttt{HAMLE: O, (4, 7)} \\
            }
        }
    }
    & & \\ \\

    \theutterance \stepcounter{utterance}  
    & & & \multicolumn{2}{p{0.3\linewidth}}{
        \cellcolor[rgb]{0.9,0.9,0.9}{
            \makecell[{{p{\linewidth}}}]{
                \texttt{\tiny{[GM$|$GM]}}
                \texttt{HAMLE: O, (4, 7)} \\
            }
        }
    }
    & & \\ \\

    \theutterance \stepcounter{utterance}  
    & & & \multicolumn{2}{p{0.3\linewidth}}{
        \cellcolor[rgb]{0.9,0.9,0.9}{
            \makecell[{{p{\linewidth}}}]{
                \texttt{\tiny{[GM$|$GM]}}
                \texttt{Ceza: (4,7) boş değil; "│" içeriyor.} \\
\\ 
\texttt{Toplamda 16 cezanın 24 kadarını biriktirdiniz.Lütfen tekrar deneyin!} \\
            }
        }
    }
    & & \\ \\

    \theutterance \stepcounter{utterance}  
    & & \multicolumn{4}{p{0.6\linewidth}}{
        \cellcolor[rgb]{0.9,0.9,0.9}{
            \makecell[{{p{\linewidth}}}]{
                \texttt{\tiny{[P1$\langle$GM]}}
                \texttt{Ceza: (4,7) boş değil; "│" içeriyor.} \\
\\ 
\texttt{Toplamda 16 cezanın 24 kadarını biriktirdiniz.Lütfen tekrar deneyin!} \\
            }
        }
    }
    & \\ \\

    \theutterance \stepcounter{utterance}  
    & \multicolumn{4}{p{0.6\linewidth}}{
        \cellcolor[rgb]{0.9,0.9,0.9}{
            \makecell[{{p{\linewidth}}}]{
                \texttt{\tiny{[P1$\rangle$GM]}}
                \texttt{HAMLE: O, (3, 4)} \\
            }
        }
    }
    & & \\ \\

    \theutterance \stepcounter{utterance}  
    & & & \multicolumn{2}{p{0.3\linewidth}}{
        \cellcolor[rgb]{0.9,0.9,0.9}{
            \makecell[{{p{\linewidth}}}]{
                \texttt{\tiny{[GM$|$GM]}}
                \texttt{HAMLE: O, (3, 4)} \\
            }
        }
    }
    & & \\ \\

    \theutterance \stepcounter{utterance}  
    & & & \multicolumn{2}{p{0.3\linewidth}}{
        \cellcolor[rgb]{0.9,0.9,0.9}{
            \makecell[{{p{\linewidth}}}]{
                \texttt{\tiny{[GM$|$GM]}}
                \texttt{O moved to (3, 4)} \\
            }
        }
    }
    & & \\ \\

    \theutterance \stepcounter{utterance}  
    & & \multicolumn{4}{p{0.6\linewidth}}{
        \cellcolor[rgb]{0.9,0.9,0.9}{
            \makecell[{{p{\linewidth}}}]{
                \texttt{\tiny{[GM$\rangle$P2]}}
                \texttt{Mesajın diğer oyuncuya iletildi.} \\
\\ 
\\ 
\texttt{Şu anda en fazla 28 turdan 13. turu oynuyorsunuz.} \\
\texttt{Toplamda 16 cezanın 24 kadarını biriktirdiniz.Diğer oyuncu kendi ızgarasında bir nesneyi hareket ettirdi.} \\
\\ 
\\ 
\texttt{Bir sonraki komutun nedir?} \\
            }
        }
    }
    & \\ \\

    \theutterance \stepcounter{utterance}  
    & & & \multicolumn{4}{p{0.6\linewidth}}{
        \cellcolor[rgb]{0.9,0.9,0.9}{
            \makecell[{{p{\linewidth}}}]{
                \texttt{\tiny{[GM$\langle$P2]}}
                \texttt{HAMLE: O, (4, 1)} \\
            }
        }
    }
     \\ \\

    \theutterance \stepcounter{utterance}  
    & & & \multicolumn{2}{p{0.3\linewidth}}{
        \cellcolor[rgb]{0.9,0.9,0.9}{
            \makecell[{{p{\linewidth}}}]{
                \texttt{\tiny{[GM$|$GM]}}
                \texttt{HAMLE: O, (4, 1)} \\
            }
        }
    }
    & & \\ \\

    \theutterance \stepcounter{utterance}  
    & & & \multicolumn{2}{p{0.3\linewidth}}{
        \cellcolor[rgb]{0.9,0.9,0.9}{
            \makecell[{{p{\linewidth}}}]{
                \texttt{\tiny{[GM$|$GM]}}
                \texttt{O moved to (4, 1)} \\
            }
        }
    }
    & & \\ \\

    \theutterance \stepcounter{utterance}  
    & & \multicolumn{4}{p{0.6\linewidth}}{
        \cellcolor[rgb]{0.9,0.9,0.9}{
            \makecell[{{p{\linewidth}}}]{
                \texttt{\tiny{[P1$\langle$GM]}}
                \texttt{O’ başarıyla (3,4) koordinatlarına taşındı. Güncellenmiş ızgaran şu şekilde:} \\
\texttt{```} \\
\texttt{ 1234567} \\
\texttt{╔═══════╗} \\
\texttt{║◌◌◌◌◌T◌║ 1} \\
\texttt{║P◌◌S◌◌◌║ 2} \\
\texttt{║◌◌AD◌┌─╢ 3} \\
\texttt{║◌◌O┌─┤◌║ 4} \\
\texttt{║◌◌◌│◌├─╢ 5} \\
\texttt{╟───┼─┘◌║ 6} \\
\texttt{║◌M◌│◌◌◌║ 7} \\
\texttt{╚═══╧═══╝} \\
\texttt{```} \\
\texttt{Şu anda en fazla 28 turdan 13. turu oynuyorsunuz.} \\
\texttt{Toplamda 16 cezanın 24 kadarını biriktirdiniz.Diğer oyuncu kendi ızgarasında bir nesneyi hareket ettirdi.} \\
\\ 
\\ 
\texttt{Bir sonraki komutun nedir?} \\
            }
        }
    }
    & \\ \\

    \theutterance \stepcounter{utterance}  
    & \multicolumn{4}{p{0.6\linewidth}}{
        \cellcolor[rgb]{0.9,0.9,0.9}{
            \makecell[{{p{\linewidth}}}]{
                \texttt{\tiny{[P1$\rangle$GM]}}
                \texttt{SÖYLE: S (4,2), O (3,4). İki nesneyi de 4. satıra yönlendirmek için (4, 2) ve (3, 4) dışında 4. satırda boş hücre olup olmadığını kontrol edelim.} \\
            }
        }
    }
    & & \\ \\

    \theutterance \stepcounter{utterance}  
    & & & \multicolumn{2}{p{0.3\linewidth}}{
        \cellcolor[rgb]{0.9,0.9,0.9}{
            \makecell[{{p{\linewidth}}}]{
                \texttt{\tiny{[GM$|$GM]}}
                \texttt{SÖYLE: S (4,2), O (3,4). İki nesneyi de 4. satıra yönlendirmek için (4, 2) ve (3, 4) dışında 4. satırda boş hücre olup olmadığını kontrol edelim.} \\
            }
        }
    }
    & & \\ \\

    \theutterance \stepcounter{utterance}  
    & & \multicolumn{4}{p{0.6\linewidth}}{
        \cellcolor[rgb]{0.9,0.9,0.9}{
            \makecell[{{p{\linewidth}}}]{
                \texttt{\tiny{[GM$\rangle$P2]}}
                \texttt{O’ başarıyla (4,1) koordinatlarına taşındı. Güncellenmiş ızgaran şu şekilde:} \\
\texttt{```} \\
\texttt{ 1234567} \\
\texttt{╔═══════╗} \\
\texttt{║◌◌◌O◌◌◌║ 1} \\
\texttt{║◌◌◌S◌◌◌║ 2} \\
\texttt{║◌◌◌M◌┌─╢ 3} \\
\texttt{║◌◌A┌─┤◌║ 4} \\
\texttt{║◌◌◌│◌├─╢ 5} \\
\texttt{╟───┼─┘◌║ 6} \\
\texttt{║D◌T│P◌◌║ 7} \\
\texttt{╚═══╧═══╝} \\
\texttt{```} \\
\texttt{Şu anda en fazla 28 turdan 14. turu oynuyorsunuz.} \\
\texttt{Toplamda 16 cezanın 24 kadarını biriktirdiniz.Diğer oyuncu şu mesajı gönderdi:} \\
\\ 
\\ 
\texttt{"S (4,2), O (3,4). İki nesneyi de 4. satıra yönlendirmek için (4, 2) ve (3, 4) dışında 4. satırda boş hücre olup olmadığını kontrol edelim."} \\
\\ 
\\ 
\texttt{Bir sonraki komutun nedir?} \\
            }
        }
    }
    & \\ \\

    \theutterance \stepcounter{utterance}  
    & & & \multicolumn{4}{p{0.6\linewidth}}{
        \cellcolor[rgb]{0.9,0.9,0.9}{
            \makecell[{{p{\linewidth}}}]{
                \texttt{\tiny{[GM$\langle$P2]}}
                \texttt{SÖYLE: 4. satırda (4, 2), (4, 3), (4, 4), (4, 5), (4, 6), (4, 7) dışında hangi hücrelerin boş olduğunu kontrol edelim. S ve O dışındaki tüm nesneleri hizalamanın adımlarını planlayalım.} \\
            }
        }
    }
     \\ \\

    \theutterance \stepcounter{utterance}  
    & & & \multicolumn{2}{p{0.3\linewidth}}{
        \cellcolor[rgb]{0.9,0.9,0.9}{
            \makecell[{{p{\linewidth}}}]{
                \texttt{\tiny{[GM$|$GM]}}
                \texttt{SÖYLE: 4. satırda (4, 2), (4, 3), (4, 4), (4, 5), (4, 6), (4, 7) dışında hangi hücrelerin boş olduğunu kontrol edelim. S ve O dışındaki tüm nesneleri hizalamanın adımlarını planlayalım.} \\
            }
        }
    }
    & & \\ \\

    \theutterance \stepcounter{utterance}  
    & & \multicolumn{4}{p{0.6\linewidth}}{
        \cellcolor[rgb]{0.9,0.9,0.9}{
            \makecell[{{p{\linewidth}}}]{
                \texttt{\tiny{[P1$\langle$GM]}}
                \texttt{Mesajın diğer oyuncuya iletildi.} \\
\\ 
\\ 
\texttt{Şu anda en fazla 28 turdan 14. turu oynuyorsunuz.} \\
\texttt{Toplamda 16 cezanın 24 kadarını biriktirdiniz.Diğer oyuncu şu mesajı gönderdi:} \\
\\ 
\\ 
\texttt{"4. satırda (4, 2), (4, 3), (4, 4), (4, 5), (4, 6), (4, 7) dışında hangi hücrelerin boş olduğunu kontrol edelim. S ve O dışındaki tüm nesneleri hizalamanın adımlarını planlayalım."} \\
\\ 
\\ 
\texttt{Bir sonraki komutun nedir?} \\
            }
        }
    }
    & \\ \\

    \theutterance \stepcounter{utterance}  
    & \multicolumn{4}{p{0.6\linewidth}}{
        \cellcolor[rgb]{0.9,0.9,0.9}{
            \makecell[{{p{\linewidth}}}]{
                \texttt{\tiny{[P1$\rangle$GM]}}
                \texttt{SÖYLE: 4. satırda (4, 2) artık S ile dolu. (4, 3) ve (4, 4) durumunu kontrol edelim.} \\
            }
        }
    }
    & & \\ \\

    \theutterance \stepcounter{utterance}  
    & & & \multicolumn{2}{p{0.3\linewidth}}{
        \cellcolor[rgb]{0.9,0.9,0.9}{
            \makecell[{{p{\linewidth}}}]{
                \texttt{\tiny{[GM$|$GM]}}
                \texttt{SÖYLE: 4. satırda (4, 2) artık S ile dolu. (4, 3) ve (4, 4) durumunu kontrol edelim.} \\
            }
        }
    }
    & & \\ \\

    \theutterance \stepcounter{utterance}  
    & & \multicolumn{4}{p{0.6\linewidth}}{
        \cellcolor[rgb]{0.9,0.9,0.9}{
            \makecell[{{p{\linewidth}}}]{
                \texttt{\tiny{[GM$\rangle$P2]}}
                \texttt{Mesajın diğer oyuncuya iletildi.} \\
\\ 
\\ 
\texttt{Şu anda en fazla 28 turdan 15. turu oynuyorsunuz.} \\
\texttt{Toplamda 16 cezanın 24 kadarını biriktirdiniz.Diğer oyuncu şu mesajı gönderdi:} \\
\\ 
\\ 
\texttt{"4. satırda (4, 2) artık S ile dolu. (4, 3) ve (4, 4) durumunu kontrol edelim."} \\
\\ 
\\ 
\texttt{Bir sonraki komutun nedir?} \\
            }
        }
    }
    & \\ \\

    \theutterance \stepcounter{utterance}  
    & & & \multicolumn{4}{p{0.6\linewidth}}{
        \cellcolor[rgb]{0.9,0.9,0.9}{
            \makecell[{{p{\linewidth}}}]{
                \texttt{\tiny{[GM$\langle$P2]}}
                \texttt{HAMLE: M, (4, 3)} \\
            }
        }
    }
     \\ \\

    \theutterance \stepcounter{utterance}  
    & & & \multicolumn{2}{p{0.3\linewidth}}{
        \cellcolor[rgb]{0.9,0.9,0.9}{
            \makecell[{{p{\linewidth}}}]{
                \texttt{\tiny{[GM$|$GM]}}
                \texttt{HAMLE: M, (4, 3)} \\
            }
        }
    }
    & & \\ \\

    \theutterance \stepcounter{utterance}  
    & & & \multicolumn{2}{p{0.3\linewidth}}{
        \cellcolor[rgb]{0.9,0.9,0.9}{
            \makecell[{{p{\linewidth}}}]{
                \texttt{\tiny{[GM$|$GM]}}
                \texttt{Ceza: (4,3) boş değil; "M" içeriyor.} \\
\\ 
\texttt{Toplamda 16 cezanın 25 kadarını biriktirdiniz.Lütfen tekrar deneyin!} \\
            }
        }
    }
    & & \\ \\

    \theutterance \stepcounter{utterance}  
    & & \multicolumn{4}{p{0.6\linewidth}}{
        \cellcolor[rgb]{0.9,0.9,0.9}{
            \makecell[{{p{\linewidth}}}]{
                \texttt{\tiny{[GM$\rangle$P2]}}
                \texttt{Ceza: (4,3) boş değil; "M" içeriyor.} \\
\\ 
\texttt{Toplamda 16 cezanın 25 kadarını biriktirdiniz.Lütfen tekrar deneyin!} \\
            }
        }
    }
    & \\ \\

    \theutterance \stepcounter{utterance}  
    & & & \multicolumn{4}{p{0.6\linewidth}}{
        \cellcolor[rgb]{0.9,0.9,0.9}{
            \makecell[{{p{\linewidth}}}]{
                \texttt{\tiny{[GM$\langle$P2]}}
                \texttt{SÖYLE: 4. satırda (4, 3) ve (4, 4) konumlarını kontrol edelim. (4, 3) ya da (4, 4) dışındaki boş hücre olup olmadığını kontrol edelim, örn. (4, 1), (4, 5), (4, 6), (4, 7).} \\
            }
        }
    }
     \\ \\

    \theutterance \stepcounter{utterance}  
    & & & \multicolumn{2}{p{0.3\linewidth}}{
        \cellcolor[rgb]{0.9,0.9,0.9}{
            \makecell[{{p{\linewidth}}}]{
                \texttt{\tiny{[GM$|$GM]}}
                \texttt{SÖYLE: 4. satırda (4, 3) ve (4, 4) konumlarını kontrol edelim. (4, 3) ya da (4, 4) dışındaki boş hücre olup olmadığını kontrol edelim, örn. (4, 1), (4, 5), (4, 6), (4, 7).} \\
            }
        }
    }
    & & \\ \\

    \theutterance \stepcounter{utterance}  
    & & \multicolumn{4}{p{0.6\linewidth}}{
        \cellcolor[rgb]{0.9,0.9,0.9}{
            \makecell[{{p{\linewidth}}}]{
                \texttt{\tiny{[P1$\langle$GM]}}
                \texttt{Mesajın diğer oyuncuya iletildi.} \\
\\ 
\\ 
\texttt{Şu anda en fazla 28 turdan 15. turu oynuyorsunuz.} \\
\texttt{Toplamda 16 cezanın 25 kadarını biriktirdiniz.Diğer oyuncu şu mesajı gönderdi:} \\
\\ 
\\ 
\texttt{"4. satırda (4, 3) ve (4, 4) konumlarını kontrol edelim. (4, 3) ya da (4, 4) dışındaki boş hücre olup olmadığını kontrol edelim, örn. (4, 1), (4, 5), (4, 6), (4, 7)."} \\
\\ 
\\ 
\texttt{Bir sonraki komutun nedir?} \\
            }
        }
    }
    & \\ \\

    \theutterance \stepcounter{utterance}  
    & \multicolumn{4}{p{0.6\linewidth}}{
        \cellcolor[rgb]{0.9,0.9,0.9}{
            \makecell[{{p{\linewidth}}}]{
                \texttt{\tiny{[P1$\rangle$GM]}}
                \texttt{SÖYLE: 4. satırda (4, 3) "D" içeriyor, (4, 4) "┌" içeriyor; dolayısıyla bu hücreler dolu. Öncelikle (4, 1) ya da (4, 7)'yi kontrol edelim.} \\
            }
        }
    }
    & & \\ \\

    \theutterance \stepcounter{utterance}  
    & & & \multicolumn{2}{p{0.3\linewidth}}{
        \cellcolor[rgb]{0.9,0.9,0.9}{
            \makecell[{{p{\linewidth}}}]{
                \texttt{\tiny{[GM$|$GM]}}
                \texttt{SÖYLE: 4. satırda (4, 3) "D" içeriyor, (4, 4) "┌" içeriyor; dolayısıyla bu hücreler dolu. Öncelikle (4, 1) ya da (4, 7)'yi kontrol edelim.} \\
            }
        }
    }
    & & \\ \\

    \theutterance \stepcounter{utterance}  
    & & \multicolumn{4}{p{0.6\linewidth}}{
        \cellcolor[rgb]{0.9,0.9,0.9}{
            \makecell[{{p{\linewidth}}}]{
                \texttt{\tiny{[GM$\rangle$P2]}}
                \texttt{Mesajın diğer oyuncuya iletildi.} \\
\\ 
\\ 
\texttt{Şu anda en fazla 28 turdan 16. turu oynuyorsunuz.} \\
\texttt{Toplamda 16 cezanın 25 kadarını biriktirdiniz.Diğer oyuncu şu mesajı gönderdi:} \\
\\ 
\\ 
\texttt{"4. satırda (4, 3) "D" içeriyor, (4, 4) "┌" içeriyor; dolayısıyla bu hücreler dolu. Öncelikle (4, 1) ya da (4, 7)'yi kontrol edelim."} \\
\\ 
\\ 
\texttt{Bir sonraki komutun nedir?} \\
            }
        }
    }
    & \\ \\

    \theutterance \stepcounter{utterance}  
    & & & \multicolumn{4}{p{0.6\linewidth}}{
        \cellcolor[rgb]{0.9,0.9,0.9}{
            \makecell[{{p{\linewidth}}}]{
                \texttt{\tiny{[GM$\langle$P2]}}
                \texttt{HAMLE: M, (4, 1)} \\
            }
        }
    }
     \\ \\

    \theutterance \stepcounter{utterance}  
    & & & \multicolumn{2}{p{0.3\linewidth}}{
        \cellcolor[rgb]{0.9,0.9,0.9}{
            \makecell[{{p{\linewidth}}}]{
                \texttt{\tiny{[GM$|$GM]}}
                \texttt{HAMLE: M, (4, 1)} \\
            }
        }
    }
    & & \\ \\

    \theutterance \stepcounter{utterance}  
    & & & \multicolumn{2}{p{0.3\linewidth}}{
        \cellcolor[rgb]{0.9,0.9,0.9}{
            \makecell[{{p{\linewidth}}}]{
                \texttt{\tiny{[GM$|$GM]}}
                \texttt{Ceza: (4,1) boş değil; "O" içeriyor.} \\
\\ 
\texttt{Toplamda 16 cezanın 26 kadarını biriktirdiniz.Lütfen tekrar deneyin!} \\
            }
        }
    }
    & & \\ \\

    \theutterance \stepcounter{utterance}  
    & & \multicolumn{4}{p{0.6\linewidth}}{
        \cellcolor[rgb]{0.9,0.9,0.9}{
            \makecell[{{p{\linewidth}}}]{
                \texttt{\tiny{[GM$\rangle$P2]}}
                \texttt{Ceza: (4,1) boş değil; "O" içeriyor.} \\
\\ 
\texttt{Toplamda 16 cezanın 26 kadarını biriktirdiniz.Lütfen tekrar deneyin!} \\
            }
        }
    }
    & \\ \\

    \theutterance \stepcounter{utterance}  
    & & & \multicolumn{4}{p{0.6\linewidth}}{
        \cellcolor[rgb]{0.9,0.9,0.9}{
            \makecell[{{p{\linewidth}}}]{
                \texttt{\tiny{[GM$\langle$P2]}}
                \texttt{SÖYLE: 4. satırda (4, 1), (4, 2), (4, 3), (4, 4), (4, 5), (4, 6), (4, 7) dışında kalan boş hücre olup olmadığını kontrol edelim. İki nesne (S ve O) da 4. satıra yerleşti; şimdi diğer nesnelerin sırayla 4. satıra taşınması stratejisini paylaşabiliriz.} \\
            }
        }
    }
     \\ \\

    \theutterance \stepcounter{utterance}  
    & & & \multicolumn{2}{p{0.3\linewidth}}{
        \cellcolor[rgb]{0.9,0.9,0.9}{
            \makecell[{{p{\linewidth}}}]{
                \texttt{\tiny{[GM$|$GM]}}
                \texttt{SÖYLE: 4. satırda (4, 1), (4, 2), (4, 3), (4, 4), (4, 5), (4, 6), (4, 7) dışında kalan boş hücre olup olmadığını kontrol edelim. İki nesne (S ve O) da 4. satıra yerleşti; şimdi diğer nesnelerin sırayla 4. satıra taşınması stratejisini paylaşabiliriz.} \\
            }
        }
    }
    & & \\ \\

    \theutterance \stepcounter{utterance}  
    & & \multicolumn{4}{p{0.6\linewidth}}{
        \cellcolor[rgb]{0.9,0.9,0.9}{
            \makecell[{{p{\linewidth}}}]{
                \texttt{\tiny{[P1$\langle$GM]}}
                \texttt{Mesajın diğer oyuncuya iletildi.} \\
\\ 
\\ 
\texttt{Şu anda en fazla 28 turdan 16. turu oynuyorsunuz.} \\
\texttt{Toplamda 16 cezanın 26 kadarını biriktirdiniz.Diğer oyuncu şu mesajı gönderdi:} \\
\\ 
\\ 
\texttt{"4. satırda (4, 1), (4, 2), (4, 3), (4, 4), (4, 5), (4, 6), (4, 7) dışında kalan boş hücre olup olmadığını kontrol edelim. İki nesne (S ve O) da 4. satıra yerleşti; şimdi diğer nesnelerin sırayla 4. satıra taşınması stratejisini paylaşabiliriz."} \\
\\ 
\\ 
\texttt{Bir sonraki komutun nedir?} \\
            }
        }
    }
    & \\ \\

    \theutterance \stepcounter{utterance}  
    & \multicolumn{4}{p{0.6\linewidth}}{
        \cellcolor[rgb]{0.9,0.9,0.9}{
            \makecell[{{p{\linewidth}}}]{
                \texttt{\tiny{[P1$\rangle$GM]}}
                \texttt{HAMLE: D, (4, 3)} \\
            }
        }
    }
    & & \\ \\

    \theutterance \stepcounter{utterance}  
    & & & \multicolumn{2}{p{0.3\linewidth}}{
        \cellcolor[rgb]{0.9,0.9,0.9}{
            \makecell[{{p{\linewidth}}}]{
                \texttt{\tiny{[GM$|$GM]}}
                \texttt{HAMLE: D, (4, 3)} \\
            }
        }
    }
    & & \\ \\

    \theutterance \stepcounter{utterance}  
    & & & \multicolumn{2}{p{0.3\linewidth}}{
        \cellcolor[rgb]{0.9,0.9,0.9}{
            \makecell[{{p{\linewidth}}}]{
                \texttt{\tiny{[GM$|$GM]}}
                \texttt{Ceza: (4,3) boş değil; "D" içeriyor.} \\
\\ 
\texttt{Toplamda 16 cezanın 27 kadarını biriktirdiniz.Lütfen tekrar deneyin!} \\
            }
        }
    }
    & & \\ \\

    \theutterance \stepcounter{utterance}  
    & & \multicolumn{4}{p{0.6\linewidth}}{
        \cellcolor[rgb]{0.9,0.9,0.9}{
            \makecell[{{p{\linewidth}}}]{
                \texttt{\tiny{[P1$\langle$GM]}}
                \texttt{Ceza: (4,3) boş değil; "D" içeriyor.} \\
\\ 
\texttt{Toplamda 16 cezanın 27 kadarını biriktirdiniz.Lütfen tekrar deneyin!} \\
            }
        }
    }
    & \\ \\

    \theutterance \stepcounter{utterance}  
    & \multicolumn{4}{p{0.6\linewidth}}{
        \cellcolor[rgb]{0.9,0.9,0.9}{
            \makecell[{{p{\linewidth}}}]{
                \texttt{\tiny{[P1$\rangle$GM]}}
                \texttt{HAMLE: D, (4, 5)} \\
            }
        }
    }
    & & \\ \\

    \theutterance \stepcounter{utterance}  
    & & & \multicolumn{2}{p{0.3\linewidth}}{
        \cellcolor[rgb]{0.9,0.9,0.9}{
            \makecell[{{p{\linewidth}}}]{
                \texttt{\tiny{[GM$|$GM]}}
                \texttt{HAMLE: D, (4, 5)} \\
            }
        }
    }
    & & \\ \\

    \theutterance \stepcounter{utterance}  
    & & & \multicolumn{2}{p{0.3\linewidth}}{
        \cellcolor[rgb]{0.9,0.9,0.9}{
            \makecell[{{p{\linewidth}}}]{
                \texttt{\tiny{[GM$|$GM]}}
                \texttt{Ceza: (4,5) boş değil; "│" içeriyor.} \\
\\ 
\texttt{Toplamda 16 cezanın 28 kadarını biriktirdiniz.Lütfen tekrar deneyin!} \\
            }
        }
    }
    & & \\ \\

    \theutterance \stepcounter{utterance}  
    & & \multicolumn{4}{p{0.6\linewidth}}{
        \cellcolor[rgb]{0.9,0.9,0.9}{
            \makecell[{{p{\linewidth}}}]{
                \texttt{\tiny{[P1$\langle$GM]}}
                \texttt{Ceza: (4,5) boş değil; "│" içeriyor.} \\
\\ 
\texttt{Toplamda 16 cezanın 28 kadarını biriktirdiniz.Lütfen tekrar deneyin!} \\
            }
        }
    }
    & \\ \\

    \theutterance \stepcounter{utterance}  
    & \multicolumn{4}{p{0.6\linewidth}}{
        \cellcolor[rgb]{0.9,0.9,0.9}{
            \makecell[{{p{\linewidth}}}]{
                \texttt{\tiny{[P1$\rangle$GM]}}
                \texttt{HAMLE: D, (4, 6)} \\
            }
        }
    }
    & & \\ \\

    \theutterance \stepcounter{utterance}  
    & & & \multicolumn{2}{p{0.3\linewidth}}{
        \cellcolor[rgb]{0.9,0.9,0.9}{
            \makecell[{{p{\linewidth}}}]{
                \texttt{\tiny{[GM$|$GM]}}
                \texttt{HAMLE: D, (4, 6)} \\
            }
        }
    }
    & & \\ \\

    \theutterance \stepcounter{utterance}  
    & & & \multicolumn{2}{p{0.3\linewidth}}{
        \cellcolor[rgb]{0.9,0.9,0.9}{
            \makecell[{{p{\linewidth}}}]{
                \texttt{\tiny{[GM$|$GM]}}
                \texttt{Ceza: (4,6) boş değil; "┼" içeriyor.} \\
\\ 
\texttt{Toplamda 16 cezanın 29 kadarını biriktirdiniz.Lütfen tekrar deneyin!} \\
            }
        }
    }
    & & \\ \\

    \theutterance \stepcounter{utterance}  
    & & \multicolumn{4}{p{0.6\linewidth}}{
        \cellcolor[rgb]{0.9,0.9,0.9}{
            \makecell[{{p{\linewidth}}}]{
                \texttt{\tiny{[P1$\langle$GM]}}
                \texttt{Ceza: (4,6) boş değil; "┼" içeriyor.} \\
\\ 
\texttt{Toplamda 16 cezanın 29 kadarını biriktirdiniz.Lütfen tekrar deneyin!} \\
            }
        }
    }
    & \\ \\

    \theutterance \stepcounter{utterance}  
    & \multicolumn{4}{p{0.6\linewidth}}{
        \cellcolor[rgb]{0.9,0.9,0.9}{
            \makecell[{{p{\linewidth}}}]{
                \texttt{\tiny{[P1$\rangle$GM]}}
                \texttt{HAMLE: D, (4, 7)} \\
            }
        }
    }
    & & \\ \\

    \theutterance \stepcounter{utterance}  
    & & & \multicolumn{2}{p{0.3\linewidth}}{
        \cellcolor[rgb]{0.9,0.9,0.9}{
            \makecell[{{p{\linewidth}}}]{
                \texttt{\tiny{[GM$|$GM]}}
                \texttt{HAMLE: D, (4, 7)} \\
            }
        }
    }
    & & \\ \\

    \theutterance \stepcounter{utterance}  
    & & & \multicolumn{2}{p{0.3\linewidth}}{
        \cellcolor[rgb]{0.9,0.9,0.9}{
            \makecell[{{p{\linewidth}}}]{
                \texttt{\tiny{[GM$|$GM]}}
                \texttt{Ceza: (4,7) boş değil; "│" içeriyor.} \\
\\ 
\texttt{Toplamda 16 cezanın 30 kadarını biriktirdiniz.Lütfen tekrar deneyin!} \\
            }
        }
    }
    & & \\ \\

    \theutterance \stepcounter{utterance}  
    & & \multicolumn{4}{p{0.6\linewidth}}{
        \cellcolor[rgb]{0.9,0.9,0.9}{
            \makecell[{{p{\linewidth}}}]{
                \texttt{\tiny{[P1$\langle$GM]}}
                \texttt{Ceza: (4,7) boş değil; "│" içeriyor.} \\
\\ 
\texttt{Toplamda 16 cezanın 30 kadarını biriktirdiniz.Lütfen tekrar deneyin!} \\
            }
        }
    }
    & \\ \\

    \theutterance \stepcounter{utterance}  
    & \multicolumn{4}{p{0.6\linewidth}}{
        \cellcolor[rgb]{0.9,0.9,0.9}{
            \makecell[{{p{\linewidth}}}]{
                \texttt{\tiny{[P1$\rangle$GM]}}
                \texttt{HAMLE: D, (4, 2)} \\
            }
        }
    }
    & & \\ \\

    \theutterance \stepcounter{utterance}  
    & & & \multicolumn{2}{p{0.3\linewidth}}{
        \cellcolor[rgb]{0.9,0.9,0.9}{
            \makecell[{{p{\linewidth}}}]{
                \texttt{\tiny{[GM$|$GM]}}
                \texttt{HAMLE: D, (4, 2)} \\
            }
        }
    }
    & & \\ \\

    \theutterance \stepcounter{utterance}  
    & & & \multicolumn{2}{p{0.3\linewidth}}{
        \cellcolor[rgb]{0.9,0.9,0.9}{
            \makecell[{{p{\linewidth}}}]{
                \texttt{\tiny{[GM$|$GM]}}
                \texttt{Ceza: (4,2) boş değil; "S" içeriyor.} \\
\\ 
\texttt{Toplamda 16 cezanın 31 kadarını biriktirdiniz.Lütfen tekrar deneyin!} \\
            }
        }
    }
    & & \\ \\

    \theutterance \stepcounter{utterance}  
    & & \multicolumn{4}{p{0.6\linewidth}}{
        \cellcolor[rgb]{0.9,0.9,0.9}{
            \makecell[{{p{\linewidth}}}]{
                \texttt{\tiny{[P1$\langle$GM]}}
                \texttt{Ceza: (4,2) boş değil; "S" içeriyor.} \\
\\ 
\texttt{Toplamda 16 cezanın 31 kadarını biriktirdiniz.Lütfen tekrar deneyin!} \\
            }
        }
    }
    & \\ \\

    \theutterance \stepcounter{utterance}  
    & \multicolumn{4}{p{0.6\linewidth}}{
        \cellcolor[rgb]{0.9,0.9,0.9}{
            \makecell[{{p{\linewidth}}}]{
                \texttt{\tiny{[P1$\rangle$GM]}}
                \texttt{HAMLE: D, (4, 1)} \\
            }
        }
    }
    & & \\ \\

    \theutterance \stepcounter{utterance}  
    & & & \multicolumn{2}{p{0.3\linewidth}}{
        \cellcolor[rgb]{0.9,0.9,0.9}{
            \makecell[{{p{\linewidth}}}]{
                \texttt{\tiny{[GM$|$GM]}}
                \texttt{HAMLE: D, (4, 1)} \\
            }
        }
    }
    & & \\ \\

    \theutterance \stepcounter{utterance}  
    & & & \multicolumn{2}{p{0.3\linewidth}}{
        \cellcolor[rgb]{0.9,0.9,0.9}{
            \makecell[{{p{\linewidth}}}]{
                \texttt{\tiny{[GM$|$GM]}}
                \texttt{D moved to (4, 1)} \\
            }
        }
    }
    & & \\ \\

    \theutterance \stepcounter{utterance}  
    & & \multicolumn{4}{p{0.6\linewidth}}{
        \cellcolor[rgb]{0.9,0.9,0.9}{
            \makecell[{{p{\linewidth}}}]{
                \texttt{\tiny{[GM$\rangle$P2]}}
                \texttt{Mesajın diğer oyuncuya iletildi.} \\
\\ 
\\ 
\texttt{Şu anda en fazla 28 turdan 17. turu oynuyorsunuz.} \\
\texttt{Toplamda 16 cezanın 31 kadarını biriktirdiniz.Diğer oyuncu kendi ızgarasında bir nesneyi hareket ettirdi.} \\
\\ 
\\ 
\texttt{Bir sonraki komutun nedir?} \\
            }
        }
    }
    & \\ \\

    \theutterance \stepcounter{utterance}  
    & & & \multicolumn{4}{p{0.6\linewidth}}{
        \cellcolor[rgb]{0.9,0.9,0.9}{
            \makecell[{{p{\linewidth}}}]{
                \texttt{\tiny{[GM$\langle$P2]}}
                \texttt{HAMLE: A, (4, 5)} \\
            }
        }
    }
     \\ \\

    \theutterance \stepcounter{utterance}  
    & & & \multicolumn{2}{p{0.3\linewidth}}{
        \cellcolor[rgb]{0.9,0.9,0.9}{
            \makecell[{{p{\linewidth}}}]{
                \texttt{\tiny{[GM$|$GM]}}
                \texttt{HAMLE: A, (4, 5)} \\
            }
        }
    }
    & & \\ \\

    \theutterance \stepcounter{utterance}  
    & & & \multicolumn{2}{p{0.3\linewidth}}{
        \cellcolor[rgb]{0.9,0.9,0.9}{
            \makecell[{{p{\linewidth}}}]{
                \texttt{\tiny{[GM$|$GM]}}
                \texttt{Ceza: (4,5) boş değil; "│" içeriyor.} \\
\\ 
\texttt{Toplamda 16 cezanın 32 kadarını biriktirdiniz.Lütfen tekrar deneyin!} \\
            }
        }
    }
    & & \\ \\

    \theutterance \stepcounter{utterance}  
    & & \multicolumn{4}{p{0.6\linewidth}}{
        \cellcolor[rgb]{0.9,0.9,0.9}{
            \makecell[{{p{\linewidth}}}]{
                \texttt{\tiny{[GM$\rangle$P2]}}
                \texttt{Ceza: (4,5) boş değil; "│" içeriyor.} \\
\\ 
\texttt{Toplamda 16 cezanın 32 kadarını biriktirdiniz.Lütfen tekrar deneyin!} \\
            }
        }
    }
    & \\ \\

    \theutterance \stepcounter{utterance}  
    & & & \multicolumn{4}{p{0.6\linewidth}}{
        \cellcolor[rgb]{0.9,0.9,0.9}{
            \makecell[{{p{\linewidth}}}]{
                \texttt{\tiny{[GM$\langle$P2]}}
                \texttt{SÖYLE: İki nesnemiz (S ve O) 4. satıra geldi. Şimdi diğer nesnelerin (A, M, D, T, P) 4. satıra nasıl taşınacağını planlayalım. Öncelikle hangi nesnelerin hangi X koordinatlarıyla taşınabileceğini kontrol edelim.} \\
            }
        }
    }
     \\ \\

    \theutterance \stepcounter{utterance}  
    & & & \multicolumn{2}{p{0.3\linewidth}}{
        \cellcolor[rgb]{0.9,0.9,0.9}{
            \makecell[{{p{\linewidth}}}]{
                \texttt{\tiny{[GM$|$GM]}}
                \texttt{SÖYLE: İki nesnemiz (S ve O) 4. satıra geldi. Şimdi diğer nesnelerin (A, M, D, T, P) 4. satıra nasıl taşınacağını planlayalım. Öncelikle hangi nesnelerin hangi X koordinatlarıyla taşınabileceğini kontrol edelim.} \\
            }
        }
    }
    & & \\ \\

    \theutterance \stepcounter{utterance}  
    & & \multicolumn{4}{p{0.6\linewidth}}{
        \cellcolor[rgb]{0.9,0.9,0.9}{
            \makecell[{{p{\linewidth}}}]{
                \texttt{\tiny{[P1$\langle$GM]}}
                \texttt{D’ başarıyla (4,1) koordinatlarına taşındı. Güncellenmiş ızgaran şu şekilde:} \\
\texttt{```} \\
\texttt{ 1234567} \\
\texttt{╔═══════╗} \\
\texttt{║◌◌◌D◌T◌║ 1} \\
\texttt{║P◌◌S◌◌◌║ 2} \\
\texttt{║◌◌A◌◌┌─╢ 3} \\
\texttt{║◌◌O┌─┤◌║ 4} \\
\texttt{║◌◌◌│◌├─╢ 5} \\
\texttt{╟───┼─┘◌║ 6} \\
\texttt{║◌M◌│◌◌◌║ 7} \\
\texttt{╚═══╧═══╝} \\
\texttt{```} \\
\texttt{Şu anda en fazla 28 turdan 17. turu oynuyorsunuz.} \\
\texttt{Toplamda 16 cezanın 32 kadarını biriktirdiniz.Diğer oyuncu şu mesajı gönderdi:} \\
\\ 
\\ 
\texttt{"İki nesnemiz (S ve O) 4. satıra geldi. Şimdi diğer nesnelerin (A, M, D, T, P) 4. satıra nasıl taşınacağını planlayalım. Öncelikle hangi nesnelerin hangi X koordinatlarıyla taşınabileceğini kontrol edelim."} \\
\\ 
\\ 
\texttt{Bir sonraki komutun nedir?} \\
            }
        }
    }
    & \\ \\

    \theutterance \stepcounter{utterance}  
    & \multicolumn{4}{p{0.6\linewidth}}{
        \cellcolor[rgb]{0.9,0.9,0.9}{
            \makecell[{{p{\linewidth}}}]{
                \texttt{\tiny{[P1$\rangle$GM]}}
                \texttt{SÖYLE: A'nın mevcut konumunu kontrol edelim.} \\
            }
        }
    }
    & & \\ \\

    \theutterance \stepcounter{utterance}  
    & & & \multicolumn{2}{p{0.3\linewidth}}{
        \cellcolor[rgb]{0.9,0.9,0.9}{
            \makecell[{{p{\linewidth}}}]{
                \texttt{\tiny{[GM$|$GM]}}
                \texttt{SÖYLE: A'nın mevcut konumunu kontrol edelim.} \\
            }
        }
    }
    & & \\ \\

    \theutterance \stepcounter{utterance}  
    & & \multicolumn{4}{p{0.6\linewidth}}{
        \cellcolor[rgb]{0.9,0.9,0.9}{
            \makecell[{{p{\linewidth}}}]{
                \texttt{\tiny{[GM$\rangle$P2]}}
                \texttt{Mesajın diğer oyuncuya iletildi.} \\
\\ 
\\ 
\texttt{Şu anda en fazla 28 turdan 18. turu oynuyorsunuz.} \\
\texttt{Toplamda 16 cezanın 32 kadarını biriktirdiniz.Diğer oyuncu şu mesajı gönderdi:} \\
\\ 
\\ 
\texttt{"A'nın mevcut konumunu kontrol edelim."} \\
\\ 
\\ 
\texttt{Bir sonraki komutun nedir?} \\
            }
        }
    }
    & \\ \\

    \theutterance \stepcounter{utterance}  
    & & & \multicolumn{4}{p{0.6\linewidth}}{
        \cellcolor[rgb]{0.9,0.9,0.9}{
            \makecell[{{p{\linewidth}}}]{
                \texttt{\tiny{[GM$\langle$P2]}}
                \texttt{SÖYLE: A'nın şu anki konumunu öğrenelim. Oyuncunun belirttiği gibi kontrollere başlayalım.} \\
            }
        }
    }
     \\ \\

    \theutterance \stepcounter{utterance}  
    & & & \multicolumn{2}{p{0.3\linewidth}}{
        \cellcolor[rgb]{0.9,0.9,0.9}{
            \makecell[{{p{\linewidth}}}]{
                \texttt{\tiny{[GM$|$GM]}}
                \texttt{SÖYLE: A'nın şu anki konumunu öğrenelim. Oyuncunun belirttiği gibi kontrollere başlayalım.} \\
            }
        }
    }
    & & \\ \\

    \theutterance \stepcounter{utterance}  
    & & \multicolumn{4}{p{0.6\linewidth}}{
        \cellcolor[rgb]{0.9,0.9,0.9}{
            \makecell[{{p{\linewidth}}}]{
                \texttt{\tiny{[P1$\langle$GM]}}
                \texttt{Mesajın diğer oyuncuya iletildi.} \\
\\ 
\\ 
\texttt{Şu anda en fazla 28 turdan 18. turu oynuyorsunuz.} \\
\texttt{Toplamda 16 cezanın 32 kadarını biriktirdiniz.Diğer oyuncu şu mesajı gönderdi:} \\
\\ 
\\ 
\texttt{"A'nın şu anki konumunu öğrenelim. Oyuncunun belirttiği gibi kontrollere başlayalım."} \\
\\ 
\\ 
\texttt{Bir sonraki komutun nedir?} \\
            }
        }
    }
    & \\ \\

    \theutterance \stepcounter{utterance}  
    & \multicolumn{4}{p{0.6\linewidth}}{
        \cellcolor[rgb]{0.9,0.9,0.9}{
            \makecell[{{p{\linewidth}}}]{
                \texttt{\tiny{[P1$\rangle$GM]}}
                \texttt{HAMLE: A, (3, 3)} \\
            }
        }
    }
    & & \\ \\

    \theutterance \stepcounter{utterance}  
    & & & \multicolumn{2}{p{0.3\linewidth}}{
        \cellcolor[rgb]{0.9,0.9,0.9}{
            \makecell[{{p{\linewidth}}}]{
                \texttt{\tiny{[GM$|$GM]}}
                \texttt{HAMLE: A, (3, 3)} \\
            }
        }
    }
    & & \\ \\

    \theutterance \stepcounter{utterance}  
    & & & \multicolumn{2}{p{0.3\linewidth}}{
        \cellcolor[rgb]{0.9,0.9,0.9}{
            \makecell[{{p{\linewidth}}}]{
                \texttt{\tiny{[GM$|$GM]}}
                \texttt{Ceza: (3,3) boş değil; "A" içeriyor.} \\
\\ 
\texttt{Toplamda 16 cezanın 33 kadarını biriktirdiniz.Lütfen tekrar deneyin!} \\
            }
        }
    }
    & & \\ \\

    \theutterance \stepcounter{utterance}  
    & & \multicolumn{4}{p{0.6\linewidth}}{
        \cellcolor[rgb]{0.9,0.9,0.9}{
            \makecell[{{p{\linewidth}}}]{
                \texttt{\tiny{[P1$\langle$GM]}}
                \texttt{Ceza: (3,3) boş değil; "A" içeriyor.} \\
\\ 
\texttt{Toplamda 16 cezanın 33 kadarını biriktirdiniz.Lütfen tekrar deneyin!} \\
            }
        }
    }
    & \\ \\

    \theutterance \stepcounter{utterance}  
    & \multicolumn{4}{p{0.6\linewidth}}{
        \cellcolor[rgb]{0.9,0.9,0.9}{
            \makecell[{{p{\linewidth}}}]{
                \texttt{\tiny{[P1$\rangle$GM]}}
                \texttt{HAMLE: A, (5, 3)} \\
            }
        }
    }
    & & \\ \\

    \theutterance \stepcounter{utterance}  
    & & & \multicolumn{2}{p{0.3\linewidth}}{
        \cellcolor[rgb]{0.9,0.9,0.9}{
            \makecell[{{p{\linewidth}}}]{
                \texttt{\tiny{[GM$|$GM]}}
                \texttt{HAMLE: A, (5, 3)} \\
            }
        }
    }
    & & \\ \\

    \theutterance \stepcounter{utterance}  
    & & & \multicolumn{2}{p{0.3\linewidth}}{
        \cellcolor[rgb]{0.9,0.9,0.9}{
            \makecell[{{p{\linewidth}}}]{
                \texttt{\tiny{[GM$|$GM]}}
                \texttt{A moved to (5, 3)} \\
            }
        }
    }
    & & \\ \\

    \theutterance \stepcounter{utterance}  
    & & \multicolumn{4}{p{0.6\linewidth}}{
        \cellcolor[rgb]{0.9,0.9,0.9}{
            \makecell[{{p{\linewidth}}}]{
                \texttt{\tiny{[GM$\rangle$P2]}}
                \texttt{Mesajın diğer oyuncuya iletildi.} \\
\\ 
\\ 
\texttt{Şu anda en fazla 28 turdan 19. turu oynuyorsunuz.} \\
\texttt{Toplamda 16 cezanın 33 kadarını biriktirdiniz.Diğer oyuncu kendi ızgarasında bir nesneyi hareket ettirdi.} \\
\\ 
\\ 
\texttt{Bir sonraki komutun nedir?} \\
            }
        }
    }
    & \\ \\

    \theutterance \stepcounter{utterance}  
    & & & \multicolumn{4}{p{0.6\linewidth}}{
        \cellcolor[rgb]{0.9,0.9,0.9}{
            \makecell[{{p{\linewidth}}}]{
                \texttt{\tiny{[GM$\langle$P2]}}
                \texttt{HAMLE: A, (4, 6)} \\
            }
        }
    }
     \\ \\

    \theutterance \stepcounter{utterance}  
    & & & \multicolumn{2}{p{0.3\linewidth}}{
        \cellcolor[rgb]{0.9,0.9,0.9}{
            \makecell[{{p{\linewidth}}}]{
                \texttt{\tiny{[GM$|$GM]}}
                \texttt{HAMLE: A, (4, 6)} \\
            }
        }
    }
    & & \\ \\

    \theutterance \stepcounter{utterance}  
    & & & \multicolumn{2}{p{0.3\linewidth}}{
        \cellcolor[rgb]{0.9,0.9,0.9}{
            \makecell[{{p{\linewidth}}}]{
                \texttt{\tiny{[GM$|$GM]}}
                \texttt{Ceza: (4,6) boş değil; "┼" içeriyor.} \\
\\ 
\texttt{Toplamda 16 cezanın 34 kadarını biriktirdiniz.Lütfen tekrar deneyin!} \\
            }
        }
    }
    & & \\ \\

    \theutterance \stepcounter{utterance}  
    & & \multicolumn{4}{p{0.6\linewidth}}{
        \cellcolor[rgb]{0.9,0.9,0.9}{
            \makecell[{{p{\linewidth}}}]{
                \texttt{\tiny{[GM$\rangle$P2]}}
                \texttt{Ceza: (4,6) boş değil; "┼" içeriyor.} \\
\\ 
\texttt{Toplamda 16 cezanın 34 kadarını biriktirdiniz.Lütfen tekrar deneyin!} \\
            }
        }
    }
    & \\ \\

    \theutterance \stepcounter{utterance}  
    & & & \multicolumn{4}{p{0.6\linewidth}}{
        \cellcolor[rgb]{0.9,0.9,0.9}{
            \makecell[{{p{\linewidth}}}]{
                \texttt{\tiny{[GM$\langle$P2]}}
                \texttt{SÖYLE: A'nın mevcut konumunu kontrol edelim. A'nın şu an hangi X ve Y koordinatında olduğunu belirleyelim; ardından 4. satırda boş hücre olup olmadığını kontrol edelim ve uygun X ile A'yı 4. satıra taşımayı deneyelim.} \\
            }
        }
    }
     \\ \\

    \theutterance \stepcounter{utterance}  
    & & & \multicolumn{2}{p{0.3\linewidth}}{
        \cellcolor[rgb]{0.9,0.9,0.9}{
            \makecell[{{p{\linewidth}}}]{
                \texttt{\tiny{[GM$|$GM]}}
                \texttt{SÖYLE: A'nın mevcut konumunu kontrol edelim. A'nın şu an hangi X ve Y koordinatında olduğunu belirleyelim; ardından 4. satırda boş hücre olup olmadığını kontrol edelim ve uygun X ile A'yı 4. satıra taşımayı deneyelim.} \\
            }
        }
    }
    & & \\ \\

    \theutterance \stepcounter{utterance}  
    & & \multicolumn{4}{p{0.6\linewidth}}{
        \cellcolor[rgb]{0.9,0.9,0.9}{
            \makecell[{{p{\linewidth}}}]{
                \texttt{\tiny{[P1$\langle$GM]}}
                \texttt{A’ başarıyla (5,3) koordinatlarına taşındı. Güncellenmiş ızgaran şu şekilde:} \\
\texttt{```} \\
\texttt{ 1234567} \\
\texttt{╔═══════╗} \\
\texttt{║◌◌◌D◌T◌║ 1} \\
\texttt{║P◌◌S◌◌◌║ 2} \\
\texttt{║◌◌◌◌A┌─╢ 3} \\
\texttt{║◌◌O┌─┤◌║ 4} \\
\texttt{║◌◌◌│◌├─╢ 5} \\
\texttt{╟───┼─┘◌║ 6} \\
\texttt{║◌M◌│◌◌◌║ 7} \\
\texttt{╚═══╧═══╝} \\
\texttt{```} \\
\texttt{Şu anda en fazla 28 turdan 19. turu oynuyorsunuz.} \\
\texttt{Toplamda 16 cezanın 34 kadarını biriktirdiniz.Diğer oyuncu şu mesajı gönderdi:} \\
\\ 
\\ 
\texttt{"A'nın mevcut konumunu kontrol edelim. A'nın şu an hangi X ve Y koordinatında olduğunu belirleyelim; ardından 4. satırda boş hücre olup olmadığını kontrol edelim ve uygun X ile A'yı 4. satıra taşımayı deneyelim."} \\
\\ 
\\ 
\texttt{Bir sonraki komutun nedir?} \\
            }
        }
    }
    & \\ \\

    \theutterance \stepcounter{utterance}  
    & \multicolumn{4}{p{0.6\linewidth}}{
        \cellcolor[rgb]{0.9,0.9,0.9}{
            \makecell[{{p{\linewidth}}}]{
                \texttt{\tiny{[P1$\rangle$GM]}}
                \texttt{HAMLE: A, (4, 3)} \\
            }
        }
    }
    & & \\ \\

    \theutterance \stepcounter{utterance}  
    & & & \multicolumn{2}{p{0.3\linewidth}}{
        \cellcolor[rgb]{0.9,0.9,0.9}{
            \makecell[{{p{\linewidth}}}]{
                \texttt{\tiny{[GM$|$GM]}}
                \texttt{HAMLE: A, (4, 3)} \\
            }
        }
    }
    & & \\ \\

    \theutterance \stepcounter{utterance}  
    & & & \multicolumn{2}{p{0.3\linewidth}}{
        \cellcolor[rgb]{0.9,0.9,0.9}{
            \makecell[{{p{\linewidth}}}]{
                \texttt{\tiny{[GM$|$GM]}}
                \texttt{A moved to (4, 3)} \\
            }
        }
    }
    & & \\ \\

    \theutterance \stepcounter{utterance}  
    & & \multicolumn{4}{p{0.6\linewidth}}{
        \cellcolor[rgb]{0.9,0.9,0.9}{
            \makecell[{{p{\linewidth}}}]{
                \texttt{\tiny{[GM$\rangle$P2]}}
                \texttt{Mesajın diğer oyuncuya iletildi.} \\
\\ 
\\ 
\texttt{Şu anda en fazla 28 turdan 20. turu oynuyorsunuz.} \\
\texttt{Toplamda 16 cezanın 34 kadarını biriktirdiniz.Diğer oyuncu kendi ızgarasında bir nesneyi hareket ettirdi.} \\
\\ 
\\ 
\texttt{Bir sonraki komutun nedir?} \\
            }
        }
    }
    & \\ \\

    \theutterance \stepcounter{utterance}  
    & & & \multicolumn{4}{p{0.6\linewidth}}{
        \cellcolor[rgb]{0.9,0.9,0.9}{
            \makecell[{{p{\linewidth}}}]{
                \texttt{\tiny{[GM$\langle$P2]}}
                \texttt{HAMLE: A, (4, 7)} \\
            }
        }
    }
     \\ \\

    \theutterance \stepcounter{utterance}  
    & & & \multicolumn{2}{p{0.3\linewidth}}{
        \cellcolor[rgb]{0.9,0.9,0.9}{
            \makecell[{{p{\linewidth}}}]{
                \texttt{\tiny{[GM$|$GM]}}
                \texttt{HAMLE: A, (4, 7)} \\
            }
        }
    }
    & & \\ \\

    \theutterance \stepcounter{utterance}  
    & & & \multicolumn{2}{p{0.3\linewidth}}{
        \cellcolor[rgb]{0.9,0.9,0.9}{
            \makecell[{{p{\linewidth}}}]{
                \texttt{\tiny{[GM$|$GM]}}
                \texttt{Ceza: (4,7) boş değil; "│" içeriyor.} \\
\\ 
\texttt{Toplamda 16 cezanın 35 kadarını biriktirdiniz.Lütfen tekrar deneyin!} \\
            }
        }
    }
    & & \\ \\

    \theutterance \stepcounter{utterance}  
    & & \multicolumn{4}{p{0.6\linewidth}}{
        \cellcolor[rgb]{0.9,0.9,0.9}{
            \makecell[{{p{\linewidth}}}]{
                \texttt{\tiny{[GM$\rangle$P2]}}
                \texttt{Ceza: (4,7) boş değil; "│" içeriyor.} \\
\\ 
\texttt{Toplamda 16 cezanın 35 kadarını biriktirdiniz.Lütfen tekrar deneyin!} \\
            }
        }
    }
    & \\ \\

    \theutterance \stepcounter{utterance}  
    & & & \multicolumn{4}{p{0.6\linewidth}}{
        \cellcolor[rgb]{0.9,0.9,0.9}{
            \makecell[{{p{\linewidth}}}]{
                \texttt{\tiny{[GM$\langle$P2]}}
                \texttt{SÖYLE: A'nın mevcut konumunu ve 4. satırda boş hücre olup olmadığını kontrol edelim. A şu anda (4, 7)'de "│", (4, 6)'da "┼", (4, 5)'de "│", (4, 4)'de "┌", (4, 3)'de "M" var. Yani 4. satırda A'yı taşımak için uygun X ve Y koordinatı bulmalıyız. Örn. (4, 2) ya da (4, 3) ya da (4, 4) ya da (4, 5) ya da (4, 6) ya da (4, 7) dışında hangi hücre boş? Kontrol edelim.} \\
            }
        }
    }
     \\ \\

    \theutterance \stepcounter{utterance}  
    & & & \multicolumn{2}{p{0.3\linewidth}}{
        \cellcolor[rgb]{0.9,0.9,0.9}{
            \makecell[{{p{\linewidth}}}]{
                \texttt{\tiny{[GM$|$GM]}}
                \texttt{SÖYLE: A'nın mevcut konumunu ve 4. satırda boş hücre olup olmadığını kontrol edelim. A şu anda (4, 7)'de "│", (4, 6)'da "┼", (4, 5)'de "│", (4, 4)'de "┌", (4, 3)'de "M" var. Yani 4. satırda A'yı taşımak için uygun X ve Y koordinatı bulmalıyız. Örn. (4, 2) ya da (4, 3) ya da (4, 4) ya da (4, 5) ya da (4, 6) ya da (4, 7) dışında hangi hücre boş? Kontrol edelim.} \\
            }
        }
    }
    & & \\ \\

    \theutterance \stepcounter{utterance}  
    & & \multicolumn{4}{p{0.6\linewidth}}{
        \cellcolor[rgb]{0.9,0.9,0.9}{
            \makecell[{{p{\linewidth}}}]{
                \texttt{\tiny{[P1$\langle$GM]}}
                \texttt{A’ başarıyla (4,3) koordinatlarına taşındı. Güncellenmiş ızgaran şu şekilde:} \\
\texttt{```} \\
\texttt{ 1234567} \\
\texttt{╔═══════╗} \\
\texttt{║◌◌◌D◌T◌║ 1} \\
\texttt{║P◌◌S◌◌◌║ 2} \\
\texttt{║◌◌◌A◌┌─╢ 3} \\
\texttt{║◌◌O┌─┤◌║ 4} \\
\texttt{║◌◌◌│◌├─╢ 5} \\
\texttt{╟───┼─┘◌║ 6} \\
\texttt{║◌M◌│◌◌◌║ 7} \\
\texttt{╚═══╧═══╝} \\
\texttt{```} \\
\texttt{Şu anda en fazla 28 turdan 20. turu oynuyorsunuz.} \\
\texttt{Toplamda 16 cezanın 35 kadarını biriktirdiniz.Diğer oyuncu şu mesajı gönderdi:} \\
\\ 
\\ 
\texttt{"A'nın mevcut konumunu ve 4. satırda boş hücre olup olmadığını kontrol edelim. A şu anda (4, 7)'de "│", (4, 6)'da "┼", (4, 5)'de "│", (4, 4)'de "┌", (4, 3)'de "M" var. Yani 4. satırda A'yı taşımak için uygun X ve Y koordinatı bulmalıyız. Örn. (4, 2) ya da (4, 3) ya da (4, 4) ya da (4, 5) ya da (4, 6) ya da (4, 7) dışında hangi hücre boş? Kontrol edelim."} \\
\\ 
\\ 
\texttt{Bir sonraki komutun nedir?} \\
            }
        }
    }
    & \\ \\

    \theutterance \stepcounter{utterance}  
    & \multicolumn{4}{p{0.6\linewidth}}{
        \cellcolor[rgb]{0.9,0.9,0.9}{
            \makecell[{{p{\linewidth}}}]{
                \texttt{\tiny{[P1$\rangle$GM]}}
                \texttt{SÖYLE: 4. satırda (4, 2) hücresinin hangi sembolü içerdiğini kontrol edelim.} \\
            }
        }
    }
    & & \\ \\

    \theutterance \stepcounter{utterance}  
    & & & \multicolumn{2}{p{0.3\linewidth}}{
        \cellcolor[rgb]{0.9,0.9,0.9}{
            \makecell[{{p{\linewidth}}}]{
                \texttt{\tiny{[GM$|$GM]}}
                \texttt{SÖYLE: 4. satırda (4, 2) hücresinin hangi sembolü içerdiğini kontrol edelim.} \\
            }
        }
    }
    & & \\ \\

    \theutterance \stepcounter{utterance}  
    & & \multicolumn{4}{p{0.6\linewidth}}{
        \cellcolor[rgb]{0.9,0.9,0.9}{
            \makecell[{{p{\linewidth}}}]{
                \texttt{\tiny{[GM$\rangle$P2]}}
                \texttt{Mesajın diğer oyuncuya iletildi.} \\
\\ 
\\ 
\texttt{Şu anda en fazla 28 turdan 21. turu oynuyorsunuz.} \\
\texttt{Toplamda 16 cezanın 35 kadarını biriktirdiniz.Diğer oyuncu şu mesajı gönderdi:} \\
\\ 
\\ 
\texttt{"4. satırda (4, 2) hücresinin hangi sembolü içerdiğini kontrol edelim."} \\
\\ 
\\ 
\texttt{Bir sonraki komutun nedir?} \\
            }
        }
    }
    & \\ \\

    \theutterance \stepcounter{utterance}  
    & & & \multicolumn{4}{p{0.6\linewidth}}{
        \cellcolor[rgb]{0.9,0.9,0.9}{
            \makecell[{{p{\linewidth}}}]{
                \texttt{\tiny{[GM$\langle$P2]}}
                \texttt{SÖYLE: 4. satırda (4, 2)'yi kontrol edelim. Hangi nesneyi ya da sembolü içerdiğini öğrenelim, böylece A'yı taşımak için uygun X değerini belirleyebiliriz.} \\
            }
        }
    }
     \\ \\

    \theutterance \stepcounter{utterance}  
    & & & \multicolumn{2}{p{0.3\linewidth}}{
        \cellcolor[rgb]{0.9,0.9,0.9}{
            \makecell[{{p{\linewidth}}}]{
                \texttt{\tiny{[GM$|$GM]}}
                \texttt{SÖYLE: 4. satırda (4, 2)'yi kontrol edelim. Hangi nesneyi ya da sembolü içerdiğini öğrenelim, böylece A'yı taşımak için uygun X değerini belirleyebiliriz.} \\
            }
        }
    }
    & & \\ \\

    \theutterance \stepcounter{utterance}  
    & & \multicolumn{4}{p{0.6\linewidth}}{
        \cellcolor[rgb]{0.9,0.9,0.9}{
            \makecell[{{p{\linewidth}}}]{
                \texttt{\tiny{[P1$\langle$GM]}}
                \texttt{Mesajın diğer oyuncuya iletildi.} \\
\\ 
\\ 
\texttt{Şu anda en fazla 28 turdan 21. turu oynuyorsunuz.} \\
\texttt{Toplamda 16 cezanın 35 kadarını biriktirdiniz.Diğer oyuncu şu mesajı gönderdi:} \\
\\ 
\\ 
\texttt{"4. satırda (4, 2)'yi kontrol edelim. Hangi nesneyi ya da sembolü içerdiğini öğrenelim, böylece A'yı taşımak için uygun X değerini belirleyebiliriz."} \\
\\ 
\\ 
\texttt{Bir sonraki komutun nedir?} \\
            }
        }
    }
    & \\ \\

    \theutterance \stepcounter{utterance}  
    & \multicolumn{4}{p{0.6\linewidth}}{
        \cellcolor[rgb]{0.9,0.9,0.9}{
            \makecell[{{p{\linewidth}}}]{
                \texttt{\tiny{[P1$\rangle$GM]}}
                \texttt{SÖYLE: 4. satırda (4, 2) 'S' içeriyor.} \\
            }
        }
    }
    & & \\ \\

    \theutterance \stepcounter{utterance}  
    & & & \multicolumn{2}{p{0.3\linewidth}}{
        \cellcolor[rgb]{0.9,0.9,0.9}{
            \makecell[{{p{\linewidth}}}]{
                \texttt{\tiny{[GM$|$GM]}}
                \texttt{SÖYLE: 4. satırda (4, 2) 'S' içeriyor.} \\
            }
        }
    }
    & & \\ \\

    \theutterance \stepcounter{utterance}  
    & & \multicolumn{4}{p{0.6\linewidth}}{
        \cellcolor[rgb]{0.9,0.9,0.9}{
            \makecell[{{p{\linewidth}}}]{
                \texttt{\tiny{[GM$\rangle$P2]}}
                \texttt{Mesajın diğer oyuncuya iletildi.} \\
\\ 
\\ 
\texttt{Şu anda en fazla 28 turdan 22. turu oynuyorsunuz.} \\
\texttt{Toplamda 16 cezanın 35 kadarını biriktirdiniz.Diğer oyuncu şu mesajı gönderdi:} \\
\\ 
\\ 
\texttt{"4. satırda (4, 2) 'S' içeriyor."} \\
\\ 
\\ 
\texttt{Bir sonraki komutun nedir?} \\
            }
        }
    }
    & \\ \\

    \theutterance \stepcounter{utterance}  
    & & & \multicolumn{4}{p{0.6\linewidth}}{
        \cellcolor[rgb]{0.9,0.9,0.9}{
            \makecell[{{p{\linewidth}}}]{
                \texttt{\tiny{[GM$\langle$P2]}}
                \texttt{HAMLE: A, (4, 3)} \\
            }
        }
    }
     \\ \\

    \theutterance \stepcounter{utterance}  
    & & & \multicolumn{2}{p{0.3\linewidth}}{
        \cellcolor[rgb]{0.9,0.9,0.9}{
            \makecell[{{p{\linewidth}}}]{
                \texttt{\tiny{[GM$|$GM]}}
                \texttt{HAMLE: A, (4, 3)} \\
            }
        }
    }
    & & \\ \\

    \theutterance \stepcounter{utterance}  
    & & & \multicolumn{2}{p{0.3\linewidth}}{
        \cellcolor[rgb]{0.9,0.9,0.9}{
            \makecell[{{p{\linewidth}}}]{
                \texttt{\tiny{[GM$|$GM]}}
                \texttt{Ceza: (4,3) boş değil; "M" içeriyor.} \\
\\ 
\texttt{Toplamda 16 cezanın 36 kadarını biriktirdiniz.Lütfen tekrar deneyin!} \\
            }
        }
    }
    & & \\ \\

    \theutterance \stepcounter{utterance}  
    & & \multicolumn{4}{p{0.6\linewidth}}{
        \cellcolor[rgb]{0.9,0.9,0.9}{
            \makecell[{{p{\linewidth}}}]{
                \texttt{\tiny{[GM$\rangle$P2]}}
                \texttt{Ceza: (4,3) boş değil; "M" içeriyor.} \\
\\ 
\texttt{Toplamda 16 cezanın 36 kadarını biriktirdiniz.Lütfen tekrar deneyin!} \\
            }
        }
    }
    & \\ \\

    \theutterance \stepcounter{utterance}  
    & & & \multicolumn{4}{p{0.6\linewidth}}{
        \cellcolor[rgb]{0.9,0.9,0.9}{
            \makecell[{{p{\linewidth}}}]{
                \texttt{\tiny{[GM$\langle$P2]}}
                \texttt{SÖYLE: 4. satırda uygun boş hücre olup olmadığını kontrol edelim. Örn. (4, 1), (4, 5), (4, 6), (4, 7). Her birinin hangi sembolü içerdiğini kontrol edelim ve A'yı taşımak için doğru X değerini belirleyelim.} \\
            }
        }
    }
     \\ \\

    \theutterance \stepcounter{utterance}  
    & & & \multicolumn{2}{p{0.3\linewidth}}{
        \cellcolor[rgb]{0.9,0.9,0.9}{
            \makecell[{{p{\linewidth}}}]{
                \texttt{\tiny{[GM$|$GM]}}
                \texttt{SÖYLE: 4. satırda uygun boş hücre olup olmadığını kontrol edelim. Örn. (4, 1), (4, 5), (4, 6), (4, 7). Her birinin hangi sembolü içerdiğini kontrol edelim ve A'yı taşımak için doğru X değerini belirleyelim.} \\
            }
        }
    }
    & & \\ \\

    \theutterance \stepcounter{utterance}  
    & & \multicolumn{4}{p{0.6\linewidth}}{
        \cellcolor[rgb]{0.9,0.9,0.9}{
            \makecell[{{p{\linewidth}}}]{
                \texttt{\tiny{[P1$\langle$GM]}}
                \texttt{Mesajın diğer oyuncuya iletildi.} \\
\\ 
\\ 
\texttt{Şu anda en fazla 28 turdan 22. turu oynuyorsunuz.} \\
\texttt{Toplamda 16 cezanın 36 kadarını biriktirdiniz.Diğer oyuncu şu mesajı gönderdi:} \\
\\ 
\\ 
\texttt{"4. satırda uygun boş hücre olup olmadığını kontrol edelim. Örn. (4, 1), (4, 5), (4, 6), (4, 7). Her birinin hangi sembolü içerdiğini kontrol edelim ve A'yı taşımak için doğru X değerini belirleyelim."} \\
\\ 
\\ 
\texttt{Bir sonraki komutun nedir?} \\
            }
        }
    }
    & \\ \\

    \theutterance \stepcounter{utterance}  
    & \multicolumn{4}{p{0.6\linewidth}}{
        \cellcolor[rgb]{0.9,0.9,0.9}{
            \makecell[{{p{\linewidth}}}]{
                \texttt{\tiny{[P1$\rangle$GM]}}
                \texttt{SÖYLE: 4. satırda (4, 1) 'T' içeriyor, (4, 5) '│', (4, 6) '┼', (4, 7) '│' içeriyor.} \\
            }
        }
    }
    & & \\ \\

    \theutterance \stepcounter{utterance}  
    & & & \multicolumn{2}{p{0.3\linewidth}}{
        \cellcolor[rgb]{0.9,0.9,0.9}{
            \makecell[{{p{\linewidth}}}]{
                \texttt{\tiny{[GM$|$GM]}}
                \texttt{SÖYLE: 4. satırda (4, 1) 'T' içeriyor, (4, 5) '│', (4, 6) '┼', (4, 7) '│' içeriyor.} \\
            }
        }
    }
    & & \\ \\

    \theutterance \stepcounter{utterance}  
    & & \multicolumn{4}{p{0.6\linewidth}}{
        \cellcolor[rgb]{0.9,0.9,0.9}{
            \makecell[{{p{\linewidth}}}]{
                \texttt{\tiny{[GM$\rangle$P2]}}
                \texttt{Mesajın diğer oyuncuya iletildi.} \\
\\ 
\\ 
\texttt{Şu anda en fazla 28 turdan 23. turu oynuyorsunuz.} \\
\texttt{Toplamda 16 cezanın 36 kadarını biriktirdiniz.Diğer oyuncu şu mesajı gönderdi:} \\
\\ 
\\ 
\texttt{"4. satırda (4, 1) 'T' içeriyor, (4, 5) '│', (4, 6) '┼', (4, 7) '│' içeriyor."} \\
\\ 
\\ 
\texttt{Bir sonraki komutun nedir?} \\
            }
        }
    }
    & \\ \\

    \theutterance \stepcounter{utterance}  
    & & & \multicolumn{4}{p{0.6\linewidth}}{
        \cellcolor[rgb]{0.9,0.9,0.9}{
            \makecell[{{p{\linewidth}}}]{
                \texttt{\tiny{[GM$\langle$P2]}}
                \texttt{HAMLE: A, (4, 4)} \\
            }
        }
    }
     \\ \\

    \theutterance \stepcounter{utterance}  
    & & & \multicolumn{2}{p{0.3\linewidth}}{
        \cellcolor[rgb]{0.9,0.9,0.9}{
            \makecell[{{p{\linewidth}}}]{
                \texttt{\tiny{[GM$|$GM]}}
                \texttt{HAMLE: A, (4, 4)} \\
            }
        }
    }
    & & \\ \\

    \theutterance \stepcounter{utterance}  
    & & & \multicolumn{2}{p{0.3\linewidth}}{
        \cellcolor[rgb]{0.9,0.9,0.9}{
            \makecell[{{p{\linewidth}}}]{
                \texttt{\tiny{[GM$|$GM]}}
                \texttt{Ceza: (4,4) boş değil; "┌" içeriyor.} \\
\\ 
\texttt{Toplamda 16 cezanın 37 kadarını biriktirdiniz.Lütfen tekrar deneyin!} \\
            }
        }
    }
    & & \\ \\

    \theutterance \stepcounter{utterance}  
    & & \multicolumn{4}{p{0.6\linewidth}}{
        \cellcolor[rgb]{0.9,0.9,0.9}{
            \makecell[{{p{\linewidth}}}]{
                \texttt{\tiny{[GM$\rangle$P2]}}
                \texttt{Ceza: (4,4) boş değil; "┌" içeriyor.} \\
\\ 
\texttt{Toplamda 16 cezanın 37 kadarını biriktirdiniz.Lütfen tekrar deneyin!} \\
            }
        }
    }
    & \\ \\

    \theutterance \stepcounter{utterance}  
    & & & \multicolumn{4}{p{0.6\linewidth}}{
        \cellcolor[rgb]{0.9,0.9,0.9}{
            \makecell[{{p{\linewidth}}}]{
                \texttt{\tiny{[GM$\langle$P2]}}
                \texttt{SÖYLE: 4. satırda uygun boş hücre olup olmadığını kontrol edelim. Hangi X koordinatlarının içeriği boş değilse, hangisinin boş olup olmadığını kontrol edelim. Örn. (4, 1) T, (4, 2) S, (4, 3) M, (4, 4) ┌. Yani 4. satırda hücrelerin içeriği: S, M, ┌ dışındaki hücreler nerede? Kontrol edelim.} \\
            }
        }
    }
     \\ \\

    \theutterance \stepcounter{utterance}  
    & & & \multicolumn{2}{p{0.3\linewidth}}{
        \cellcolor[rgb]{0.9,0.9,0.9}{
            \makecell[{{p{\linewidth}}}]{
                \texttt{\tiny{[GM$|$GM]}}
                \texttt{SÖYLE: 4. satırda uygun boş hücre olup olmadığını kontrol edelim. Hangi X koordinatlarının içeriği boş değilse, hangisinin boş olup olmadığını kontrol edelim. Örn. (4, 1) T, (4, 2) S, (4, 3) M, (4, 4) ┌. Yani 4. satırda hücrelerin içeriği: S, M, ┌ dışındaki hücreler nerede? Kontrol edelim.} \\
            }
        }
    }
    & & \\ \\

    \theutterance \stepcounter{utterance}  
    & & \multicolumn{4}{p{0.6\linewidth}}{
        \cellcolor[rgb]{0.9,0.9,0.9}{
            \makecell[{{p{\linewidth}}}]{
                \texttt{\tiny{[P1$\langle$GM]}}
                \texttt{Mesajın diğer oyuncuya iletildi.} \\
\\ 
\\ 
\texttt{Şu anda en fazla 28 turdan 23. turu oynuyorsunuz.} \\
\texttt{Toplamda 16 cezanın 37 kadarını biriktirdiniz.Diğer oyuncu şu mesajı gönderdi:} \\
\\ 
\\ 
\texttt{"4. satırda uygun boş hücre olup olmadığını kontrol edelim. Hangi X koordinatlarının içeriği boş değilse, hangisinin boş olup olmadığını kontrol edelim. Örn. (4, 1) T, (4, 2) S, (4, 3) M, (4, 4) ┌. Yani 4. satırda hücrelerin içeriği: S, M, ┌ dışındaki hücreler nerede? Kontrol edelim."} \\
\\ 
\\ 
\texttt{Bir sonraki komutun nedir?} \\
            }
        }
    }
    & \\ \\

    \theutterance \stepcounter{utterance}  
    & \multicolumn{4}{p{0.6\linewidth}}{
        \cellcolor[rgb]{0.9,0.9,0.9}{
            \makecell[{{p{\linewidth}}}]{
                \texttt{\tiny{[P1$\rangle$GM]}}
                \texttt{SÖYLE: 4. satırda (4, 2) S, (4, 3) M, (4, 4) ┌, içerdiğini kontrol edelim.} \\
            }
        }
    }
    & & \\ \\

    \theutterance \stepcounter{utterance}  
    & & & \multicolumn{2}{p{0.3\linewidth}}{
        \cellcolor[rgb]{0.9,0.9,0.9}{
            \makecell[{{p{\linewidth}}}]{
                \texttt{\tiny{[GM$|$GM]}}
                \texttt{SÖYLE: 4. satırda (4, 2) S, (4, 3) M, (4, 4) ┌, içerdiğini kontrol edelim.} \\
            }
        }
    }
    & & \\ \\

    \theutterance \stepcounter{utterance}  
    & & \multicolumn{4}{p{0.6\linewidth}}{
        \cellcolor[rgb]{0.9,0.9,0.9}{
            \makecell[{{p{\linewidth}}}]{
                \texttt{\tiny{[GM$\rangle$P2]}}
                \texttt{Mesajın diğer oyuncuya iletildi.} \\
\\ 
\\ 
\texttt{Şu anda en fazla 28 turdan 24. turu oynuyorsunuz.} \\
\texttt{Toplamda 16 cezanın 37 kadarını biriktirdiniz.Diğer oyuncu şu mesajı gönderdi:} \\
\\ 
\\ 
\texttt{"4. satırda (4, 2) S, (4, 3) M, (4, 4) ┌, içerdiğini kontrol edelim."} \\
\\ 
\\ 
\texttt{Bir sonraki komutun nedir?} \\
            }
        }
    }
    & \\ \\

    \theutterance \stepcounter{utterance}  
    & & & \multicolumn{4}{p{0.6\linewidth}}{
        \cellcolor[rgb]{0.9,0.9,0.9}{
            \makecell[{{p{\linewidth}}}]{
                \texttt{\tiny{[GM$\langle$P2]}}
                \texttt{HAMLE: A, (4, 5)} \\
            }
        }
    }
     \\ \\

    \theutterance \stepcounter{utterance}  
    & & & \multicolumn{2}{p{0.3\linewidth}}{
        \cellcolor[rgb]{0.9,0.9,0.9}{
            \makecell[{{p{\linewidth}}}]{
                \texttt{\tiny{[GM$|$GM]}}
                \texttt{HAMLE: A, (4, 5)} \\
            }
        }
    }
    & & \\ \\

    \theutterance \stepcounter{utterance}  
    & & & \multicolumn{2}{p{0.3\linewidth}}{
        \cellcolor[rgb]{0.9,0.9,0.9}{
            \makecell[{{p{\linewidth}}}]{
                \texttt{\tiny{[GM$|$GM]}}
                \texttt{Ceza: (4,5) boş değil; "│" içeriyor.} \\
\\ 
\texttt{Toplamda 16 cezanın 38 kadarını biriktirdiniz.Lütfen tekrar deneyin!} \\
            }
        }
    }
    & & \\ \\

    \theutterance \stepcounter{utterance}  
    & & \multicolumn{4}{p{0.6\linewidth}}{
        \cellcolor[rgb]{0.9,0.9,0.9}{
            \makecell[{{p{\linewidth}}}]{
                \texttt{\tiny{[GM$\rangle$P2]}}
                \texttt{Ceza: (4,5) boş değil; "│" içeriyor.} \\
\\ 
\texttt{Toplamda 16 cezanın 38 kadarını biriktirdiniz.Lütfen tekrar deneyin!} \\
            }
        }
    }
    & \\ \\

    \theutterance \stepcounter{utterance}  
    & & & \multicolumn{4}{p{0.6\linewidth}}{
        \cellcolor[rgb]{0.9,0.9,0.9}{
            \makecell[{{p{\linewidth}}}]{
                \texttt{\tiny{[GM$\langle$P2]}}
                \texttt{SÖYLE: 4. satırda uygun boş hücre olup olmadığını kontrol edelim. 4. satırda (4, 1) T, (4, 2) S, (4, 3) M, (4, 4) ┌, (4, 5) │, (4, 6) ┼, (4, 7) │. Yani hücreler dolu! Şimdi ne yapılabilir? Belki (1,4), (2,4) gibi başka satırlara taşınması gerekebilir? Ama hedef 4. satır! Stratejimizi gözden geçirelim.} \\
            }
        }
    }
     \\ \\

    \theutterance \stepcounter{utterance}  
    & & & \multicolumn{2}{p{0.3\linewidth}}{
        \cellcolor[rgb]{0.9,0.9,0.9}{
            \makecell[{{p{\linewidth}}}]{
                \texttt{\tiny{[GM$|$GM]}}
                \texttt{SÖYLE: 4. satırda uygun boş hücre olup olmadığını kontrol edelim. 4. satırda (4, 1) T, (4, 2) S, (4, 3) M, (4, 4) ┌, (4, 5) │, (4, 6) ┼, (4, 7) │. Yani hücreler dolu! Şimdi ne yapılabilir? Belki (1,4), (2,4) gibi başka satırlara taşınması gerekebilir? Ama hedef 4. satır! Stratejimizi gözden geçirelim.} \\
            }
        }
    }
    & & \\ \\

    \theutterance \stepcounter{utterance}  
    & & \multicolumn{4}{p{0.6\linewidth}}{
        \cellcolor[rgb]{0.9,0.9,0.9}{
            \makecell[{{p{\linewidth}}}]{
                \texttt{\tiny{[P1$\langle$GM]}}
                \texttt{Mesajın diğer oyuncuya iletildi.} \\
\\ 
\\ 
\texttt{Şu anda en fazla 28 turdan 24. turu oynuyorsunuz.} \\
\texttt{Toplamda 16 cezanın 38 kadarını biriktirdiniz.Diğer oyuncu şu mesajı gönderdi:} \\
\\ 
\\ 
\texttt{"4. satırda uygun boş hücre olup olmadığını kontrol edelim. 4. satırda (4, 1) T, (4, 2) S, (4, 3) M, (4, 4) ┌, (4, 5) │, (4, 6) ┼, (4, 7) │. Yani hücreler dolu! Şimdi ne yapılabilir? Belki (1,4), (2,4) gibi başka satırlara taşınması gerekebilir? Ama hedef 4. satır! Stratejimizi gözden geçirelim."} \\
\\ 
\\ 
\texttt{Bir sonraki komutun nedir?} \\
            }
        }
    }
    & \\ \\

    \theutterance \stepcounter{utterance}  
    & \multicolumn{4}{p{0.6\linewidth}}{
        \cellcolor[rgb]{0.9,0.9,0.9}{
            \makecell[{{p{\linewidth}}}]{
                \texttt{\tiny{[P1$\rangle$GM]}}
                \texttt{SÖYLE: Hedef 4. satır, ama tüm hücreler dolu. Şimdi veya başka bir hedef düşünelim (örn. 3. satır)?} \\
            }
        }
    }
    & & \\ \\

    \theutterance \stepcounter{utterance}  
    & & & \multicolumn{2}{p{0.3\linewidth}}{
        \cellcolor[rgb]{0.9,0.9,0.9}{
            \makecell[{{p{\linewidth}}}]{
                \texttt{\tiny{[GM$|$GM]}}
                \texttt{SÖYLE: Hedef 4. satır, ama tüm hücreler dolu. Şimdi veya başka bir hedef düşünelim (örn. 3. satır)?} \\
            }
        }
    }
    & & \\ \\

    \theutterance \stepcounter{utterance}  
    & & \multicolumn{4}{p{0.6\linewidth}}{
        \cellcolor[rgb]{0.9,0.9,0.9}{
            \makecell[{{p{\linewidth}}}]{
                \texttt{\tiny{[GM$\rangle$P2]}}
                \texttt{Mesajın diğer oyuncuya iletildi.} \\
\\ 
\\ 
\texttt{Şu anda en fazla 28 turdan 25. turu oynuyorsunuz.} \\
\texttt{Toplamda 16 cezanın 38 kadarını biriktirdiniz.Diğer oyuncu şu mesajı gönderdi:} \\
\\ 
\\ 
\texttt{"Hedef 4. satır, ama tüm hücreler dolu. Şimdi veya başka bir hedef düşünelim (örn. 3. satır)?"} \\
\\ 
\\ 
\texttt{Bir sonraki komutun nedir?} \\
            }
        }
    }
    & \\ \\

    \theutterance \stepcounter{utterance}  
    & & & \multicolumn{4}{p{0.6\linewidth}}{
        \cellcolor[rgb]{0.9,0.9,0.9}{
            \makecell[{{p{\linewidth}}}]{
                \texttt{\tiny{[GM$\langle$P2]}}
                \texttt{SÖYLE: Stratejimizi gözden geçirelim. Alternatif hedef olarak 3. satırda her türden nesneyi toplama fikrini düşünelim. 3. satırda (3,1), (3,2), (3,3), (3,4), (3,5) gibi hücreler hangisi boş? D, T, M, A gibi.} \\
            }
        }
    }
     \\ \\

    \theutterance \stepcounter{utterance}  
    & & & \multicolumn{2}{p{0.3\linewidth}}{
        \cellcolor[rgb]{0.9,0.9,0.9}{
            \makecell[{{p{\linewidth}}}]{
                \texttt{\tiny{[GM$|$GM]}}
                \texttt{SÖYLE: Stratejimizi gözden geçirelim. Alternatif hedef olarak 3. satırda her türden nesneyi toplama fikrini düşünelim. 3. satırda (3,1), (3,2), (3,3), (3,4), (3,5) gibi hücreler hangisi boş? D, T, M, A gibi.} \\
            }
        }
    }
    & & \\ \\

    \theutterance \stepcounter{utterance}  
    & & \multicolumn{4}{p{0.6\linewidth}}{
        \cellcolor[rgb]{0.9,0.9,0.9}{
            \makecell[{{p{\linewidth}}}]{
                \texttt{\tiny{[P1$\langle$GM]}}
                \texttt{Mesajın diğer oyuncuya iletildi.} \\
\\ 
\\ 
\texttt{Şu anda en fazla 28 turdan 25. turu oynuyorsunuz.} \\
\texttt{Toplamda 16 cezanın 38 kadarını biriktirdiniz.Diğer oyuncu şu mesajı gönderdi:} \\
\\ 
\\ 
\texttt{"Stratejimizi gözden geçirelim. Alternatif hedef olarak 3. satırda her türden nesneyi toplama fikrini düşünelim. 3. satırda (3,1), (3,2), (3,3), (3,4), (3,5) gibi hücreler hangisi boş? D, T, M, A gibi."} \\
\\ 
\\ 
\texttt{Bir sonraki komutun nedir?} \\
            }
        }
    }
    & \\ \\

    \theutterance \stepcounter{utterance}  
    & \multicolumn{4}{p{0.6\linewidth}}{
        \cellcolor[rgb]{0.9,0.9,0.9}{
            \makecell[{{p{\linewidth}}}]{
                \texttt{\tiny{[P1$\rangle$GM]}}
                \texttt{SÖYLE: 3. satırda (3,1) D, (3,2) simidi görülemiyorum, (3,3) A, (3,4) değerleri hangisi?} \\
            }
        }
    }
    & & \\ \\

    \theutterance \stepcounter{utterance}  
    & & & \multicolumn{2}{p{0.3\linewidth}}{
        \cellcolor[rgb]{0.9,0.9,0.9}{
            \makecell[{{p{\linewidth}}}]{
                \texttt{\tiny{[GM$|$GM]}}
                \texttt{SÖYLE: 3. satırda (3,1) D, (3,2) simidi görülemiyorum, (3,3) A, (3,4) değerleri hangisi?} \\
            }
        }
    }
    & & \\ \\

    \theutterance \stepcounter{utterance}  
    & & \multicolumn{4}{p{0.6\linewidth}}{
        \cellcolor[rgb]{0.9,0.9,0.9}{
            \makecell[{{p{\linewidth}}}]{
                \texttt{\tiny{[GM$\rangle$P2]}}
                \texttt{Mesajın diğer oyuncuya iletildi.} \\
\\ 
\\ 
\texttt{Şu anda en fazla 28 turdan 26. turu oynuyorsunuz.} \\
\texttt{Toplamda 16 cezanın 38 kadarını biriktirdiniz.Diğer oyuncu şu mesajı gönderdi:} \\
\\ 
\\ 
\texttt{"3. satırda (3,1) D, (3,2) simidi görülemiyorum, (3,3) A, (3,4) değerleri hangisi?"} \\
\\ 
\\ 
\texttt{Bir sonraki komutun nedir?} \\
            }
        }
    }
    & \\ \\

    \theutterance \stepcounter{utterance}  
    & & & \multicolumn{4}{p{0.6\linewidth}}{
        \cellcolor[rgb]{0.9,0.9,0.9}{
            \makecell[{{p{\linewidth}}}]{
                \texttt{\tiny{[GM$\langle$P2]}}
                \texttt{SÖYLE: 3. satırda hücreleri kontrol edelim. (3,1) D, (3,2) {-} ✅ boş (simidi yok), (3,3) A, (3,4) values hangi? Örn. (3,4) {-} hangi sembol ya da harf? Bize anlat.} \\
            }
        }
    }
     \\ \\

    \theutterance \stepcounter{utterance}  
    & & & \multicolumn{2}{p{0.3\linewidth}}{
        \cellcolor[rgb]{0.9,0.9,0.9}{
            \makecell[{{p{\linewidth}}}]{
                \texttt{\tiny{[GM$|$GM]}}
                \texttt{SÖYLE: 3. satırda hücreleri kontrol edelim. (3,1) D, (3,2) {-} ✅ boş (simidi yok), (3,3) A, (3,4) values hangi? Örn. (3,4) {-} hangi sembol ya da harf? Bize anlat.} \\
            }
        }
    }
    & & \\ \\

    \theutterance \stepcounter{utterance}  
    & & \multicolumn{4}{p{0.6\linewidth}}{
        \cellcolor[rgb]{0.9,0.9,0.9}{
            \makecell[{{p{\linewidth}}}]{
                \texttt{\tiny{[P1$\langle$GM]}}
                \texttt{Mesajın diğer oyuncuya iletildi.} \\
\\ 
\\ 
\texttt{Şu anda en fazla 28 turdan 26. turu oynuyorsunuz.} \\
\texttt{Toplamda 16 cezanın 38 kadarını biriktirdiniz.Diğer oyuncu şu mesajı gönderdi:} \\
\\ 
\\ 
\texttt{"3. satırda hücreleri kontrol edelim. (3,1) D, (3,2) {-} ✅ boş (simidi yok), (3,3) A, (3,4) values hangi? Örn. (3,4) {-} hangi sembol ya da harf? Bize anlat."} \\
\\ 
\\ 
\texttt{Bir sonraki komutun nedir?} \\
            }
        }
    }
    & \\ \\

    \theutterance \stepcounter{utterance}  
    & \multicolumn{4}{p{0.6\linewidth}}{
        \cellcolor[rgb]{0.9,0.9,0.9}{
            \makecell[{{p{\linewidth}}}]{
                \texttt{\tiny{[P1$\rangle$GM]}}
                \texttt{SÖYLE: 3. satırda (3,4) AD’den önce AD rüyasında değil, AD (A, D) hücresi, yani hücrede 'A' ve 'D' var, dolayısıyla hücre dolu.} \\
            }
        }
    }
    & & \\ \\

    \theutterance \stepcounter{utterance}  
    & & & \multicolumn{2}{p{0.3\linewidth}}{
        \cellcolor[rgb]{0.9,0.9,0.9}{
            \makecell[{{p{\linewidth}}}]{
                \texttt{\tiny{[GM$|$GM]}}
                \texttt{SÖYLE: 3. satırda (3,4) AD’den önce AD rüyasında değil, AD (A, D) hücresi, yani hücrede 'A' ve 'D' var, dolayısıyla hücre dolu.} \\
            }
        }
    }
    & & \\ \\

    \theutterance \stepcounter{utterance}  
    & & \multicolumn{4}{p{0.6\linewidth}}{
        \cellcolor[rgb]{0.9,0.9,0.9}{
            \makecell[{{p{\linewidth}}}]{
                \texttt{\tiny{[GM$\rangle$P2]}}
                \texttt{Mesajın diğer oyuncuya iletildi.} \\
\\ 
\\ 
\texttt{Şu anda en fazla 28 turdan 27. turu oynuyorsunuz.} \\
\texttt{Toplamda 16 cezanın 38 kadarını biriktirdiniz.Diğer oyuncu şu mesajı gönderdi:} \\
\\ 
\\ 
\texttt{"3. satırda (3,4) AD’den önce AD rüyasında değil, AD (A, D) hücresi, yani hücrede 'A' ve 'D' var, dolayısıyla hücre dolu."} \\
\\ 
\\ 
\texttt{Bir sonraki komutun nedir?} \\
            }
        }
    }
    & \\ \\

    \theutterance \stepcounter{utterance}  
    & & & \multicolumn{4}{p{0.6\linewidth}}{
        \cellcolor[rgb]{0.9,0.9,0.9}{
            \makecell[{{p{\linewidth}}}]{
                \texttt{\tiny{[GM$\langle$P2]}}
                \texttt{SÖYLE: 3. satırda hücreleri tekrar kontrol edelim. (3,1) D, (3,2) boş, (3,3) A, (3,4) AD (yani A ve D değil dolu). Yani 3. satırda boş hücre olup olmadığını kontrol edelim: (3,1) D, (3,2) {-} ✅ boş, (3,3) A, (3,4) AD. Yani 3. satırda boş hücre (3,2). Bu hücreye hangi nesneyi taşımayı düşünüyorsunuz?} \\
            }
        }
    }
     \\ \\

    \theutterance \stepcounter{utterance}  
    & & & \multicolumn{2}{p{0.3\linewidth}}{
        \cellcolor[rgb]{0.9,0.9,0.9}{
            \makecell[{{p{\linewidth}}}]{
                \texttt{\tiny{[GM$|$GM]}}
                \texttt{SÖYLE: 3. satırda hücreleri tekrar kontrol edelim. (3,1) D, (3,2) boş, (3,3) A, (3,4) AD (yani A ve D değil dolu). Yani 3. satırda boş hücre olup olmadığını kontrol edelim: (3,1) D, (3,2) {-} ✅ boş, (3,3) A, (3,4) AD. Yani 3. satırda boş hücre (3,2). Bu hücreye hangi nesneyi taşımayı düşünüyorsunuz?} \\
            }
        }
    }
    & & \\ \\

    \theutterance \stepcounter{utterance}  
    & & \multicolumn{4}{p{0.6\linewidth}}{
        \cellcolor[rgb]{0.9,0.9,0.9}{
            \makecell[{{p{\linewidth}}}]{
                \texttt{\tiny{[P1$\langle$GM]}}
                \texttt{Mesajın diğer oyuncuya iletildi.} \\
\\ 
\\ 
\texttt{Şu anda en fazla 28 turdan 27. turu oynuyorsunuz.} \\
\texttt{Toplamda 16 cezanın 38 kadarını biriktirdiniz.Diğer oyuncu şu mesajı gönderdi:} \\
\\ 
\\ 
\texttt{"3. satırda hücreleri tekrar kontrol edelim. (3,1) D, (3,2) boş, (3,3) A, (3,4) AD (yani A ve D değil dolu). Yani 3. satırda boş hücre olup olmadığını kontrol edelim: (3,1) D, (3,2) {-} ✅ boş, (3,3) A, (3,4) AD. Yani 3. satırda boş hücre (3,2). Bu hücreye hangi nesneyi taşımayı düşünüyorsunuz?"} \\
\\ 
\\ 
\texttt{Bir sonraki komutun nedir?} \\
            }
        }
    }
    & \\ \\

    \theutterance \stepcounter{utterance}  
    & \multicolumn{4}{p{0.6\linewidth}}{
        \cellcolor[rgb]{0.9,0.9,0.9}{
            \makecell[{{p{\linewidth}}}]{
                \texttt{\tiny{[P1$\rangle$GM]}}
                \texttt{HAMLE: A, (3, 2)} \\
            }
        }
    }
    & & \\ \\

    \theutterance \stepcounter{utterance}  
    & & & \multicolumn{2}{p{0.3\linewidth}}{
        \cellcolor[rgb]{0.9,0.9,0.9}{
            \makecell[{{p{\linewidth}}}]{
                \texttt{\tiny{[GM$|$GM]}}
                \texttt{HAMLE: A, (3, 2)} \\
            }
        }
    }
    & & \\ \\

    \theutterance \stepcounter{utterance}  
    & & & \multicolumn{2}{p{0.3\linewidth}}{
        \cellcolor[rgb]{0.9,0.9,0.9}{
            \makecell[{{p{\linewidth}}}]{
                \texttt{\tiny{[GM$|$GM]}}
                \texttt{A moved to (3, 2)} \\
            }
        }
    }
    & & \\ \\

    \theutterance \stepcounter{utterance}  
    & & \multicolumn{4}{p{0.6\linewidth}}{
        \cellcolor[rgb]{0.9,0.9,0.9}{
            \makecell[{{p{\linewidth}}}]{
                \texttt{\tiny{[GM$\rangle$P2]}}
                \texttt{Mesajın diğer oyuncuya iletildi.} \\
\\ 
\\ 
\texttt{Şu anda en fazla 28 turdan 28. turu oynuyorsunuz.} \\
\texttt{Toplamda 16 cezanın 38 kadarını biriktirdiniz.Diğer oyuncu kendi ızgarasında bir nesneyi hareket ettirdi.} \\
\\ 
\\ 
\texttt{Bir sonraki komutun nedir?} \\
            }
        }
    }
    & \\ \\

    \theutterance \stepcounter{utterance}  
    & & & \multicolumn{4}{p{0.6\linewidth}}{
        \cellcolor[rgb]{0.9,0.9,0.9}{
            \makecell[{{p{\linewidth}}}]{
                \texttt{\tiny{[GM$\langle$P2]}}
                \texttt{SÖYLE: Oyunun sona ermesini istiyorum. SÖYLE: bitti mi?} \\
            }
        }
    }
     \\ \\

    \theutterance \stepcounter{utterance}  
    & & & \multicolumn{2}{p{0.3\linewidth}}{
        \cellcolor[rgb]{0.9,0.9,0.9}{
            \makecell[{{p{\linewidth}}}]{
                \texttt{\tiny{[GM$|$GM]}}
                \texttt{SÖYLE: Oyunun sona ermesini istiyorum. SÖYLE: bitti mi?} \\
            }
        }
    }
    & & \\ \\

    \theutterance \stepcounter{utterance}  
    & & & \multicolumn{2}{p{0.3\linewidth}}{
        \cellcolor[rgb]{0.9,0.9,0.9}{
            \makecell[{{p{\linewidth}}}]{
                \texttt{\tiny{[GM$|$GM]}}
                \texttt{Invalid format: head is not empty: 'SÖYLE: Oyunun sona ermesini istiyorum. '} \\
            }
        }
    }
    & & \\ \\

    \theutterance \stepcounter{utterance}  
    & & & \multicolumn{2}{p{0.3\linewidth}}{
        \cellcolor[rgb]{0.9,0.9,0.9}{
            \makecell[{{p{\linewidth}}}]{
                \texttt{\tiny{[GM$|$GM]}}
                \texttt{Maximum number of penalties exceeded} \\
            }
        }
    }
    & & \\ \\

    \theutterance \stepcounter{utterance}  
    & & & \multicolumn{2}{p{0.3\linewidth}}{
        \cellcolor[rgb]{0.9,0.9,0.9}{
            \makecell[{{p{\linewidth}}}]{
                \texttt{\tiny{[GM$|$GM]}}
                \texttt{Initial states:} \\
\texttt{```} \\
\texttt{Player 1:} \\
\texttt{ 1234567} \\
\texttt{╔═══════╗} \\
\texttt{║◌◌◌◌◌T◌║ 1} \\
\texttt{║P◌◌◌◌◌◌║ 2} \\
\texttt{║◌◌AD◌┌─╢ 3} \\
\texttt{║◌◌◌┌─┤S║ 4} \\
\texttt{║O◌◌│◌├─╢ 5} \\
\texttt{╟───┼─┘◌║ 6} \\
\texttt{║◌M◌│◌◌◌║ 7} \\
\texttt{╚═══╧═══╝} \\
\\ 
\\ 
\texttt{Player 2:} \\
\texttt{ 1234567} \\
\texttt{╔═══════╗} \\
\texttt{║◌◌◌◌◌S◌║ 1} \\
\texttt{║◌◌◌◌O◌◌║ 2} \\
\texttt{║◌◌◌M◌┌─╢ 3} \\
\texttt{║◌◌A┌─┤◌║ 4} \\
\texttt{║◌◌◌│◌├─╢ 5} \\
\texttt{╟───┼─┘◌║ 6} \\
\texttt{║D◌T│P◌◌║ 7} \\
\texttt{╚═══╧═══╝} \\
\\ 
\texttt{```} \\
            }
        }
    }
    & & \\ \\

    \theutterance \stepcounter{utterance}  
    & & & \multicolumn{2}{p{0.3\linewidth}}{
        \cellcolor[rgb]{0.9,0.9,0.9}{
            \makecell[{{p{\linewidth}}}]{
                \texttt{\tiny{[GM$|$GM]}}
                \texttt{End states:} \\
\texttt{```} \\
\texttt{Player 1:} \\
\texttt{ 1234567} \\
\texttt{╔═══════╗} \\
\texttt{║◌◌◌D◌T◌║ 1} \\
\texttt{║P◌AS◌◌◌║ 2} \\
\texttt{║◌◌◌◌◌┌─╢ 3} \\
\texttt{║◌◌O┌─┤◌║ 4} \\
\texttt{║◌◌◌│◌├─╢ 5} \\
\texttt{╟───┼─┘◌║ 6} \\
\texttt{║◌M◌│◌◌◌║ 7} \\
\texttt{╚═══╧═══╝} \\
\\ 
\\ 
\texttt{Player 2:} \\
\texttt{ 1234567} \\
\texttt{╔═══════╗} \\
\texttt{║◌◌◌O◌◌◌║ 1} \\
\texttt{║◌◌◌S◌◌◌║ 2} \\
\texttt{║◌◌◌M◌┌─╢ 3} \\
\texttt{║◌◌A┌─┤◌║ 4} \\
\texttt{║◌◌◌│◌├─╢ 5} \\
\texttt{╟───┼─┘◌║ 6} \\
\texttt{║D◌T│P◌◌║ 7} \\
\texttt{╚═══╧═══╝} \\
\\ 
\texttt{```} \\
            }
        }
    }
    & & \\ \\

    \theutterance \stepcounter{utterance}  
    & & & \multicolumn{2}{p{0.3\linewidth}}{
        \cellcolor[rgb]{0.9,0.9,0.9}{
            \makecell[{{p{\linewidth}}}]{
                \texttt{\tiny{[GM$|$GM]}}
                \texttt{* success: False} \\
\texttt{* lose: False} \\
\texttt{* aborted: True} \\
\texttt{{-}{-}{-}{-}{-}{-}{-}} \\
\texttt{* Shifts: 3.00} \\
\texttt{* Max Shifts: 12.00} \\
\texttt{* Min Shifts: 6.00} \\
\texttt{* End Distance Sum: 29.45} \\
\texttt{* Init Distance Sum: 31.75} \\
\texttt{* Expected Distance Sum: 29.33} \\
\texttt{* Penalties: 39.00} \\
\texttt{* Max Penalties: 16.00} \\
\texttt{* Rounds: 27.00} \\
\texttt{* Max Rounds: 28.00} \\
\texttt{* Object Count: 7.00} \\
            }
        }
    }
    & & \\ \\

\end{supertabular}
}

\end{document}
